\subsection{The split local estimate at the conductor primes and at three}
\label{subsec:split-local-length}
\hypertarget{source-small-prime-lengths}{}

At the small split primes, the depth of the cyclotomic level enters the colength of the marked lattice. We keep this depth in the cohomology estimates before normalizing the length. Let \(\ell\) be
the coefficient prime at a retained localization and let
\(O\) be its DVR, with fraction field \(K_\ell\).
Write \(f,m,r,t,R\) for the host degree parameter, the number
of lifts, the rank of the lower block, the rank of the
character block, and the total rank. The parameter range is
\[
 1\le m\le f,\qquad 0\le r\le t,\qquad t,R\ge1.
\]

The CM field has degree \(2f\) over \(\mathbb Q\).
Use the level exponent \(s=\max(1,v_\ell(N))\) for the modulus defining the auxiliary field \(N=Q_Sqd\). At the small primes \(s\) is the
local level exponent in the row comparison. The same is true
at \(3\) when \(3\nmid Q_S\), with level exponent \(1\).

The primes in this comparison are
\[
 \mathcal S_{\rm sp}=\{\ell:\ell\mid Q_S\}\cup\{3\},\qquad
 Q_S=\operatorname{lcm}(2^{e+2},f_\psi).
\]
For \(\ell=2\), one has \(s\ge2\) and \(-d\equiv1\pmod8\).
An odd \(\ell\in\mathcal S_{\rm sp}\) is either 3 or a divisor of
the fixed conductor \(f_\psi\). Eventually every such prime is at
most \(w\) in Proposition~\ref{prop:quadratic-prime-pair}. Hence
\[
     -d\equiv1\pmod\ell\quad(\ell\text{ odd}),\qquad
     \mathbb Q(\sqrt{-d})\otimes\mathbb Q_\ell
          \simeq\mathbb Q_\ell\times\mathbb Q_\ell
                  \quad(\ell\in\mathcal S_{\rm sp}).
\]
Hensel's lemma gives the splitting at odd primes, and the dyadic
unit-square criterion gives it at two. Thus the row and column lie at conjugate places
even when \(K_\psi\) ramifies at \(\ell\). For each such completion
\(L\), put \(N'=N/\ell^{v_\ell(N)}\). Equation~\eqref{eq:host-character-compositum} gives
\[
    L\subseteq F_s=F_0(\mu_{\ell^s}),\qquad
    F_0=\mathbb Q_\ell(\mu_{N'}),\qquad
    F_0/\mathbb Q_\ell\text{ unramified},
\]
since \(\ell\nmid N'\) \cite[Corollary~10.16]{sutherland-local-fields}.
In particular \(s=1\) at three if \(3\nmid Q_S\).

Use the cohomology conditions on these completions. They include the finite-class lifting criterion, the bound on invariants over the cyclotomic subfield and the finite character criterion. The
split row and column constructions below invoke these
criteria on the restricted original lattices. In the
finite-class criterion the target class comes from
integral finite cohomology and the map of crystalline
lattices is surjective. In the character criterion of
Proposition~\ref{prop:cyclotomic-scalar-finite-condition},
the auxiliary extension is crystalline and has a
trivial rank-one quotient with a chosen integral lift
of its generator. \label{cor:small-prime-lengths-1}

\begin{lemma}[Invariants on the character string]\label{lem:small-prime-scalar-invariants}
At a place of \(E\) in the split case,
for \(1\le a<t\) write \(F\) for this completion. The normalized length satisfies
\[
 \operatorname{len}\bigl((K_\ell/O)(-a)^{G_F}\bigr)\le s+v_\ell(a)\le s+\lfloor\log_\ell t\rfloor ,
\]
uniformly in the twist and place.
\end{lemma}

\begin{proof}
For a place of \(E\) over \(\ell\), the prescribed cyclotomic
inertia image contains \(1+\ell^s\). On
\((K_\ell/O)(-a)\) with \(1\le a<t\), invariants lie in the
kernel of multiplication by
\((1+\ell^s)^a-1\). This kernel has normalized length
\[
     v_\ell((1+\ell^s)^a-1)=s+v_\ell(a).
\]
For the equality, the binomial formula gives
\(v_\ell((1+u)^b-1)=v_\ell(u)\) if \(\ell\nmid b\).

For a step of exponent \(\ell\), the first term \(\ell u\)
has valuation \(1+v_\ell(u)\). Each later term has greater
valuation when \(\ell\) is odd and \(v_\ell(u)\ge1\), or
when \(\ell=2\) and \(v_\ell(u)\ge2\).
Iterating through the \(\ell\)-part of \(a\) proves the
formula. The maximum invariant length over the indicated
twists and places is at most
\(s+\lfloor\log_\ell t\rfloor\). In particular the loss at two depends on its actual local depth. We retain that depth in passing from these invariant lengths to the split local conditions.
\end{proof}

\subsubsection*{The split row and column local conditions}

On the two sides of a conjugate split pair, the same good character supplies a column and a Tate-dual row. Their finite local conditions have different roles: the column provides classes to lift, and the row measures the colength of their images.

\begin{lemma}[Finite-class rank and colength on the small split row]\label{lem:small-split-finite-colength}
For every \(\ell\in\mathcal S_{\rm sp}\), with the split host
completions and cyclotomic depth \(s=\max(1,v_\ell(N))\)
of Section~\ref{subsec:split-local-length}, use the paired
places \(u,w\) of the same good character \(\pi\), the rank-\(r\) complement
\(W\) on the column and \(W^\vee(1)\) on the row.
The rational finite dimensions are at most
\([E_u:\mathbb Q_\ell]m\) on both sides. For the marked lattice
surjection, if the finite lift uses \(\ell^{e'}\) and the row
invariants are killed by \(\ell^{c_0(u)}\), the normalized row colength satisfies
\[
 \operatorname{colength}_{\rm row}\le [E_u:\mathbb Q_\ell]m e'+rc_0(u).
\]
\end{lemma}

\begin{proof}
\proofstep{Count the finite dimensions on the two signed jump lists.}
Use the same good Hecke character \(\pi\) on the retained
family. Write \(R_\pi^\circ\) for the full twisted
representation and \(W\) for its rank \(r\) constituent
after removing the direct sum of the \(t\) cyclotomic
characters. Write
\[
 T_2=W^\vee(1),
\]
using the convention
that the Tate twist by \(1\) has filtration jump \(-1\).
The local crystalline data give the following jumps,
with multiplicity one, for \(R_\pi^\circ\) at each
embedding of the completion into a splitting coefficient
field. The cyclotomic block has jumps \(0,\ldots,t-1\).

Of the other \(r\) jumps, \(k-1\) are in an interval of
strictly positive jumps at every embedding of a column
place and in an interval of strictly negative jumps at
every embedding of its conjugate row place. The remaining
\(m\) jumps can have either sign. These sets are disjoint,
and the two sets for \(W\) avoid both \(0\) and \(-1\).
Conjugation of the embedding type gives the switch of
sign on the row places. The local description of \(T_2\)
uses the following two exclusions at the same embedding:
\begin{equation}\label{eq:split-row-zero-jump}
 J(T_2)=\{-a-1:a\in J(W)\},\qquad
 0\notin J(W),\qquad -1\notin J(W)\ \Longrightarrow\ 0\notin J(T_2).
\end{equation}

At a column place \(u\), \(W\) is crystalline and
\(H^0(E_u,W)=0\), since a trivial subrepresentation
would contribute a zero jump. For a crystalline representation
\(V\) over a splitting coefficient field, the
Bloch--Kato local dimension formula is
\[
 \dim H^1_f(E_u,V)=
 \dim H^0(E_u,V)+
 \sum_{\tau:E_u\hookrightarrow K_\ell}
    \dim \bigl(D_{\mathrm{dR},\tau}(V)/
                       \operatorname{Fil}^0 D_{\mathrm{dR},\tau}(V)\bigr).
\]
Here and below one may enlarge the coefficient field to
contain all local embeddings. The filtration quotient
counts the negative jumps. Additivity of the crystalline
filtration under the split decomposition removes the
cyclotomic jumps from the full list. Thus
\(\dim H^1_f(E_u,W)\le [E_u:\mathbb Q_\ell]m\).

On the conjugate row place \(w\), \(T_2\) also has zero
rational invariants. Under the passage to the Tate dual the
\(k-1\) negative integer jumps of \(W\) become
nonnegative jumps of \(T_2\), by \(j\mapsto -j-1\).
Thus at most \(m\) jumps on each row embedding contribute:
\[
 \dim H^1_f(E_w,T_2)
   \le \sum_{\tau:E_w\hookrightarrow K_\ell} m
   =[E_w:\mathbb Q_\ell]m=[E_u:\mathbb Q_\ell]m .
\]
The comparison uses the Tate-dual row and its actual conjugate
column, even though their degrees agree.
\cite[Corollary 3.8.4]{bloch-kato}

\proofstep{Lift finite classes on the actual crystalline lattices and return by corestriction.}
For the integral condition, Proposition~\ref{prop:full-cyclotomic-finite-lift} gives an exponent \(e'\). Its input at a
column place is a surjection of stable lattices whose
rational representations are subquotients of finite
direct sums of the \(R_x^\circ\) in the family. At a row place the constituent family is instead obtained by the operation
\[
 R_x^\circ\longmapsto (R_x^\circ)^\vee(1).
\]

The family members are crystalline at these places,
and the absolute jumps of both families lie in the
bounded range used to choose \(e'\). Crystallinity and
the jump window pass to rational subquotients by
exactness and strictness of the crystalline filtration.
The local field is inside the prescribed cyclotomic
extension at level \(N\). The estimate says that the
cokernel of the map on integral \(H^1_f\) is killed
by \(\ell^{e'}\).

For this use of the Wach lattice estimate, the local
field \(L\) at a row or column is extended to
\[
 L'=\mathbb Q_\ell(\mu_{\operatorname{lcm}(\ell,N)}).
\]

The coefficients contain the corresponding roots of
unity. The extension \(L'/L\) is Galois and has
\(\ell\)-valuation of its degree at most
\(\lfloor\log_\ell N\rfloor+1\). Indeed
\[
 \ell^{v_\ell([L':L])}\le[L':L]\le[L':\mathbb Q_\ell]\le\ell N
 \quad\Longrightarrow\quad
 v_\ell([L':L])\le1+\lfloor\log_\ell N\rfloor .
\]
Restriction of a crystalline representation preserves
its jump window by base change at each embedding.
On \(L'\) the full cyclotomic level is \(s\).
The crystalline lattice comparison there gives the
lifting loss \(32(2hs+1)\) for jumps in \([-h,h]\).
The restriction--corestriction identity on the
finite local groups adds at most
\(\lfloor\log_\ell N\rfloor+1\) to the loss.

For \(c\in H^1_f(L,T)\) choose a source class \(d\)
over \(L'\) with \(f_*(d)=\ell^{32(2hs+1)}{\rm res}(c)\).
Naturality and restriction--corestriction give
\[
 f_*{\rm cor}(d)
   =\ell^{32(2hs+1)}{\rm cor}({\rm res}(c))
   =\ell^{32(2hs+1)}[L':L]c .
\] The prime-to-\(\ell\)
factor of the degree is a unit in the coefficients.
Thus the exponent \(e'\) in the conversion below
suffices. An invariant over \(L\) is already invariant
over \(L'\), so the invariant-killing argument uses
restriction alone.

\proofstep{Add the torsion length in the original coefficient normalization.}
The free rank is known; the integral row still contributes torsion.

Apply this at a row place to a lattice surjection
onto the target lattice \(T\) in \(T_2\) in its tied
matrix basis. The source sits rationally in a sum of the
dual-twisted family, and the target is therefore a
rational subquotient of that sum. The cokernel is a
quotient of \(H^1_f(E_w,T)/\ell^{e'}H^1_f(E_w,T)\).
This integral group is the preimage of the rational
finite part in \(H^1(E_w,T)\); by local cohomology
finiteness it is finitely generated. By the definition as an inverse image under rationalization,
\[
 \operatorname{rank}_O H^1_f(E_w,T)\le\dim_{K_\ell}H^1_f(E_w,T_2).
\]
The torsion lies in the image
of \(H^0(E_w,T_2/T)\) by the cohomology sequence.

If \(\ell^{c_0(u)}\) kills the torsion invariants, they embed into
\[
 \ell^{-c_0(u)}T/T\simeq T/\ell^{c_0(u)}T.
\]
The normalized length of the right side is \(r c_0(u)\).
Dividing lengths by the ramification degree of the
coefficient DVR gives the row colength bound
\[
       [E_u:\mathbb Q_\ell]m e' + r c_0(u).
\]
This computation gives the same bound on any of the
tied extensions of the coefficient field. Apply the degree-zero conclusion of
Lemma~\ref{lem:cyclotomic-invariant-loss} to \(W\) and \(T_2\).
The two exclusions in \eqref{eq:split-row-zero-jump} persist
under restriction to the full cyclotomic field, since a field
extension repeats each labelled list at the embeddings above it. The union of lists over
the embeddings of \(E_w\) has at most
\([E_w:\mathbb Q_\ell]r\) entries, so the bound uses the
degree of the base field and the rank, independently of
the degree of the splitting coefficient field.

The rank and colength comparison uses the crystalline lifts and invariant bounds for the specified local data.
The global lattice-index argument now uses these
local estimates on the same family.
\label{cor:small-prime-lengths-2}
\end{proof}

\paragraph{How the distinct jumps enter the local exponents.}
The two estimates involving the set of jumps use distinct
jump values. For a window \(U\subset [-h,h]\) with
\(|U|\le n\), put \(M=2s+n_0+1\), where \(n_0\) is least
with \(2h\le(\ell-1)\ell^{n_0}\), and put
\(\lambda=\lambda_N(h)=s+\lfloor\log_\ell(2h)\rfloor+
\lfloor\log_\ell N\rfloor+1\). The degree term is
\(D=v_\ell([F_M:F_s])=M-s\) on the full cyclotomic field.
The jump contribution to the local length is the positive integer
\begin{equation}\label{eq:local-jump-loss}
 C(k)=D+\sum_{i=0}^{k-1}\left(2+\max\left(1,
       \sum_{a\in U,\ a\ne i}(M+v_\ell(a-i))\right)\right),
 \qquad 1\le k\le h+1.
\end{equation}
Indeed \(n_0\le\lfloor\log_\ell(2h)\rfloor+1\), and
\(0<|a-i|\le2h\) in the inner sum. Hence
\(C(k)\le 2\lambda+3k(n+1)\lambda\).

For the invariant bound the fine loss is \(C(1)\),
so \(20(n+1)\lambda\) suffices. For a character of
jump \(e\), where \(0\le e<t\le h\), the full-field
specialization uses \(C(e+1)\) plus
\[
 (e+2)s+v_\ell(e)+\left\lfloor\frac{e}{\ell-1}\right\rfloor
       +1+s+v_\ell(e+1),
\]
with the zero term \(v_\ell(e)\) interpreted as \(0\) at
\(e=0\). The restriction--corestriction loss at the smaller
field is at most \(\lfloor\log_\ell N\rfloor+1\).
Together these are bounded by \(40(t+1)(n+1)\lambda\),
since \(e+1\le t\), \(s\le\lambda\) and each valuation
at a positive integer is at most that integer.
 It uses the
fine lifting and invariant losses of the local lattice
comparison. \label{cor:small-prime-lengths-3}

\subsubsection*{The integral Hodge--Tate estimate at full cyclotomic level}
\label{subsec:ramified-ht}
\hypertarget{source-ramified-cyclotomic}{}

The small-prime comparison needs an integral bound for invariants. We first pass to the field containing all roots of unity at the chosen cyclotomic level and keep track of the power of the coefficient prime needed to return to the original lattice.

Let \(F_0/\mathbb Q_\ell\) be finite unramified and
\(F_j=F_0(\mu_{\ell^j})\). The local lattices in this
argument are first restricted to \(F_s\), with \(s\ge1\)
and \(s\ge2\) at two. Include \(F_s\) in the finite
coefficient field. Normalize \(v(\ell)=1\) and let
\(L\subset V\) be the original stable lattice.
Write its jump interval as \([-a,b]\), and choose the depth by
\[
 H=a+b,\qquad n_0=\min\{n\ge0:H\le(\ell-1)\ell^n\},\qquad M=2s+n_0+1.
\]

\begin{proposition}[Integral cyclotomic comparison at the full ramified level]
\label{prop:ramified-integral-wach-comparison}\label{hyp:ramified-integral-wach-comparison}
Put \(S=\mathcal O_{F_0}[[\pi]]\), \(X=(1+\pi)^{\ell^{s-1}}-1\), \(Q=\varphi(X)/X\), and \(\Gamma_s=\operatorname{Gal}(F_\infty/F_s)\), where \(F_\infty=\bigcup_jF_j\). Let \(T\) be a crystalline lattice over \(F_s\) with jumps in \([-a,b]\). Its cyclotomic module \(N=N(T)\) is finite free over \(S\), is functorial in \(T\), and generates the associated integral étale module. For a generator \(\gamma\) of \(\Gamma_s\),
\[
 (\gamma-1)N\subset XN,\qquad
 Q^a\varphi^*N\subset N,\qquad Q^bN\subset\varphi^*N.
\]
Here \(\varphi^*N\) denotes the image of linearized Frobenius.

Inside the common period space the corresponding bounds are
\[
  X^b(A_{\inf}\otimes T)\subset A_{\inf}N
                     \subset X^{-a}(A_{\inf}\otimes T).
\]
At every root of \(Q\), write \(t=\log(1+\pi)\) for a local uniformizer. In a basis \((v_j)_j\) adapted to the filtration, indexed with multiplicity by its jumps, the rational completed lattice is \(\bigoplus_j F'[[t]]t^{-j}v_j\) over a splitting field \(F'\). The module of a kernel is the kernel on integral modules.
\end{proposition}

\begin{proof}
\proofstep{Specialize the Kisin--Ren classification.}
Use the multiplicative Lubin--Tate group over \(L=\mathbb Q_\ell\) with uniformizer \(\ell\). Its \(\ell^s\)-division points are \(\zeta-1\) for \(\zeta^{\ell^s}=1\), and \(K=F_s\) is contained in \(F_0L_\infty\). Thus one may take \(m=s\) in the construction of Kisin and Ren. Their three elements are
\[
 [\ell](\pi)=(1+\pi)^\ell-1,\quad
 P=[\ell^{s-1}](\pi)=X,\quad
 [\ell^s](\pi)/[\ell^{s-1}](\pi)=Q .
\]
Since \(K\otimes_{\mathbb Q_\ell}L=K\), their crystalline condition imposes no further filtration condition. The Lie algebra condition of \(L\)-analyticity is automatic for \(L=\mathbb Q_\ell\). Their integral equivalence therefore applies to the underlying \(\mathbb Z_\ell\)-lattices of all the representations here. We use the covariant module, equivalently the contravariant construction on the dual lattice. \cite[Sections~2.1, 2.4 and Corollary~3.3.8]{kisin-ren-lubin-tate}

Any extra coefficient action comes along by functoriality. For coefficient integers \(O\) containing \(\mathcal O_{F_0}\), the unramified embeddings give
\[
 O\otimes_{\mathbb Z_\ell} S\simeq
   \prod_{\tau:F_0\hookrightarrow\operatorname{Frac}O} O[[\pi]] .
\] On each summand, reduction of \(N\) modulo \(\pi\) is \(\ell\)-torsion-free and hence free over \(O\). Lift a basis. It generates by Nakayama's lemma and is independent by successive reduction modulo \(\pi\), since \(N\) has no \(\pi\)-torsion.

\proofstep{Read the filtration and the two integral bounds.}
The filtered-module construction makes modifications at the roots of \(Q,\varphi Q,\varphi^2 Q,\ldots\). At a root of \(Q\) the defining completed lattice is the degree-zero filtration on \(F'((t))\otimes D_{\rm dR}\); the preceding root reached by Frobenius is unmodified. An adapted basis therefore gives
\[
   N^{\wedge}=\bigoplus_j F'[[t]]t^{-j}v_j,\qquad
   (\varphi^*N)^{\wedge}=\bigoplus_j F'[[t]]v_j .
\]
Here \(Q\) has a simple zero, so these relative exponents give both height bounds after inverting \(\ell\). Integrally, the Frobenius matrix and its inverse already have entries in \(S[1/Q]\). The required products by powers of \(Q\) have no poles at any factor of \(Q\) by the calculation just made. Hence they lie in \(S\): division by the distinguished monic polynomial \(Q^k\) has an integral remainder, and a remainder vanishing after inverting \(\ell\) is zero. \cite[Section~2.2, Lemmas~2.2.1 and~2.2.5]{kisin-ren-lubin-tate}

The lower height bound gives
\[
      \varphi^*(X^a N)=X^aQ^a\varphi^*N\subseteq X^aN .
\] Lemma~3.2.1 of Kisin and Ren evaluates every integral functional of its étale lattice on this module in their integral period ring \(S^{\rm ur}\subset A_{\inf}\). Multiplication by \(X\) is invertible in the étale coefficient ring. Evaluating on a basis of \(T^\vee\) thus gives \(X^aN\subset A_{\inf}\otimes T\). The dual module has jumps in \([-b,a]\); the identical argument yields
\[
     X^bN^\vee\subset A_{\inf}\otimes T^\vee
        \quad\Longrightarrow\quad
     X^b(A_{\inf}\otimes T)\subset A_{\inf}N .
\]
The implication uses the perfect evaluation pairing for a finite free module. \cite[Lemma~3.2.1 and Corollary~3.3.8]{kisin-ren-lubin-tate}

\proofstep{Pass from the parameter \(\pi\) to the full divisor \(X\).}
The classification gives trivial action modulo \(\pi\). Consider any \(z=\zeta-1\) with \(\zeta^{\ell^{s-1}}=1\). The group \(\Gamma_s\) fixes \(z\), and \(Q(z)=\ell\ne0\). Over the fraction field, Frobenius identifies the fibre at \(z\) with the coefficient twist of the fibre at \(\zeta^\ell-1\). Frobenius commutes with \(\gamma\). Iteration from the fibre at zero shows that \(\gamma\) acts trivially on all these fibres. They are the distinct roots of \(X\), so monic division gives
\[
                         (\gamma-1)N\subset XN.
\]

For a surjection \(T_1\to T_2\) with kernel \(T'\), the rational modified lattices are exact. Away from their modifications this follows from exactness of the crystalline modules. At each modification the completed lattices are the filtered modules above, and the induced map is strictly surjective for the filtration. This also proves exactness on \(S[1/\ell]\): any nonunit irreducible there is represented by a distinguished polynomial and is tested in the corresponding local completion.

The integral module for \(T'\) is the intersection of its étale lattice with its rational module, as in the proof of Corollary~3.3.8 of Kisin and Ren. Combining rational and étale exactness yields
\[
                         N(T')=\ker(N(T_1)\longrightarrow N(T_2)).
\]
The integral cokernel \(C\) is killed by a power of \(\ell\) by rational surjectivity. After reduction modulo \(\ell\) it vanishes upon inverting \(\pi\), by étale surjectivity. Hence \(\pi^d C\subset\ell C\) for some \(d\); iteration shows that some power of \(\pi\) kills \(C\). Section~\ref{subsec:ramified-lift} removes the dependence on this preliminary power. \cite[Proposition~2.2.6 and Corollary~3.3.8]{kisin-ren-lubin-tate}
\end{proof}

\begin{lemma}[Finite classes and restricted Herr pairs]\label{lem:ramified-restricted-herr}
For \(T\) as in Proposition~\ref{prop:ramified-integral-wach-comparison} with jumps in \([-h,h]\), put \(N=N(T)\). Every cocycle in
\[
   Q^{-h}N\oplus XN,\qquad
   d(y,z)=(\gamma-1)y-(\varphi-1)z ,
\]
defines an integral class with crystalline rational extension. After inverting \(\ell\), these pairs represent all finite classes, and two such pairs have the same class precisely when their difference is \(((\varphi-1)w,(\gamma-1)w)\) with \(w\in N[1/\ell]\).
\end{lemma}

\begin{proof}
Given a cocycle \((y,z)\), adjoin one generator \(v\) to \(N\). Define its Frobenius and cyclotomic action by
\[
      \varphi(v)=v+y,\qquad\gamma(v)=v+z .
\]
The cocycle equation says exactly that the operators commute. Write \(\Phi_N\) for the Frobenius matrix on \(N\). After inverting \(Q\), the linearized Frobenius has inverse
\[
   \begin{pmatrix}\Phi_N & y\\ 0 & 1\end{pmatrix}^{-1}
     =\begin{pmatrix}\Phi_N^{-1}&-\Phi_N^{-1}y\\0&1\end{pmatrix}.
\] The cyclotomic generator acts trivially modulo \(\pi\) because \(z\in XN\). Its powers extend continuously to \(\Gamma_s\simeq\mathbb Z_\ell\). Indeed, \(\gamma\) acts trivially on \(I/I^2\) and on the module modulo \(I\), where \(I=(\ell,\pi)\). For \(M=N\oplus Sv\) this gives
\[
  (\gamma-1)I^kM\subset I^{k+1}M,\qquad
  \lim_{n\to\infty}\gamma^{\ell^n}=1\quad\text{on each }M/I^kM
\]
by the binomial expansion.

This is an object of the Kisin--Ren category for \(L=\mathbb Q_\ell\); its associated lattice is crystalline after tensoring with \(\mathbb Q_\ell\). In the étale module the two differences of \(v\) are its usual Herr representative. For the new lift \(v+w\) its pair is
\[
    (y,z)+\bigl((\varphi-1)w,(\gamma-1)w\bigr).
\] \cite[Theorem~1.6 and Corollary~3.3.8]{kisin-ren-lubin-tate}

Conversely, a rational finite class is a crystalline extension of the trivial representation by \(T[1/\ell]\). It admits a stable lattice extension after scaling the class: a continuous cocycle on the compact Galois group has bounded image. On the corresponding cyclotomic modules the surjection to the trivial module is rationally onto by Proposition~\ref{prop:ramified-integral-wach-comparison}. The kernel identity and the bounds on the extension module put the differences of any such lift \(v\) in
\[
  (\varphi-1)v\in Q^{-h}N[1/\ell],\qquad
  (\gamma-1)v\in XN[1/\ell].
\]

Indeed, if \(N'\) is the extension module, multiplication by \(Q^h\) reduces the first intersection to the kernel identity. For the second, \(Xn'\) in the kernel implies that the image of \(n'\) in the free trivial quotient is zero. Thus the intersections inside the rational étale module are
\[
 Q^{-h}N'[1/\ell]\cap D_{\rm et}(T[1/\ell])=Q^{-h}N[1/\ell],
 \qquad XN'[1/\ell]\cap D_{\rm et}(T[1/\ell])=XN[1/\ell].
\]
Here \(D_{\rm et}\) denotes the cyclotomic étale module.

An equality of rational extension classes gives an isomorphism of extensions inducing the identity on kernel and quotient. Full faithfulness of the cyclotomic equivalence carries it back to the modules. It changes the chosen lift by a single \(w\in N[1/\ell]\), giving the asserted description of coboundaries. \cite[Corollaries~2.4.4 and~3.3.8]{kisin-ren-lubin-tate}
\end{proof}

At the full level, Proposition~\ref{prop:ramified-integral-wach-comparison} gives the integral Hodge--Tate bounds below. Lemma~\ref{lem:ramified-restricted-herr} will enter when lifting finite classes across a surjection.\label{cor:ramified-cyclotomic-1}

\begin{proposition}[Integral Hodge--Tate coordinates at full cyclotomic level]\label{prop:full-level-ht-coordinates}
With the module of Proposition~\ref{prop:ramified-integral-wach-comparison}, over \(F_s=F_0(\mu_{\ell^s})\), with \(F_0/\mathbb Q_\ell\)
unramified and \(s\ge2\) at two, for each embedding \(\tau\)
there are coordinates with image lattice \(B'\) and semilinear generator matrix \(U\) at level \(M\) such that
\[
 B_r\subseteq B'\subseteq\ell^{-1}B_r,\qquad v((U-1)_{ij})>1/(\ell-1)\quad\text{for all }i,j.
\] A vector of the rational
space whose image belongs to \(B'\) lies in the original
stable lattice, and the graded de Rham comparisons identify
the twist coordinates with the Hodge--Tate jumps at \(\tau\).
\end{proposition}

\begin{proof}
\proofstep{Specialize the two height inclusions in the same Wach basis.}
In the Wach comparison the height sandwich is
\[
 X^b (A_{\inf}\otimes L)\ \subseteq\ A_{\inf}N
    \ \subseteq\ X^{-a}(A_{\inf}\otimes L),
 \qquad X=(1+\pi)^{\ell^{s-1}}-1.
\]
Here the coefficient tensor products and the images
in the common period module are understood.
Use the Wach basis after multiplying its columns
by \(X^a\), so they can be read in the \(A_{\inf}\)
lattice.

Choose compatible primitive roots with \(\zeta_{\ell^0}=1\)
and put \(\epsilon=(1,\zeta_\ell,\zeta_{\ell^2},\ldots)\) and
\(\pi=[\epsilon]-1\). For every \(k\ge0\), Fontaine's map
\(\theta\) and the Witt Frobenius give
\[
  (\theta\circ\varphi^{-k})([\epsilon])
       =(\epsilon^{1/\ell^k})^\sharp=\zeta_{\ell^k}.
\]
The sharp map takes the limit of the \(\ell^j\)-th powers of
the lifted \(j\)-th components. Here every such power equals
\(\zeta_{\ell^k}\). We specialize at the absolute depth
\(d_0=s+n_0\). Thus
\begin{equation}\label{eq:full-level-fontaine-map}
 \pi\longmapsto\zeta_{\ell^{d_0}}-1,\qquad
 X\longmapsto c=\zeta_{\ell^{n_0+1}}-1,\qquad
 v(c)=\frac{1}{(\ell-1)\ell^{n_0}} .
\end{equation}

On \(\mathcal O_{F_0}\) the map acts by inverse unramified
Frobenius to power \(d_0\). It permutes the unramified embeddings.
In an integral \(\mathbb Z_\ell\)-basis the action of the extra coefficient
integers is by endomorphisms of the underlying Wach module.
Their matrix entries lie in \(\mathcal O_{F_0}[[\pi]]\).
Evaluation therefore puts these entries in
\(\mathcal O_{F_{d_0}}\subset\mathcal O_{F_M}\).
The change of basis to the original period lattice itself can
have entries in \(\mathcal O_{\mathbb C_\ell}\).

Taking the Wach basis columns times a
\(\mathbb Z_\ell\)-basis of the coefficient integers gives
a square matrix \(D\) on the underlying \(\mathbb Q_\ell\)
coordinates. The specialized sandwich puts its ball
image between \(c^H B_d\) and \(B_d\), where
\(B_d=\mathcal O_{\mathbb C_\ell}^d\).
Since
\[
 v(\ell)-v(c^H)=1-\frac{H}{(\ell-1)\ell^{n_0}}\ge0,
\]
it follows that \(D\) is invertible and
\[
                       \ell B_d\subseteq DB_d\subseteq B_d.
\]
The action and the coefficient multiplication in this
basis have matrices over \(\mathcal O_{F_M}\).
Indeed their columns are the specialized
Wach-coordinate expressions in the same basis.
The multiplication by \(X^a\) cancels in the action
comparison because \(c\) is fixed at level \(M\).
On the cyclotomic kernel the matrix part of the
action is the identity, by the \'etale--Wach
comparison. \label{cor:ramified-cyclotomic-2}

\proofstep{Increase the depth until the generator is strictly small in the integral model.}
For the generator estimate, start with a
cyclotomic generator at depth \(s\), with action
\(\gamma_0\), and set \(I=(\ell,\pi)\) in the
Wach coordinate ring. On \(S/(\pi)=\mathcal O_{F_0}\) the cyclotomic
generator acts trivially. Its action on \(N/\pi N\) is trivial
by \((\gamma_0-1)N\subset XN\) in
Proposition~\ref{prop:ramified-integral-wach-comparison}, and \(\gamma_0(\pi)/\pi\equiv1\pmod I\).
The Leibniz rule for a semilinear difference gives
\[
   (\gamma_0-1)(I^j N)\subset I^{j+1}N
           \qquad(j\ge 0).
\]
Let \(k=s+n_0+1\). Expand the operator on \(N\):
\[
  \gamma_0^{\ell^k}-1
       =\sum_{j=1}^{\ell^k}\binom{\ell^k}{j}(\gamma_0-1)^j .
\]
For \(j<\ell^{k-1}\), \(v_\ell(j)\le k-2\), so the small-index binomial terms satisfy
\[
 v_\ell\binom{\ell^k}{j}
 =k-v_\ell(j)+v_\ell\binom{\ell^k-1}{j-1}\ge2
 \qquad(1\le j<\ell^{k-1}).
\]

For \(j\ge\ell^{k-1}\), the
operator \((\gamma_0-1)^j\) takes values in
\(I^{\ell^{k-1}}N\). By \eqref{eq:full-level-fontaine-map}, the least valuation of a
generator of the specialized ideal \(I\) is
\(1/((\ell-1)\ell^{s+n_0-1})\). Since \(k-1=s+n_0\),
the valuations for the two ranges are respectively at least
\[
     2,\qquad
     \frac{\ell^{k-1}}{(\ell-1)\ell^{s+n_0-1}}
                =\frac{\ell}{\ell-1}.
\]
In both ranges the strict logarithm threshold holds because
\[
 \min\{2,\ell/(\ell-1)\}>1/(\ell-1).
\]

Consequently \(\gamma=\gamma_0^{\ell^k}\) at level \(M\)
has action matrix \(U\) with this entrywise bound on \(U-1\).
For \(u=\chi(\gamma_0)\) and a nonzero integer \(j\), the generator depth is computed by
\[
 v(u^j-1)=v(u-1)+v_\ell(j).
\]
For \(j\) prime to \(\ell\)
this follows by expanding the binomial and for
multiplication of \(j\) by \(\ell\) the valuation
increases by one. The latter step requires
\(v(u-1)>1/(\ell-1)\), which holds for the admitted
depths. Negative \(j\) give the same valuation by
inversion.

\proofstep{Project by the separability idempotent and choose a maximal minor.}
Fix a field embedding \(\tau:F_s\to L_{\rm coeff}\).
The separability idempotent \(E_\tau\) of
\(F_s\otimes_{\mathbb Q_\ell}\mathbb C_\ell\), acting
through \(\tau\) on the coefficients and the standard
embedding on the scalar factor, selects the
\(\tau\)-component. The conjugated idempotent
\(D^{-1}E_\tau D\) has entries in \(F_M\).
Write
\[
 E'=D^{-1}E_\tau D.
\]
Select a nonzero
maximal-rank minor of largest absolute value.
Multiply the selected rows by the inverse of
that minor's matrix to obtain a row map \(R\).
Its entries are integral by the minor ratio rule.

For the selected column inclusion \(S_0\), we have
\(RS_0=1\), \(RE'=R\), and \(E'=E'S_0R\).
Thus \(S_0\) is an integral section of \(R\).
The induced action on \(P=RD^{-1}\) is computed by
\(R U_g S_0\). This is well defined since \(\ker R=\ker E'\)
is stable and \(S_0\) lifts every coordinate.
It gives integral matrices and preserves the strict
bound for \(U_\gamma-1\).
In these coordinates the image \(B'=P B_d\) satisfies
\[
                         B_r\subseteq B'\subseteq \ell^{-1}B_r.
\]
The first inclusion follows from \(D B_d\subseteq B_d\)
and the integral section \(S_0\), and the second follows
from \(\ell D^{-1} B_d\subseteq B_d\).

To detect integrality on \(V\), extend \(\tau^{-1}\)
to an embedding of the coefficient field in
\(\mathbb C_\ell\). Evaluation in this embedding sends
\(B_d\) into the unit ball of the coefficient
coordinates and factors through \(P\) since it
factors through \(E_\tau\).
On a rational vector these evaluations have the
valuations of its coordinates in an integral
coefficient basis of \(L\). For \(v=\sum_i v_i e_i\) in this basis, the integrality test reads
\[
 Pv\in B'\quad\Longrightarrow\quad v_\ell(v_i)\ge0\ (\text{all }i)\quad\Longrightarrow\quad v\in L.
\]
Finally the graded de Rham
comparison, using the same separability idempotent
on both factors, gives a basis of the coordinate
space of vectors fixed after twists by jumps at \(\tau\).
Indeed adapted filtered vectors give their symbols in
\(\operatorname{gr} B_{\rm dR}\), and the projection
selects the \(\tau\)-part of each graded piece.
These are the Hodge--Tate coordinates of the
specialized lattice. \label{cor:ramified-cyclotomic-3}
\end{proof}

\begin{lemma}[Cohomology and invariant loss at full cyclotomic level]\label{lem:cyclotomic-invariant-loss}
In the cyclotomic comparison of Proposition~\ref{prop:full-level-ht-coordinates}, put
\[
 N_i=\max\left(1,\sum_{a\in J\setminus\{i\}}
                    (M+v_\ell(a-i))\right)
\]
for the indicated distinct jumps \(J\) in the marked comparison.
Let \(B(i)\) be the integral coordinate ball with twisted action,
and keep the projected jump set \(J\) at the chosen embedding.
The degree-one bound holds for every \(i\). For the second bound
take an embedding giving the integrality test of
Proposition~\ref{prop:full-level-ht-coordinates} with \(0\notin J\).
Then, for the original lattice \(L\subset V\),
\[
 \ell^{1+N_i} H^1(G_{F_M},B(i))_{\rm tors}=0,\qquad
 0\notin J\ \Longrightarrow\ \ell^{2+N_0}(V/L)^{G_{F_M}}=0.
\]
These exponents are independent of the coefficient degree.
\end{lemma}

\begin{proof}
\proofstep{Use the integral trace on the cyclotomic kernel.}
Set \(H_\infty=G_{\bigcup_jF_j}\).
For a continuous cocycle of \(H_\infty\) in the scalar
ball \(\mathcal O_{\mathbb C_\ell}\), multiplication by
\(\ell\) gives an integral coboundary. In fact Tate's
trace theorem for a finite extension of the
cyclotomic tower puts every element of positive
valuation in the integral trace image. To approximate
the coboundary, choose an open normal subgroup on
which the cocycle \(f\) has norm at most \(\varepsilon\).
Take an integral \(b\) of trace \(\ell\) in the fixed
field of this subgroup. For a set \(S\) of coset
representatives put \(x=-\sum_{\sigma\in S}\sigma(b)f(\sigma)\).
In the difference \(hx-x\), replace each \(h\sigma\)
by its representative times an element of the subgroup.

If \(h\sigma=\tau n_\sigma\) with \(\tau\in S\), the
cocycle identity gives
\[
 h f(\sigma)=f(\tau)+ \tau f(n_\sigma)-f(h).
\]
Since \(n_\sigma(b)=b\) and the
\(\tau\)'s permute \(S\), we obtain
\[
       hx-x=\ell f(h)-\sum_{\sigma\in S}
                       \tau(b)\tau f(n_\sigma).
\]
The last sum has norm at most \(\varepsilon\) by the
ultrametric inequality and integrality of \(b\), uniformly
in \(h\).
Also \(\|x\|\le\sup_h\|f(h)\|\).
Apply this construction successively to the residual
cocycle divided by \(\ell\), with the error bounds
tending to zero geometrically. The successive corrections
then converge in the ball to the integral cobounding
vector. Apply the construction coordinatewise on
\(H_\infty\), where the matrix action and the
cyclotomic twist are trivial. Tate's fixed-field
comparison gives
\[
 \mathbb C_\ell^{H_\infty}=\widehat{F_\infty}.
\]

\cite[\S3, Propositions 9--10]{tate-pdivisible}

\proofstep{Split off the base by the compatible normalized trace.}
Let \(A_g=\chi(g)^iU_g g\) denote the semilinear
action in twist \(i\). On the completed integers of
\(F_\infty\), split off \(\mathcal O_{F_M}\) by the
normalized trace. At finite level \(j\ge M\) the
basis \(\zeta_{\ell^j}^{a}\) for
\(0\le a<\ell^{j-M}\) is integral over
\(\mathcal O_{F_M}\) by the Eisenstein uniformizer
description, with normalized trace the constant
coordinate. The integral projections are compatible
as \(j\) grows. Write \(Y\) for the kernel of the extended normalized trace. Compatibility gives the integral splitting
\[
 \widehat{\mathcal O}_{F_\infty}=\mathcal O_{F_M}\oplus Y.
\]
Both this module and the base are
stable under \(A_\gamma\).

On \(Y\), \(\gamma-1\) is invertible after inverting
\(\ell\), with inverse taking \((\zeta_\ell-1)Y\)
into \(Y\). To see the uniform bound, at each finite
level the nonzero basis exponents are permuted by
multiplication by \(u=\chi(\gamma)\) modulo
\(\ell^{j-M}\), with multipliers in the base roots of
unity. If \(v_\ell(a)=b_0\), the length of its orbit
is \(d_a=\ell^{\max(0,j-2M-b_0)}\). The return multiplier
is \(\zeta_{\ell^j}^{a(u^{d_a}-1)}\), which is nontrivial
since
\[
 v_\ell(a(u^{d_a}-1))=b_0+M+\max(0,j-2M-b_0)=\max(M+b_0,j-M)<j.
\]
Let \(c_a=\zeta_{\ell^j}^{a(u^{d_a}-1)}\) be this
return multiplier, whose difference from one has valuation
at most \(1/(\ell-1)\).

On the orbit module,
\[
   (\gamma-1)\sum_{i=0}^{d_a-1}\gamma^i
                          =\gamma^{d_a}-1=(c_a-1)\mathrm{id}.
\]
The sum of powers is integral, so the inverse obtained by
dividing by \(c_a-1\) sends \((\zeta_\ell-1)\) times
the orbit module into the orbit module. This proves the
integral bound and the inverses agree on the nested
finite trace-zero spaces. Since they commute with
the base uniformizer, they extend through completion
with the same integral bound. Now
\(u^iU_\gamma-1\) has strictly higher valuation than
\(\zeta_\ell-1\) entrywise. On this summand write \(D=\gamma-1\) and \(E=(A_\gamma-1)-D\).
The denominator of \(D^{-1}\) costs at most the valuation of
\(\zeta_\ell-1\), whereas \(E\) has strictly higher valuation. Hence
\(D^{-1}E\) raises the base-uniformizer filtration and
\[
 (A_\gamma-1)^{-1}=\sum_{n\ge0}(-D^{-1}E)^nD^{-1}
\]
converges in that filtration. In particular
\((A_\gamma-1)^{-1}\ell Y^r\subseteq Y^r\).

The Hodge--Tate basis vectors have coordinates fixed
by \(H_\infty\) and are fixed by \(A_\gamma\)
in their individual twists. By the trace-zero
invertibility just proved, their coordinates are
in the base field \(F_M\). Thus \(U_\gamma\) has
an annihilating polynomial with the distinct roots
\(u^{-a}\) for the set \(J\) of their jumps.
For twist \(i\), let
\[
 Q_i(Z)=\prod_{a\in J\setminus\{i\}}(Z-u^{-a}),\qquad
 N_i=\max\left(1,\sum_{a\in J\setminus\{i\}}
                   (M+v_\ell(a-i))\right).
\]
If \(w\) is integral on the base and
\(\ell^k w=(u^iU_\gamma-1)w'\) there, then
\(Q_i(U_\gamma)w=0\).
Integral polynomial division gives
\[
 Q_i(Z)-Q_i(u^{-i})=(Z-u^{-i})D_i(Z),\qquad D_i\in \mathcal O_{F_M}[Z].
\]
Substitute \(Z=U_\gamma\) and apply to \(w\). Since \(u^i\) is a unit, this gives
\[
 Q_i(u^{-i})w\in(A_\gamma-1)B(i),
\] where \(B(i)\) denotes the integral ball with the twisted action. The valuation of this scalar is the sum in \(N_i\). Together with the trace-zero
bound, this shows that \(\ell^{N_i}\) kills the
torsion in the integral generator coinvariants.

\proofstep{Return to integral cohomology, then kill invariants modulo the original lattice.}
The coinvariant bound is now used to clear denominators of invariant vectors.

For an integral cocycle \(f\) of \(G_{F_M}\) whose
class is \(\ell\)-power torsion, subtract the
integral coboundary from \(\ell f\) on \(H_\infty\).
The result descends to the procyclic quotient, with
values fixed by \(H_\infty\). Its torsion cobounding
vector is fixed there as well, since the cocycle
vanishes on the kernel. The generator calculation
makes the descended value at \(\gamma\) an integral
coboundary after multiplication by \(\ell^{N_i}\).
Equality at the generator extends by the cocycle
law and continuity. Restriction to the kernel and the generator calculation together give
\[
 \ell^{1+N_i}H^1(G_{F_M},B(i))_{\rm tors}=0.
\]

\label{cor:ramified-cyclotomic-4}

For an invariant class \(v\bmod L\), write \(y=Pv\).
The cocycle \(g\mapsto\ell(gy-y)\) in twist zero is
integral in the ball and a sufficiently large
\(\ell\)-power makes it an integral coboundary.
Applying the bound yields an integral vector whose
coboundary is that of \(\ell^{2+N_0}y\).
When the jumps avoid zero, the difference vanishes
by Tate's nonzero-twist invariant theorem over the
finite local field \(F_M\).
Thus
\[
 \ell^{2+N_0}y\in B_r\subseteq B'
\]
and \(\ell^{2+N_0}v\in L\).
This proves the invariant loss used in the row
estimate with no dependence on the coefficient degree.
\cite[Theorem 2]{tate-pdivisible}
\end{proof}


\subsubsection*{Finite-class lifting at full cyclotomic level}
\label{subsec:ramified-lift}

The preceding invariant bound treats a single lattice. Lifting finite cohomology classes across a surjection requires an additional estimate. We make the lift at the same full cyclotomic level before applying corestriction.

\begin{established}[Full cyclotomic field and crystalline height]\label{hyp:full-cyclotomic-height}
Let \(F_0/\mathbb Q_\ell\) be finite unramified, and set
\(F_s=F_0(\mu_{\ell^s})\) for \(s\ge1\), with \(s\ge2\)
at \(\ell=2\). The coefficient field contains \(F_s\),
as in Section~\ref{subsec:ramified-ht}. Let \(T_1\to T_2\)
be a surjection of integral crystalline lattices with
source jumps in \([-h,h]\) at every embedding. Put
\(H=2h\). For every embedding \(\tau\) and every generator \(\gamma\) of the full cyclotomic tower,
\[
 \operatorname{HT}_\tau(\ker f)\cup\operatorname{HT}_\tau(T_2)
 \subseteq[-h,h],\qquad v_\ell(\chi(\gamma)-1)=s,
 \qquad f:T_1\twoheadrightarrow T_2 .
\]
The integral comparison holds for these lattices.
\end{established}
\begin{proof}
The filtered de Rham functor is exact with strict maps on crystalline representations. Thus the graded pieces of the kernel and quotient vanish outside the source window. Since \(F_0\) is unramified and the cyclotomic extension of \(\mathbb Q_\ell\) is totally ramified, restriction identifies the group above \(F_s\) with \(1+\ell^s\mathbb Z_\ell\). In the indicated range, including \(s\ge2\) at two, the logarithm takes this group isomorphically onto \(\ell^s\mathbb Z_\ell\) and preserves valuations. Any generator therefore has precisely the asserted valuation. Proposition~\ref{prop:ramified-integral-wach-comparison} and Lemma~\ref{lem:ramified-restricted-herr} identify the fixed lattices over this full cyclotomic base. In their notation
\[
 X=(1+\pi)^{\ell^{s-1}}-1,\qquad Q=\varphi(X)/X .
\]
\label{cor:full-cyclotomic-lifting-1}
\end{proof}

\begin{proposition}[Finite lifting over the full cyclotomic field]\label{prop:full-cyclotomic-finite-lift}
For the crystalline surjection in Lemma~\ref{hyp:full-cyclotomic-height}, write \(f_*\) for its map on integral finite cohomology. Uniformly in the coefficient degree,
\[
 \ell^{32(2hs+1)}\operatorname{coker}(f_*)=0.
\]
The proof first lifts a representative after a Wach cokernel
loss, then corrects its actual Herr defect in the kernel.
At \(\ell=2\) the depth satisfies \(s\ge2\).
\end{proposition}

\begin{proof}
\proofstep{Control the cokernel of the integral Wach map.}
Write \(A\) for the integral Wach coefficient ring
and \(N_i\) for the Wach lattices. On each coefficient
letter, let \(\mathfrak a\) annihilate the quotient of
\(N_2\) by the image of \(N_1\).
The Wach length comparison gives
\(\pi^c\in\mathfrak a\) for some \(c\).
The height inclusions
\(Q^h\varphi^*N_i\subset N_i\) and
\(Q^hN_i\subset\varphi^*N_i\) give
\[
 Q^H\mathfrak a\subset A\varphi(\mathfrak a),
\]
with the transported ideal on the Frobenius letter.

In fact multiplying by \(Q^h\) on either side of
the original inclusion through the image passes from
the Frobenius target to the original target and
then back to the Frobenius source. The Frobenius
coordinate basis over \(\varphi(A)\) is
\(1,\epsilon,\ldots,\epsilon^{\ell-1}\), with
\(\epsilon=1+\pi\), and \(\varphi(\epsilon)=\epsilon^\ell\).
Thus the Frobenius extension of the annihilator is
the ideal with precisely the original ideal as all
its coefficient components. These statements use
the same Wach map in the étale and height comparisons.

The coefficient extraction can be written on the original and transported letters as
\[
 X^H\sum_j\epsilon^j\varphi(b_j)\in\mathfrak a\quad\Longrightarrow\quad X^H b_j\in\mathfrak a_{\rm prev}\quad(\text{each }j),
\]
where \(\mathfrak a_{\rm prev}\) is the corresponding annihilator on the preceding letter. Indeed multiply by \(Q^H\) before extracting the Frobenius components. Repeatedly perform
this operation on \(X^H(\epsilon-1)^c\).
Write \(f(\epsilon)=(\epsilon-1)^c\), and at each
extraction write
\[
 f(\epsilon)=\sum_{j=0}^{\ell-1}\epsilon^j f_j(\epsilon^\ell).
\]
The leading term lies in \(j=c\bmod\ell\). After
taking that component its coefficient is still one and its
degree becomes \(\lfloor c/\ell\rfloor\).

Iterate on that
component; each extraction preserves the property that
\(X^H\) times the current polynomial belongs to the
annihilator on the corresponding letter. If \(d_j\) denotes the successive degrees, the monic extraction gives
\[
 d_0=c,\qquad d_{j+1}=\lfloor d_j/\ell\rfloor.
\]
Thus the degree becomes zero after finitely many steps and the terminal monic constant is one.
Thus \(X^H\) lies in the annihilator. This step is uniform
even if \(c\) depends on the lattice.

Put \(u=\chi(\gamma)\).
The full-level action on \(N_i/XN_i\) is trivial and \(\gamma(X)/X\equiv u\pmod X\), so
\[
 \gamma|_{X^jN_i/X^{j+1}N_i}=u^j\operatorname{id}.
\]
Apply the polynomials with roots
\(u^0,\ldots,u^{H-1}\) and \(u^H,\ldots,u^{2H-1}\)
to \(X^H y\) for \(y\in N_2\).
The first result lies in \(X^H\operatorname{im}N_1\)
by first lifting \(X^Hy\), then using the first
polynomial on that lift. The other polynomial gives a vector in the deeper filtration step, and the common annihilator gives
\[
 X^{2H}N_2=X^H(X^HN_2)\subseteq X^H\operatorname{im}N_1.
\]

For \(H>0\), interpolate an inverse of the second
polynomial at the roots of the first.
Every denominator divides a product of the
differences between one root \(u^i\) and all the
other roots in \(0\le j<2H\).
Its valuation is at most
\[
 (2H-1)s+v_\ell(i!)+v_\ell((2H-1-i)!)\le
                 E_W:=2Hs+v_\ell((2H)!).
\]
Here the valuation of each nonzero difference is
\(s+v_\ell(i-j)\), by the admitted cyclotomic depths.
Write \(L_0(Z)=\prod_{i=0}^{H-1}(Z-u^i)\) and
\(L_1(Z)=\prod_{i=H}^{2H-1}(Z-u^i)\).

At each root of \(L_0\), Lagrange interpolation for
\(\ell^{E_W}/L_1(u^i)\) gives a polynomial \(B\) with
integral coefficients by the displayed denominator bound.
Then \(\ell^{E_W}-BL_1\) vanishes at all the roots of
\(L_0\), and monic division gives an integral \(A\) with
\[
    A L_0+B L_1=\ell^{E_W}.
\]
Applying this relation at \(\gamma\) to \(X^H N_2\) puts its
image inside \(X^H\operatorname{im}N_1\). Cancelling \(X^H\)
gives
\[
 \ell^{E_W}N_2\subseteq\operatorname{im}N_1.
\]

For \(H=0\), \(X^H=1\) already gives surjectivity.
\label{cor:full-cyclotomic-lifting-2}

\proofstep{Correct a defect in the kernel through the cyclotomic graded pieces.}

For the Wach lattice \(N\) of the kernel let
\(N_Q=Q^{-h}N\) in the étale module.
The restricted Herr pairs lie in \(N_Q\oplus XN\),
with differential
\[
 d(y,z)=(\gamma-1)y-(\varphi-1)z
\]
taking values in \(XN_Q\). This uses
\(\gamma(Q)/Q\equiv1\bmod X\) and the height
inclusion. The exactness comparison identifies
the kernel of the map on these coordinates by
intersection with the original Wach lattices.
The following additional inputs of that comparison
are used in the correction.

\begin{itemize}
\item \(N\) is \(\ell\)-saturated in \(N_Q\), and
  \(Xy\in N\) for \(y\in N_Q\) implies \(y\in N\).
  These follow by division for the monic cyclotomic
  factors of \(Q\) in the free Wach lattices.
\item \(1-\varphi\) is onto on \(X^{h+1}N\).
  Indeed
\[
 \varphi(X^{h+1}N)\subseteq QX^{h+1}N,
\]

  and \(\varphi(Q)\) and all its iterates lie in
  the maximal ideal. The geometric series converges
  in the complete Wach lattice.
\item With \(U\) the union of source jumps, put
  \(J'=\bigl((U+\{1,\ldots,h\})\setminus\{0\}\bigr)
  \cup\{1,\ldots,h\}\) and
\[
 P(Z)=\prod_{j\in J'}(Z-u^{-j}).
\]

  The polynomial \((Z-1)P(Z)\) annihilates
  \(N_Q/N\) under \(\gamma\).
  The following local calculation explains both the distinct factors and the integral conclusion.
\end{itemize}
To see why repeated jumps do not require repeated polynomial factors, invert \(\ell\) and separate the simple roots of \(Q\). Each is fixed by \(\gamma\), since its cyclotomic coordinate has order dividing \(\ell^s\). In the completed component take a filtration-adapted crystalline basis. In the logarithmic parameter \(t\) the generators are \(t^{-b}v_b\) for the jumps \(b\), with \(\gamma(t)=ut\). Modulo the original lattice, \(Q^{-h}N\) contributes the individual monomials \(t^{-b-j}v_b\), \(1\le j\le h\). In particular
\[
         \gamma(t^{-b-j}v_b)=u^{-(b+j)}t^{-b-j}v_b .
\]
The eigenvectors for equal exponents have the same eigenvalue. The single factor for that eigenvalue kills their whole span. Monic division and freeness make \(Q^{-h}N/N\) torsion-free over the coefficient DVR, so the same annihilating polynomial works integrally. Strictness places the kernel jumps among the source jumps. \label{cor:full-cyclotomic-lifting-3}

Let \(c_G=\prod_{j=1}^h(u^j-1)\).
Successively correcting powers \(X^j\) for
\(1\le j\le h\) shows that for \(z\in XN\) one has
\[
 c_Gz=(\gamma-1)x+t,\qquad x\in XN,\quad t\in X^{h+1}N.
\]
In each degree the correction replaces multiplication
by \(u^j-1\) by \(\gamma-1\) modulo the next degree.
Suppose \(\omega\in XN_Q\) and
\(\ell^k\omega=d(y,z)\) for a restricted pair.
Then
\[
 c_G\ell^k\omega
   =(\gamma-1)\bigl(c_Gy-(\varphi-1)x\bigr)
                      -(\varphi-1)t.
\]
Here the first argument is in \(N_Q\) and the last
term is in \(N\). Applying \(P(\gamma)\) puts the
right side in \(N\) by the quotient comparison.
Saturation gives \(P(\gamma)\omega\in N\), and
since \(XN_Q\) is stable this lies in \(XN\).

Apply the graded correction again to this vector.
The remaining term in \(X^{h+1}N\) is a
\((1-\varphi)\)-image there.
Finally \(P(\gamma)-P(1)\) is divisible by
\(\gamma-1\) with integral polynomial quotient.
It follows that \(c_G P(1)\omega\) is the
differential of a restricted pair.
The valuation loss is
\[
 E_U = \sum_{j\in J'}(s+v_\ell(j))+hs+v_\ell(h!).
\]

\proofstep{Lift a finite extension, then kill the defect of its Herr representative.}
Lifting an integral vector in a Wach module alone need not give the desired finite cohomology class. We apply the preceding correction to an extension representing that class and then pass to its Herr representative.


A finite target cocycle defines a crystalline
extension of the trivial lattice by the target
lattice. A lift of \(\ell^{E_W}\) in the Wach module
of the extension exists by the same cokernel
bound, since the extension adds only the zero jump.
Its Frobenius and cyclotomic differences give the
restricted Herr pair representing the scaled
class by the connecting comparison.
A restricted pair lifted termwise from
the target has a defect in \(XN_Q\) in the kernel,
provided its target pair is a cocycle.
For this defect \(z\), the rational comparison supplies some \(k\ge0\) and a restricted kernel pair \((a,b)\) with
\[
 \ell^k z=d(a,b).
\]

Indeed, for rational crystalline representations \(V_1\to V_2\) in the sequence the crystalline and de Rham functors are exact, the latter strictly filtered. In the Bloch--Kato presentation
\[
 H^1_f(F_s,V_i)=\operatorname{coker}\!\left[
 D_{\rm cris}(V_i)\xrightarrow{(1-\varphi,\,\mathrm{can})}
 D_{\rm cris}(V_i)\oplus D_{\rm dR}(V_i)/\mathrm{Fil}^0\right]
\]
the terms to the right are therefore surjective from \(i=1\) to \(i=2\). Lift the target class rationally. By Lemma~\ref{lem:ramified-restricted-herr}, adjust for the difference of its representative and the given target pair by a coboundary. Equivalently the differential of the termwise lift has zero class in the rational kernel complex. Clearing the finitely many coefficient denominators yields the displayed equality for some \(k\). \cite[Section~3]{bloch-kato-finite}
\label{cor:full-cyclotomic-lifting-4}

The cokernel bound lifts the two coordinates
of the target representative after another
\(\ell^{E_W}\). Subtract the kernel correction
from the scaled lifted pair. This is now a
restricted cocycle in the source and represents
the desired finite lift by naturality.
The lift of the extension, the lift of its coordinates and the correction of the defect require the respective multipliers
\[
 \ell^{E_W}\ell^{E_W}\ell^{E_U}=\ell^{2E_W+E_U}.
\]

Indeed \(J'\subseteq[-H,H]\setminus\{0\}\), so
\[
 \sum_{j\in J'}v_\ell(j)
 \le 2\sum_{j=1}^H v_\ell(j)
 =2\sum_{a\ge1}\lfloor H/\ell^a\rfloor\le2H.
\]
Also \(|J'|\le 2H\), \(h\le H\) and
\(v_\ell((2H)!)\le2H\).
These bounds give \(2E_W+E_U\le14Hs\) for \(s\ge1\)
and in particular the stated allowance \(32(2hs+1)\).
All bounds are independent of the coefficient
degree and of the Wach cokernel's initial
\(\pi\)-power. The restriction--corestriction
argument returns the finite lift to the
original local field and original lattices.
\end{proof}


\subsubsection*{The finite condition after projecting to a character}
\label{subsec:ramified-character}

A finite class lifted through the crystalline lattices must still satisfy the required condition after projection to a character. The extension class below isolates the difference between an integral lift of a generator and a Galois-invariant one.

Let \(F_s=F_0(\mu_{\ell^s})\) be the full cyclotomic field of
Section~\ref{subsec:ramified-ht}, and let \(O\) be the coefficient DVR.
Choose a crystalline lattice extension and a lift \(v\) of its quotient generator:
\[
 0\longrightarrow T\longrightarrow\widetilde T\xrightarrow{p}O\longrightarrow0,\qquad p(v)=1.
\]
Put \(f(g)=gv-v\in T\).
Suppose that \(\psi:T\to E/O(-e)\) is an equivariant
map, where \(E\) is the fraction field of the
coefficient integers \(O\), and \(0\le e<t\le h\).
The kernel jumps lie in a common set
\(U\subseteq[-h,h]\) by strictness of the crystalline
filtration. We bound the loss to make
\(\psi f\) a finite class. Use \(M\) from
Section~\ref{subsec:ramified-ht} with the interval
\([-h,h]\).

\paragraph{The truncated integral periods.}
Write \(A_{\inf}\) for the integral period ring,
\[
 \xi=([\epsilon]-1)/\varphi^{-1}([\epsilon]-1)
\]
for the cyclotomic generator of the kernel of
Fontaine's map, and \(t_{\rm dR}\) for the
cyclotomic period.
Modulo \(\xi^{e+1}\), the cocycle
\(f\otimes1\) is a rational coboundary.
In fact strictness gives a lift of one in filtration
zero of the de Rham invariants of \(\widetilde T\).
Subtract it from \(v\) in the de Rham tensor.
The result lies in filtration zero of the tensor
with \(T[1/\ell]\), and has coboundary \(f\otimes1\).
The quotient \(A_{\inf}/\xi^{e+1}\) after inverting
\(\ell\) is the corresponding de Rham jet, so some
\(\ell\)-power gives an integral cobounding vector.

At level \(M\), the degree \(i\) piece of the jet
is the integral scalar ball in twist \(i\).
For Fontaine's map \(\vartheta\), the ratio of the two generators satisfies
\[
 \vartheta(t_{\rm dR}/\xi)=\zeta_\ell-1.
\]
The difference between the two
generators in degree \(i\) rescales the ball by an
element fixed by \(G_{F_M}\). Use the unprojected
frame \(D\) obtained at absolute depth \(s+n_0\) in
\eqref{eq:full-level-fontaine-map}. The map \(\vartheta\) in the jet
formula is the original depth-zero quotient. Indeed, writing
\(\mu=[\epsilon]-1\), one has
\[
  t_{\rm dR}/\xi
   =\bigl(\log(1+\mu)/\mu\bigr)\varphi^{-1}(\mu)
   \ \xrightarrow{\ \vartheta\ }\ 1\cdot(\zeta_\ell-1).
\]

Both the scalar in this formula and the scalar \(c\) in
\eqref{eq:full-level-fontaine-map} are fixed by \(G_{F_M}\).
The transformed lattice lies between the ball and
its \(\ell^{-1}\) multiple. Multiply a torsion cocycle
once by \(\ell\) to enter the ball. A rational
cobounding vector can also be brought into the ball
by a sufficiently large \(\ell\)-power.
The degree-one assertion of Lemma~\ref{lem:cyclotomic-invariant-loss}
applies even when the auxiliary extension has a zero jump. It bounds
the additional loss by
\[
   1+\max\left(1,\sum_{a\in U,\ a\ne i}(M+v_\ell(a-i))\right).
\]
This gives the summand of \(C(e+1)\) in degree \(i\)
in the original integral tensor lattice.

These losses add through the truncation.
For successive rational jets \(J'\twoheadrightarrow J\), the necessary lifting statement is
\[
 H^0(G_{F_M},J')\longrightarrow H^0(G_{F_M},J)\longrightarrow0 .
\]
An adapted filtration basis in the de Rham comparison reduces surjectivity to scalar filtered periods. In a scalar interval the
invariants come from its degree-zero piece, if
present, by Tate's vanishing in nonzero degrees.
This piece lifts by constants in the field \(F_M\).
After clearing denominators this is an integral
lift up to a power of \(\ell\).

For two successive quotients, kill the projected torsion class and subtract its integral coboundary. The remaining graded cocycle is torsion: some power of \(\ell\) gives a cobounding vector in the whole jet. Its
projection is invariant; lift that invariant
after clearing denominators and subtract the lift.
The remaining cobounding vector belongs to the
graded piece. If the projected class is killed by \(\ell^u\) and the recovered graded torsion by \(\ell^v\), the original class \(c\) satisfies
\[
 \ell^v(\ell^u c)=0.
\]
This is the additivity of the successive exponent bounds. Restriction to \(F_M\) and
corestriction back add
\[
 M-s=v_\ell([F_M:F_s])
\]
once. Thus \(\ell^{C(e+1)}(f\otimes1)\) is a
coboundary over \(A_{\inf}/\xi^{e+1}\).
\label{cor:cyclotomic-character-lifting-1}

\begin{lemma}[Integral dual exponentials of positive twists at two]
\label{lem:dyadic-integral-period}\label{hyp:positive-character-integral-period}
Let \(F_0/\mathbb Q_2\) be finite unramified, \(F_s=F_0(\mu_{2^s})\), \(s\ge2\), and \(e\ge1\). Identify the de Rham line for \(\mathbb Q_2(-e)\) with \(F_s\) in the Tate basis \(t_{\rm dR}^e\). The integral cohomology classes have dual exponentials satisfying
\[
  2^{1-e-v_2((e-1)!)}\mathcal O_{F_s}
           \subseteq\exp^*\!\bigl(H^1(F_s,\mathbb Z_2(-e))\bigr).
\]
\end{lemma}

\begin{proof}
\proofstep{Construct integral classes before taking their moments.}
Choose compatible primitive roots \(\zeta_r\) of order \(2^r\). On \(\mathcal O_{F_0}[[X]]\) write \(Y=1+X\) and let \(\varphi\) act by arithmetic Frobenius \(\sigma\) on coefficients and by \(Y\mapsto Y^2\). Its left inverse \(\psi\) extracts the even powers in the decomposition \(f=\varphi(f_0)+Y\varphi(f_1)\), so \(\psi(f)=f_0\).

At level \(r\ge2\), put \(u=1+2^r\). For a positive odd integer \(a\) and \(b\in\mathcal O_{F_0}\) set
\[
    g=b(Y^a-Y^{au}),\qquad y=\sum_{n\ge0}\varphi^n(g).
\]
Since \(Y^{2^n}-1\in(2,X)^{n+1}\), the series converges in \(\mathcal O_{F_0}[[X]]\). Both exponents in \(g\) are odd, hence \(\psi(g)=0\) and \(\psi(y)=y\). The integral Iwasawa comparison for the trivial lattice attaches to \(y\) a system of integral cohomology classes. Its moment of order \(-e\) at level \(r\) belongs to \(H^1(F_r,\mathbb Z_2(-e))\). At two the torsion-free cyclotomic tower used for this comparison starts at \(r=2\). \cite[Proposition~II.1.2 and Theorem~II.1.3]{cherbonnier-colmez-iwasawa}

\proofstep{Evaluate a finite-level moment.}
Let \(y^{\sigma^{-r}}\) denote the series obtained by acting on its coefficients. In the chosen Tate basis, the dual exponential of the moment is the coefficient of \(t^e\) in
\[
   2^{-r}y^{\sigma^{-r}}(\zeta_r\exp(t/2^r)-1),\qquad t=t_{\rm dR}.
\]
The dual exponential formula of Cherbonnier and Colmez gives this series at all sufficiently high levels. To descend one step from \(j\) to \(j-1\ge2\), trace the series. The two conjugates of \(\zeta_j\) have opposite signs. Their average on \(y\) retains its even part \(\varphi(\psi y)=\varphi(y)\). Hence
\[
 \operatorname{Tr}_{F_j/F_{j-1}}\!
     \left(2^{-j}y^{\sigma^{-j}}(\zeta_j\exp(t/2^j)-1)\right)
 =2^{-(j-1)}y^{\sigma^{-(j-1)}}(\zeta_{j-1}\exp(t/2^{j-1})-1).
\]
Corestriction of the moment commutes with the dual exponential and this unnormalized trace. The coefficient formula therefore holds for each \(r\ge2\). \cite[Theorem~IV.2.1 and its proof]{cherbonnier-colmez-iwasawa}

After division by the odd unit \(a^e\), the computed value is \(C_r S_{r,a}(b)\), where
\[
 C_r=\frac{1-u^e}{2^{r(e+1)}e!},\qquad
 S_{r,a}(b)=\sum_{n\ge0}2^{ne}\sigma^{n-r}(b)\zeta_r^{a2^n}.
\]
Indeed the powers \(a2^n\) and \(au2^n\) give the same root of unity at level \(r\), while the coefficient of \(t^e\) in their exponentials introduces their \(e\)-th powers. The series for this coefficient converges since \(e>0\). For \(r\ge2\), the lifting-the-exponent formula yields
\[
    v_2(C_r)=r+v_2(e)-r(e+1)-v_2(e!)
            =-re-v_2((e-1)!).
\]

\proofstep{Recover each integral direction in the cyclotomic field.}
Replace \(a\) by \(a+2^{r-1}\). In the difference of the two normalized moments all the terms with \(n\ge1\) cancel. Since \(b\) is arbitrary,
\[
 C_r\bigl(S_{r,a}(b)-S_{r,a+2^{r-1}}(b)\bigr)
                   =2 C_r\sigma^{-r}(b)\zeta_r^a .
\]
Thus the image at level \(r\) contains \(2^{1-re-v_2((e-1)!)}\mathcal O_{F_0}\zeta_r^a\) for each odd \(a\).

The ground-field direction is also needed. At \(r=2\), add the normalized moments for \(a=1,3\). Their leading terms cancel, the roots in the next terms equal \(-1\), and all later roots equal one. Consequently
\[
 C_2\bigl(S_{2,1}(b)+S_{2,3}(b)\bigr)
 =2^{e+1}C_2\left(-\sigma^{-1}(b)+
             \sum_{n\ge2}2^{(n-1)e}\sigma^{n-2}(b)\right).
\]
The endomorphism in parentheses is congruent to \(-\sigma^{-1}\) modulo \(2\), hence is bijective on the finite free \(\mathbb Z_2\)-module \(\mathcal O_{F_0}\). The image contains \(2^{1-e-v_2((e-1)!)}\mathcal O_{F_0}\).

Restriction from \(F_r\) to \(F_s\) for \(2\le r\le s\) preserves integrality and includes the dual exponential value into \(F_s\). Each power of \(\zeta_s\) is \(\pm1\) or an odd power of a primitive \(\zeta_r\) with \(2\le r\le s\). These powers span \(\mathcal O_{F_0}[\zeta_s]=\mathcal O_{F_s}\) over \(\mathcal O_{F_0}\). The exponent \(1-re-v_2((e-1)!)\) is at most \(1-e-v_2((e-1)!)\), proving the assertion.
\end{proof}

\begin{proposition}[Returning scalar cocycles to the finite condition]\label{prop:cyclotomic-scalar-finite-condition}
Take the cyclotomic module of Proposition~\ref{prop:ramified-integral-wach-comparison} for the ramified lattice. Let \(f(g)=gv-v\) represent the crystalline lattice extension over \(F_s\),
and let \(\psi:T\to(E/O)(-e)\) be equivariant, with \(0\le e<t\le h\).
With the integer loss \(C(k)\) of \eqref{eq:local-jump-loss}, put
\[
 A_e=(e+2)s+v_\ell(e)+\lfloor e/(\ell-1)\rfloor+1,
 \qquad B_e=C(e+1)+s+v_\ell(e+1)+A_e,
\]
where the \(v_\ell(e)\) term in \(A_e\) is omitted when \(e=0\). Then the pushed-out class satisfies
\[
 \ell^{B_e}[\psi\circ f]\in H^1_f(F_s,(E/O)(-e)).
\]
Here the finite subgroup is the image of the rational finite condition.
For a class defined over the original local field at level \(N\),
corestriction adds at most \(1+\lfloor\log_\ell N\rfloor\) to the exponent.
\end{proposition}

\begin{proof}
\proofstep{Use the Bockstein allowance to lift the divisible cocycle to rational coefficients.}
Lift \(\psi\) on an integral
\(\mathbb Z_\ell\)-basis of \(T\) to the rational
target. The lift need only be equivariant modulo the
target lattice. Expand in an integral
\(\mathbb Z_\ell\)-basis of this target lattice.
Each coordinate of \(\psi f\) modulo \(\mathbb Z_\ell\)
is now a cocycle in
\(\mathbb Q_\ell/\mathbb Z_\ell(-e)\).

By local Tate duality,
\(H^2(F_s,\mathbb Z_\ell(-e))\) is dual to the
invariants in \(\mathbb Q_\ell/\mathbb Z_\ell(e+1)\).
The cyclotomic generator has \(v_\ell(\chi(\gamma)^{e+1}-1)=s+v_\ell(e+1)\). Duality therefore gives
\[
 \ell^{s+v_\ell(e+1)}H^2(F_s,\mathbb Z_\ell(-e))=0.
\]
The integral Bockstein sequence
then gives a rational cocycle representing the scaled
divisible cocycle in each coordinate. Adjust by a
rational coboundary lifting the discrepancy between
the divisible cocycles. The representatives then
agree pointwise modulo
\(\mathbb Z_\ell\).
The equivariance errors of the lifted coefficient
map are integral, so applying it to the jet
coboundary preserves the congruence modulo integral
periods. After scaling further by \(\ell^{C(e+1)}\),
each rational cocycle \(z\) so obtained satisfies
\[
 z(g)\equiv \chi(g)^{-e}g b-b
       \pmod {A_{\inf}+t_{\rm dR}^{e+1}B_{\rm dR}^+}
 \qquad (b\in B_{\rm dR}^+).
\]
The congruence means that the error is in the
image of \(A_{\inf}\) in the jet.

Use the dual exponential in the Tate basis \(t_{\rm dR}^e\). Its value \(\alpha\in F_s\) is
specified by the remaining term
\(\log\chi(g)\alpha t_{\rm dR}^e\) after a rational
twisted coboundary in that jet. To see this, filter the rational jet
\((B_{\rm dR}^+/t_{\rm dR}^{e+1})(-e)\) by powers of \(t_{\rm dR}\).
Its degree-\(i\) quotient is \(\mathbb C_\ell(i-e)\) for \(0\le i\le e\).
For \(i<e\), Tate's theorem gives
\(H^1(F_s,\mathbb C_\ell(i-e))=0\), so the cocycle in that quotient is
a coboundary. In degree \(e\) the twist is trivial, and Tate's calculation gives
\[
 H^1(F_s,\mathbb C_\ell)=F_s\,[\log\chi].
\]

When the dual exponential is zero, the Bloch--Kato
comparison gives a cobounding vector over
\(B_{\mathrm{cris}}\otimes_{\mathbb Q_\ell}K_\ell(-e)\),
where \(K_\ell\) is the coefficient field. This identifies
the corrected class with the crystalline finite condition.
\cite{tate-pdivisible,bloch-kato}\label{cor:cyclotomic-character-lifting-2}

\proofstep{Estimate the dual exponential in the truncated period coordinates.}
The Bockstein argument leaves a rational lift. Control of its projection to the finite subgroup now comes from the integral periods entering the dual exponential.

The displayed integral congruence bounds \(\alpha\).
At degree \(i<e\) use the integral multiplier
\[
 \ell^{1+s+v_\ell(i-e)}
\]
on the modified cobounding vector, then subtract an integral lift of its leading coefficient.
The degree \(i\) scalar belongs to the ball
modulo a twisted invariant, after this multiplication.
This is the rank-one ball argument at depth \(s\)
in twist \(i-e\), using the valuation of
\(\chi(\gamma)^{i-e}-1\).
In the \(t_{\rm dR}^i\) coordinates the integral
leading coefficient is multiplied by
\((\zeta_\ell-1)^i\), from \(t_{\rm dR}/\xi\).
Thus it can be lifted by Fontaine's surjection
and the power \(\xi^i\).
After eliminating all lower degrees the logarithm coefficient is
\[
 \ell^{e(s+1)+v_\ell(e!)}(\zeta_\ell-1)^e\alpha.
\]
Its two factors record the scalar elimination costs and the generator conversion, respectively.

If \(\alpha_1\log\chi(g)+g b_1-b_1\) is integral
and \(\alpha_1\in F_s\), apply the integral
coboundary estimate on the cyclotomic kernel
after multiplying by \(\ell\).
The remaining cobounding vector is fixed by that
kernel. Take the integral normalized trace at the
generator. It kills the difference of this
vector and its translate and gives
\(\ell\alpha_1\log\chi(\gamma)\) integral.
Since \(v_\ell(\log\chi(\gamma))=s\), one obtains
\[
 v_\ell(\alpha)\ge
 -e(s+1)-v_\ell(e!)-1-s-\frac{e}{\ell-1}.
\]

At weight zero the relevant integral image can be computed by local class field theory. Write \(E_0=(\ell-1)\ell^{s-1}\) for the ramification index of \(F_s/F_0\). The different is generated by \(\Phi'_{\ell^s}(\zeta_{\ell^s})\); its absolute valuation, with \(v_\ell(\ell)=1\), is \(s-1/(\ell-1)\). For \(v_\ell(z)\ge1/E_0\) the logarithmic series gives
\[
 v_\ell(\log(1+z))\ge \min_{n\ge1}\{n/E_0-v_\ell(n)\}
                  =1+1/(\ell-1)-s .
\]
For fixed \(v_\ell(n)=j\) the minimum occurs at \(n=\ell^j\); comparing consecutive values gives the equality. For the inverse different \(\mathfrak D^{-1}=\mathfrak D_{F_s/\mathbb Q_\ell}^{-1}\) this gives
\[
          \ell^{-1}\mathcal O_{F_s}\log(1+\mathfrak m_{F_s})
                              \subseteq\mathfrak D^{-1}.
\]

For \(a\in\ell^{-1}\mathcal O_{F_s}\) define on principal units
\[
       h_a(u)=\operatorname{Tr}_{F_s/\mathbb Q_\ell}(a\log u)\in\mathbb Z_\ell .
\]
Extend it by zero on prime-to-\(\ell\) roots of unity and a uniformizer. Local reciprocity gives an integral Galois character. The duality of the Bloch--Kato exponential with Kummer theory identifies its dual exponential with \(a\), up to the common reciprocity sign: pairing on principal units on both sides is the displayed trace. This proves the required inclusion for \(e=0\). \cite[Section~3]{bloch-kato-finite}

For odd \(p=\ell\) and \(e>0\), put \(m=e+1\) and \(L=F_s\). For \(m\ge2\) the rational Bloch--Kato exponential is an isomorphism. In its standard Tate basis put
\[
 \mathcal H_m=\exp_{L,m}^{-1}\!\left(\operatorname{im}\big[H^1(L,\mathbb Z_p(m))
                           \longrightarrow H^1(L,\mathbb Q_p(m))\big]\right)\subset L .
\] Let \(\sigma\) be Frobenius on \(F_0\), and let \(f\in\mathcal O_{F_0}[[X]]\) range over the power series satisfying \(\sum_{\zeta^p=1}f(\zeta(1+X)-1)=0\). Denote their module by \(R\).

The integral Coleman sequence and its explicit evaluation at level \(s\) give
\[
          p^{s+v_p(e)}\mathcal H_m\ \subseteq\ e!\,\Xi_{m,s}(R).
\]
Indeed the quotient of the full inverse-limit classes by the classes from principal units, after passing to the coinvariants for \(F_s\), is \(H^0(F_s,\mathbb Q_p/\mathbb Z_p(e))\). A generator with cyclotomic value \(1+p^s\) acts on that quotient through the relation
\[
       ((1+p^s)^e-1)x=0,\qquad v_p((1+p^s)^e-1)=s+v_p(e).
\] For the principal-unit classes the evaluation diagram sends each integral class to \((m-1)!\,\Xi_{m,s}(f)\) in exponential coordinates. Here \(L=F_s\) is already the entire cyclotomic level, so the trace from that level to \(L\) in the diagram is the identity. \cite[Sections~2.3.1--2.3.2, (4)--(5)]{benois-nguyen-tamagawa}

The explicit evaluation gives more than its determinant. In the formula for \(\Xi_{m,s}\), division by \(p^{m-s}\) makes the terms of index \(0\le k<s-1\) divisible by \(p^m\): their coefficients are \(p^{m(s-1-k)}\). The term for \(k=s-1\) and the constant correction give
\[
 p^{s-m}\Xi_{m,s}(f)\equiv
 f^{\sigma^{-1}}(\zeta_p-1)-(1-p^m\sigma^{-1})^{-1}f^{\sigma^{-1}}(0)
                      \pmod{p^m\mathcal O_L}.
\]
The inverse in the second term is the convergent geometric series and is the identity modulo \(p^m\). Evaluation of an integral series at \(\zeta_p-1\) agrees with its constant coefficient modulo \(\zeta_p-1\). The calculation also covers \(s=1\), when the earlier sum is empty. \cite[Section~2.2.1, definition of \(\Xi_{m,n}\)]{benois-nguyen-tamagawa}

In particular every vector in \(\mathcal H_m\) has valuation at least
\[
 A=v_p(e!)+m-2s-v_p(e)+1/(p-1).
\]
Under local Tate duality the image of \(H^1(L,\mathbb Z_p(1-m))\) in the dual de Rham line is the trace dual of \(\mathcal H_m\). Indeed the integral cup product on the torsion-free quotients is perfect, and the Bloch--Kato exponential and dual exponential are adjoint for this pairing. In the dual Tate bases their de Rham pairing is \((a,b)\mapsto\operatorname{Tr}_{L/\mathbb Q_p}(ab)\). \cite[Section~3]{bloch-kato-finite}

Put \(\beta=s-1-e-v_p((e-1)!)\). The valuation threshold for the inverse different is \(-s+1/(p-1)\); with \(m=e+1\) one has
\[
        \beta+A=-s+1/(p-1).
\]
Thus \(p^\beta\mathcal O_L\mathcal H_m\subseteq\mathfrak D^{-1}\). Every vector of \(p^\beta\mathcal O_L\) has integral trace pairing with \(\mathcal H_m\), proving the required inclusion at every odd prime.

For \(\ell=2\) and \(e>0\), the exponent in Lemma~\ref{lem:dyadic-integral-period} is independent of the cyclotomic level. Since \(s\ge2\), the required ball satisfies
\[
  2^{s-1-e-v_2((e-1)!)}\mathcal O_{F_s}
     \subseteq 2^{1-e-v_2((e-1)!)}\mathcal O_{F_s}
     \subseteq\exp^* H^1(F_s,\mathbb Z_2(-e)).
\]
\label{cor:cyclotomic-character-lifting-3}

Set
\[
 A_e=(e+2)s+v_\ell(e)+\lfloor e/(\ell-1)\rfloor+1,
\]
omitting the \(v_\ell(e)\) term at \(e=0\).
The bound for \(\alpha\) gives
\(\ell^{A_e}\alpha\in\ell^\beta\mathcal O_{F_s}\).
For positive \(e\) use
\(v_\ell(e!)=v_\ell(e)+v_\ell((e-1)!)\) and
\(e/(\ell-1)<\lfloor e/(\ell-1)\rfloor+1\).
For \(e=0\) it follows directly.
Subtract an integral cocycle with this exponential
from \(\ell^{A_e}z\). The difference has zero
exponential and represents the same divisible
class. Reassemble the scalar coordinates in the
coefficient lattice. This puts the character image
in the divisible finite condition after loss
\(C(e+1)+s+v_\ell(e+1)+A_e\).
Corestriction adds at most the original local degree
loss, giving the character allowance used above.
\end{proof}



\subsubsection*{Why the retained localization has an irreducible block}
\label{subsec:retained-galois}
\hypertarget{source-retained-galois}{}

The local estimates on the character summands leave a complementary Galois block. Its residual irreducibility is needed for the integral matrix algebra used in the congruence argument. The two tame type conditions test that irreducibility on the retained support.

\begin{paramdefinition}[Good theta point and the two tame type tests]\label{hyp:retained-type-tests}
For a candidate type at a theta point of Section~\ref{subsec:selected-cut}, let \(g\in\{g_0,g_1\}\) be a type prime and let \(r\) be
the rank of the complementary summand. At each place \(v\mid g\), require that the degree-\(r\) polynomial \(P_v\) have simple nonzero roots. Call the tuple good at this prime when it satisfies all three tests.
\begin{itemize}
\item For each possible constituent dimension, one place excludes it:
\[
 \forall a\in\mathbb Z\cap[1,r/2]\quad\exists v\mid g
 \quad\nexists Q\in\mathbb F_g[X]:\quad Q\mid P_v,\quad \deg Q=a .
\]
\item Away from the bad prime set for \((r,g)\), projection to the
prime-to-\(\ell\) parts is injective on the roots of each \(P_v\).
At every \(\ell\ne g\), those projected roots avoid the projected
labels of the character block and their inverses.
\item At one place there is an irreducible factor of prime degree
\(p\) with \(r/2<p\le r-3\).
\end{itemize}
Retain a pair of types when it is good at both primes. The bad prime set is
\[
 \{g\}\cup\{\ell\text{ prime}:\ g^d<\ell^{4v_\ell(g^d-1)}
                    \text{ for some }1\le d\le r\}.
\]
The projector in Definitions~\ref{def:localized-type-support} and~\ref{def:good-projector-support} retains exactly the jointly good labels within its theta support. Proposition~\ref{prop:two-type-primes} supplies type primes with disjoint exceptional sets. For a root of unity \(z\) of order
\(\ell^a b\), with \(\ell\nmid b\), its prime-to-\(\ell\) part
is \(z^e\), where \(e\equiv0\pmod {\ell^a}\) and
\(e\equiv1\pmod b\).  \label{cor:retained-galois-1}
\end{paramdefinition}

\begin{proposition}[The irreducible rank-complement]\label{prop:retained-rank-complement}
At a good theta point satisfying Definition~\ref{hyp:retained-type-tests},
the complement of the prescribed string is cuspidal of rank \(r\).
Its attached integral representation has absolutely irreducible
reduction at the coefficient prime \(\ell\). Choose a marked type
prime \(g\) with \(\ell\notin B_r(g)\), with \(B_r(g)\) as in Definition~\ref{def:type-exceptional-primes}. Write \(\bar W\) for the reduction over an algebraic closure of its residue field. For every open normal subgroup \(H\) with solvable quotient containing the inertia groups at the selected places over \(g\),
\[
 0\ne U\subseteq\bar W,\quad HU=U\quad\Longrightarrow\quad U=\bar W
\]
for every subspace \(U\).
\end{proposition}

\begin{proof}
\proofstep{Apply the common divisor-degree test to the possible isobaric partitions.}
First the parameter left after removal of the specified character
block is cuspidal of rank \(r\). Indeed, if its parts have ranks
\(a_j\) and multiplicities \(d_j\), the degree calculation is
\[
 r=\deg\!\left(\prod_j P_{v,j}^{d_j}\right)=\sum_j d_j\deg P_{v,j}\le\sum_j d_j a_j=r.
\]
Thus \(\deg P_{v,j}=a_j\) for every part at every place.
The labels have no repetitions, so every \(d_j=1\). If there were at least two parts, choose a smallest one of rank \(a\). Its factor would satisfy
\[
 1\le a\le r/2,\qquad P_{v,j}\mid P_v,\qquad\deg P_{v,j}=a\quad(\text{every }v).
\]
This contradicts the
first test.

\proofstep{At a good type prime, use the distinct prime-to-residue-characteristic inertia labels.}
Choose \(g\) for which \(\ell\) is outside the bad set. The attached
integral representation for this regular algebraic conjugate
self-dual cusp has the following inertia polynomials. In the
depth-zero dictionary a root \(z\) gives the corresponding tame
root of unity on a fixed generator \(T_v\). Transport to
\(\overline{\mathbb Q}_\ell\) followed by reduction kills its
\(\ell\)-part and is injective on the roots of order prime to
\(\ell\). For the prime-to-\(\ell\) order \(b\), these roots remain distinct because
\[
 z^b=1\quad\Longrightarrow\quad b z^{b-1}\ne0
 \quad\text{in residue characteristic }\ell.
\] The unramified
normalizing twist has value one on inertia. Thus the eigenvalues of
\(\bar\rho(T_v)\) are the images of the projected roots just
described. The depth-zero parameter has trivial wild inertia.
For a wild element the polynomial \((X-1)^r\) makes its residual image unipotent, hence of \(\ell\)-power order. Its order is also a power of \(g\), so
\[
 g\ne\ell\quad\Longrightarrow\quad\operatorname{ord}(\bar\rho(\text{wild element}))=1.
\]

Arithmetic Frobenius conjugates \(T_v\) to \(T_v^g\) modulo wild
inertia, since every marked type place splits completely and has
residue field \(\mathbb F_g\). Consequently Frobenius permutes the
residual eigenlines with the same orbit lengths as the roots of
\(P_v\) under \(z\mapsto z^g\). Projection commutes with this
power map and the roots remain distinct. On lines, taking the
inverse permutation gives the same orbit lengths.

\proofstep{Use the two permutation tests and then restrict to the auxiliary subgroup.}
The tame generator has separated the residual inertia lines. Thus an invariant subspace would select a set of lines stable under the permutation induced by Frobenius.

An invariant subspace over \(\overline{\mathbb F}_\ell\)
is a sum of the eigenspaces for \(T_v\), by the polynomial
projectors for its distinct eigenvalues. It is Frobenius stable.
Thus for each selected \(v\) its dimension and codimension lie in
the set
\[
     \left\{\sum_{C\in S}|C|:
          S\text{ a subset of the Frobenius cycles at }v\right\}.
\]
This set consists of the degrees of divisors of \(P_v\)
over \(\mathbb F_g\). The first test excludes a common degree
between \(1\) and \(r/2\), so no proper nonzero subspace
is possible. Scalar normalization preserves invariant subspaces.
At each split unramified Frobenius outside the finite level, denote the full polynomial by \(P_{\rm full}\) and the monic character factors by \(P_1,\dots,P_t\). The comparison with the cusp reads
\[
 P_{\rm full}/(P_1\cdots P_t)=P_{\rm cusp}.
\]
Chebotarev on a common finite quotient and Brauer--Nesbitt give
the same semisimplified reduction. Its irreducibility implies
irreducibility of the actual reduction.

Let \(H\) now be the open normal subgroup containing the selected
inertia groups, with solvable quotient. At the prime-degree
test place a minimal nonzero \(H\)-invariant space is a span of
coordinate lines of \(T_v\). Its translates are also minimal
and distinct ones have zero intersection. Full irreducibility makes these translates a partition of the eigenlines into \(h\) blocks of size \(b\), hence
\[
 \#\{\text{eigenlines}\}=hb=r.
\]
A power of the Frobenius permutation is a single
\(p\)-cycle: the other cycles have lengths smaller than \(p\).
For \(h,b>1\), the group preserving the blocks has order
\((b!)^h h!\), prime to \(p\) since
\[
 h,b\le r/2<p.
\]

This excludes those values of \(h,b\).

If \(b=1\), the action on the lines is transitive and factors
through the solvable quotient. It is primitive by the same
order argument applied to any invariant nontrivial partition.
A minimal nontrivial normal subgroup is elementary abelian
and transitive, and faithfulness makes that abelian action
regular. Hence the action is affine on \(\mathbb F_q^u\)
for a prime \(q\) with \(r=q^u\). Our \(p\)-cycle has a
nonempty affine fixed space, so
\[
 r-p=q^d.
\]
In the test
range \(3\le r-p<p\), giving \(d\ge1\). But then
\[
               p=q^d(q^{u-d}-1)
\]
would force \(q^d=p\), a contradiction. The remaining
partition has \(h=1\), so the original minimal subspace
is the whole space.

This proves the restriction assertion for \(H\). \label{cor:retained-galois-2}
\end{proof}

\subsubsection*{Constructing the two integral projections}

With residual irreducibility established, the next step is integral. The matrix algebra must contain projections that distinguish the complementary block from a character line after specialization.

\begin{proposition}[Integral projections at a retained localization]\label{prop:retained-integral-projections}
At a retained theta localization with absolutely irreducible
residual complement of rank \(r>1\), the local matrix algebra contains orthogonal idempotents \(e_B,e_1\). For the complement \(W\) and every character line \(L\) at a theta point, their ranks are
\[
 (\operatorname{rk}e_B|_W,\operatorname{rk}e_1|_W)=(0,1),\qquad (\operatorname{rk}e_B|_L,\operatorname{rk}e_1|_L)=(1,0).
\]
The marked corner operator may be chosen with simple spectrum at
the retained fibres, keeping its prescribed value at the selected
point and perturbing by any given power of the coefficient uniformizer
in the ideal vanishing there.
\end{proposition}

\begin{proof}
Let \(O\) be the complete coefficient DVR and \(F\) its fraction
field. At a localization containing a good point write \(X\)
for the finite set of systems and \(\mathcal T\subset O^X\)
for the local Hecke algebra. At every \(x\) there is a
continuous semisimple integral model \(R_x\), with characteristic
coefficients in \(\mathcal T\) as functions of \(x\).
Here an integral model means a stable lattice in the semisimple
rational representation. A compact image stabilizes such a lattice
by taking the span of the orbit of an initial lattice. The
coefficient splitting in the selected norm data realizes the
constituents over \(F\). At theta points use the rational
decomposition into the rank-\(r\) complement and the \(t\)
prescribed lines.

Let \(k\) be the residue field and \(\mathfrak m\) the unique maximal
ideal of \(\mathcal T\). Each residual evaluation
\[
 \overline{\operatorname{ev}}_x:\mathcal T\longrightarrow k
\]
is onto, since it contains the reduction of the constant scalars. Its kernel is therefore
\(\mathfrak m\), and the induced isomorphism \(\mathcal T/\mathfrak m\to k\)
is the same for all \(x\) because it fixes the residue scalars. Thus,
for every \(g\) and any two systems \(x,x'\), reduction of the
characteristic coefficients gives
\[
 \det(1-T\overline R_x(g))=\det(1-T\overline R_{x'}(g))
       \quad\text{in }k[T],
 \qquad
 \overline R_x^{\,\mathrm{ss}}\simeq \overline R_{x'}^{\,\mathrm{ss}}
\]
by Brauer--Nesbitt. This includes the systems outside theta.
In particular their distinct simple modules consist of the
absolutely irreducible rank-\(r\) module above and some
one-dimensional characters. Repeated residual characters are
allowed. Here \(r>1\) by the type tests.

\proofstep{The residual algebra and its idempotents.}
The integral residual action is obtained from the simultaneous matrix order
\[
 A=\sum_g\mathcal T\,(R_x(g))_x.
\]

It is a finite torsion-free \(O\)-algebra inside a product of
matrix rings. In its residue action on the distinct simple
modules the image is the product of their full matrix rings.
Indeed each action is onto by absolute irreducibility and the
density theorem, and inequivalent simples have different
maximal two-sided annihilators, which are pairwise comaximal.
This also follows by applying the semisimple algebra theorem
to the joint action. The simple actions for \(A\) here are the
ones for the residual Galois matrices, because the extra Hecke
scalars have common reduction. In the product of the rank-\(r\) matrix factor with the character factors, the desired residual projections are
\[
 (0,1,\ldots,1),\qquad (E_{11},0,\ldots,0).
\]


Let \(k=O/\varpi O\) and let \(J\) be the Jacobson radical
of \(A/\varpi A\). Its semisimple quotient surjects onto
the product acting on the indicated simples. Choose the
character projection and the matrix unit there and give
both zero on any unused simple factors. For an approximate idempotent \(a\), correct its defect by
\[
 d=a^2-a,\qquad b=3a^2-2a^3=a-(2a-1)d.
\]
The polynomial identities
\[
       b-a=-(2a-1)d,\qquad
       b^2-b=d^2(4d-3)
\]
show that the error is squared at each iteration.
This first lifts across the nilpotent ideal \(J\) and then
converges in the \(\varpi\)-adically complete ring \(A\).
After lifting the character projection \(e_B\), lift the
second idempotent \(e_1\) inside the corner with identity
\(1-e_B\). This preserves orthogonality.

For a stable integral fibre, an idempotent \(e\) gives
\(R_x=eR_x\oplus(1-e)R_x\) on its lattice. Reduction
therefore preserves its rank, and applying \(e\) on an
invariant composition series adds the ranks on the
residual factors. There are \(t\) one-dimensional factors \(L_i\)
counted with multiplicity, and one rank-\(r\) factor \(W\).
The chosen residual actions give
\[
 \begin{aligned}
 \operatorname{rk}(e_B R_x)&=\sum_{i=1}^{t}\dim e_B L_i
                                      +\dim e_B W=t,\\
 \operatorname{rk}(e_1 R_x)&=\sum_{i=1}^{t}\dim e_1 L_i
                                      +\dim e_1 W=1 .
 \end{aligned}
\]
This calculation depends only on the residual composition
factors and hence holds for every admitted fibre. At the marked theta point the \(\mathcal T\)-linear
matrix expressions restrict individually to each of the
rational constituents. Applying the same rank check on
each given stable lattice gives zero and rank one on the
rank-\(r\) summand, respectively, and identity and zero
on every scalar line. These are the characteristic-zero
projections asserted above.

\proofstep{An inertia element with the prescribed depth.}
At the split small primes the local host lies inside the compositum of the unramified field and the level-\(s\) cyclotomic field. Hence one can choose
\[
 g_0\in I_\ell,\qquad \epsilon_{\rm cyc}(g_0)=1+\ell^s,
\]
where \(I_\ell\) denotes host inertia and \(\epsilon_{\rm cyc}\) the cyclotomic character. When a type prime \(g_k=\ell\), start with
\(\sigma\) in inertia above \(\ell\) whose cyclotomic value
is \(1+\ell^{s_0}\), with \(s_0\ge1\) the prescribed local
depth.

In the closed procyclic subgroup generated by \(\sigma\), define
\[
 g_0=\operatorname{pr}_{\ell}(\sigma),
\]
where the projection is to the \(\ell\)-primary component in each finite quotient. The original
inertia group is closed, hence contains \(g_0\). At every
place above \(\ell\) whose inertia contains \(g_0\), its
image in tame inertia is trivial since that quotient has
prime-to-\(\ell\) order in every finite quotient. The
cyclotomic image of the closed cyclic group lies in the
pro-\(\ell\) group \(1+\ell^{s_0}\mathbb Z_\ell\).
Projection therefore preserves its value, in particular the
required exact depth
\[
 v_\ell(\epsilon(g_0)-1)=s_0.
\]

In particular the marked \(t\) eigenvalues are distinct,
since consecutive character values differ by the fixed
cyclotomic ratio of infinite order.

\proofstep{Preserving the value at the marked point.}
Lifting an idempotent separates the blocks in the integral order. The determinant calculation also requires a corner element with a prescribed specialization at the marked point.

We construct \(X_0\) in the corner of \(e_B\), specializing
to \(e_BR_\pi(g_0)e_B\) at \(\pi\) and with simple spectrum
on the corner at all \(x\). Over an algebraic closure of
\(F\) the evaluation algebra at each \(x\) is the product
of full matrix algebras on distinct irreducible constituents.
Their multiplicities are one by the weight test of
Section~\ref{subsec:selected-cut}. Its \(e_B\)-corner can
therefore realize a diagonal matrix with \(t\) distinct
entries. Over \(F\), separation by \(\mathcal T\) gives all
individual coordinate functions by interpolation. Hence the
rational span of
\[
  J_\pi e_B A e_B,\qquad
  J_\pi=\{f\in\mathcal T:f(\pi)=0\},
\]
gives all the corner directions away from \(\pi\).
Write \(V_x\) for the rational fibre at \(x\).

On the perturbation lattice consider the polynomial
\[
    F(s)=\prod_{x\ne\pi}\operatorname{disc}
       \det\bigl(T-(e_BR(g_0)e_B+s)|_{e_B V_x}\bigr).
\]
Each factor is nonzero as a polynomial, because the directions
at that point include matrices with distinct eigenvalues.
Thus the product is nonzero, including the empty product. This may be checked after
extending the fraction field by the preceding matrix
description. For any prescribed \(L\ge0\), the lattice scaled
by \(\varpi^L\) is Zariski dense in its rational span. Indeed a
polynomial vanishing on every integral linear combination of
a spanning list is zero there by repeated univariate
interpolation over the infinite ring \(O\). Writing \(D_x\) for the discriminant of the perturbed corner at \(x\), density supplies an integral perturbation \(s\) in the chosen scaled lattice with
\[
 s(\pi)=0,\qquad \prod_{x\ne\pi}D_x(s)\ne0,
\]
with the empty product interpreted as one.
This gives all three marked matrices and the perturbation in
the integral congruence module. The caller takes \(L=0\).

\proofstep{The finite choice of coefficient field.}
To arrange roots in the actual coefficients, first take the
system at a preliminary coefficient level which splits the
constituents and admits the theta decompositions. At each live
localization make the choices just proved and record the roots
of its corner polynomials at all points and all theta marks,
using the tied embedding in \(\mathbb C\). This is a finite set
for that localization. Write \(S_i\) for the finite root set at a live localization. At a coefficient place \(w\) take
\[
 S_w=\bigcup_{j=0}^{m}\ \bigcup_{i\mid(j,w)}S_i,
\]
with the inner union over the live localizations there. The construction of the coefficient rings
allows any prescribed finite set at each such place, along
with the elements required to identify the systems and
their integral Hecke actions at all larger levels.

At the resulting level every live localization lies over
a preliminary one at the same coefficient place. The Hecke
character and theta factorization descend along this
comparison. If \(k\subset k'\) are the two residue fields, cancellation of the common character factors at split Frobenius elements identifies the semisimplified rank-\(r\) complements:
\[
 (\bar W_{\rm old}^{\rm ss}\otimes_k k')^{\rm ss}\simeq \bar W_{\rm new}^{\rm ss}.
\]
Thus a live point upstairs lies over a live one downstairs. The marked roots occur among
the prescribed elements at its place. The integral Hecke
span extends by the identified coefficient map, so the
old matrices and their polynomial factorizations give the
marks there. Restricting to a cut restricts these functions
and matrices, keeps every good point and can only decrease
the finite set of root conditions.  It applies to the coefficients allowed
in the original choice of a norm. No estimate on the degree
of this splitting field enters the normalized lengths.
\label{cor:retained-galois-3}
\end{proof}

\subsubsection*{Inertia tests in the localized Hecke algebra}

The projections separate the two blocks. The congruence presentation also needs inertia relations inside the localized Hecke algebra, where different character blocks can have identical restrictions to inertia.

Use the same localization \(i\), coefficient DVR, character
\(\chi_1\), and the sets \(S,Q\) of its norm choice. All their
places are away from \(\ell\). Write \(\gamma_v\) for the
chosen lift of a topological tame generator. Unramifiedness of the cyclotomic character gives
\[
 \chi_c|_{I_v}=(\chi_1\epsilon^{1-c})|_{I_v}=\chi_1|_{I_v}.
\]

There are two different applications of the type dictionary.

At a type place \(v\in S\), put \(u=\chi_1(\gamma_v)\).
One needs separation of \(\bar u\) from the residual complement
even if \(\ell\) belongs to the bad prime set for this particular
type prime. Here \(\ell\ne g\). The local admissibility condition excludes equality of the prime-to-\(\ell\) part of any root of the complementary label and that of the specified character label. The test also holds for the inverse label. The
change to a conjugate marked place can invert either value in
the comparison.

If \(z\) is a projected complementary label and \(b\) the projected character label, the two exclusions give
\[
 z^\delta\ne b^{\delta'}\qquad(\delta,\delta'\in\{1,-1\}),
\]
by simultaneous inversion as necessary. Changing the tame generator
acts on the finite cyclic image by an automorphism and preserves
the exclusions. The residual Frobenius comparison identifies the
complement in the actual integral lattice with the scalar twist
of the good cuspidal one. This scalar is unramified at the type
places away from \(\ell\). Let \(D^{\rm res}_v\) be the residual characteristic
polynomial of the rank-\(r\) complement at \(\gamma_v\).
Thus
\[
      D_v^{\rm res}(\bar u)\ne0,\qquad
      \bar F_v(X)=(X-\bar u)^tD_v^{\rm res}(X).
\]
The multiplicity of \(\bar u\) is exactly \(t\) on every
system with the same residual characteristic polynomials.
This uses all the character factors and equality of the residual
characteristic polynomials.

At a retained point the type comparison gives the rank-\(t\)
inertia summand on which every \(g\in I_v\) acts by
\(\chi_1(g)\), and an invariant complement of rank \(r\).
Indeed the occurring full label has the character block of
Section~\ref{subsec:selected-cut}. The unramified normalization
preserves it. Its distinct tame labels give the inertia
characters and its monodromy vanishes. The local parameter supplies \(t\) scalar copies and reduction can only increase multiplicity. Thus
\[
 t\le\operatorname{mult}_{u}(F_v)\le\operatorname{mult}_{\bar u}(\overline F_v)=t,
\]
where \(F_v\) is the full inertia characteristic polynomial. Local--global
compatibility passes this decomposition to each Galois fibre
after scalar extension. The common complement is invariant for
the entire inertia group.

Form the full characteristic polynomial \(F_v(X)\in\mathcal T[X]\)
of \(R(\gamma_v)\) and divide by the monic polynomial
\((X-u)^t\). In every fraction-field fibre the remainder
vanishes by the decomposition. Injectivity into the product of fibres kills the remainder already over \(\mathcal T\). Thus
\[
 F_v(X)=(X-u)^t D_v(X),\qquad D_v\in\mathcal T[X]\ \text{monic of degree }r.
\]
Its evaluation
at a fibre is the characteristic polynomial on the complement.
The multiplicity computation shows that \(D_v(u)\) has nonzero
residue in the local Hecke algebra.

On each generic fibre write
\[
 V_x=V_{x,B}\oplus V_{x,C}
\]
for the invariant block and complement. Then
\[
 D_v(R_x(\gamma_v))|_{V_{x,B}}=D_v(u)\,\mathrm{id},
 \qquad D_v(R_x(\gamma_v))|_{V_{x,C}}=0,
\]
the latter by Cayley--Hamilton. On the rank-\(t\)
space every \(R(g)-\chi_1(g)\) vanishes. As the first space
is stable under \(R(g)\), the product
\[
       D_v(R(\gamma_v))(R(g)-\chi_1(g))=0
       \qquad(g\in I_v)
\]
holds integrally on the whole family. This constructs the single
polynomial with unit character value needed for the row cocycles.
The unit test uses the residual multiplicity at the localization.
Squarefreeness in characteristic zero by itself only controls
the rational eigenspaces. \label{cor:retained-galois-4}

\begin{lemma}[Inertia at the auxiliary and order-fifteen places]\label{lem:retained-inertia-types}
At \(Q\) and at the additional distinguished places of \(S\),
the zero-monodromy and character conditions give, on every characteristic-zero fibre,
\[
 R_x|_{I_v}\simeq \mathbf1^{\oplus r}\oplus(\chi_1|_{I_v})^{\oplus t}.
\] The power of the chosen inertia generator representing
\(R_x(g)\), for \(g\in I_v\), can be chosen independently of \(x\).
At the distinguished \(S\)-places \(\zeta_v-1\) is a unit, including
in coefficient characteristics three and five.
\end{lemma}

\begin{proof}
At \(Q\) and at the additional distinguished places of \(S\),
the local parameter has zero monodromy and inertia polynomial
\((X-1)^r(X-\chi_1(g))^t\) for each \(g\in I_v\).
These are the unramified rank-\(r\) and prescribed character
types in the norm tuple. Zero monodromy gives finite inertia
image for each fibre by the Weil--Deligne comparison away from
\(\ell\). The displayed polynomial identifies the representation
of inertia with \(1^r\oplus\chi_1^t\): both representations
factor through a finite quotient and their characteristic
polynomials agree there, so Maschke and Brauer--Nesbitt apply
in characteristic zero. In particular one decomposition per
fibre works for all inertia elements. Merely checking the
spectrum of a single matrix would not prove this assertion.

At \(v\in Q\) the first twist factor is unramified and the ramified auxiliary factor has
\[
 \#\chi_1(I_v)=\ell^a\qquad(a\ge0).
\]
At the extra
distinguished prime of \(S\) the auxiliary character is
unramified, and the tame factor or its inverse has inertia
image of order \(15\), with residue norm \(1\bmod 15\).
The construction uses the second distinguished prime when
the first equals \(\ell\). In both cases wild inertia is
trivial on the character: the wild pro-residue-characteristic
group has no quotient of the indicated order. Hence the image
of inertia is the finite cyclic group generated by
\(\zeta_v=\chi_1(\gamma_v)\).

At \(Q\) ramification forces
\(\zeta_v\ne1\). For \(g\in I_v\) choose an integer \(a\)
such that
\[
 \chi_1(g)=\zeta_v^a.
\]
In the simultaneous
decomposition just proved,
\[
 R_x(g)=R_x(\gamma_v)^a
\]
at every \(x\), using the same
\(a\) for them all. At the additional \(S\)-places
\(\zeta_v\) is primitive of order \(15\).
The identity \(\Phi_{15}(1)=1\) makes \(\zeta_v-1\) a unit:
subtract \(\Phi_{15}(\zeta_v)=0\) and factor the difference.
This argument also covers coefficient characteristics three
and five. The common exponent of the tame generator and the unit at the additional places give the two inertia controls needed in the ensuing congruence presentation. \label{cor:retained-galois-5}
\end{proof}

\subsubsection*{A determinant supported on the theta points}
\label{subsec:congruence-determinant}
\hypertarget{source-congruence-determinant}{}

The projections and inertia relations describe the relevant blocks of the integral Galois matrix algebra. A determinant in the Hecke algebra will convert their cohomology lengths into a congruence estimate supported on the theta points. Fix one coefficient DVR \(O\) with
fraction field \(F\). In this calculation write
\(\operatorname{len}\) for the unnormalized \(O\)-length.
Use a finite parameter set \(X\), the theta subset \(Z\), a
Hecke algebra \(\mathcal T\subset O^X\) and a good point
\(\pi\in Z\). Write \(I_Z\) for the ideal of elements of
\(\mathcal T\) vanishing outside \(Z\). The congruence
length before normalization is
\(\operatorname{len}(O/I_Z(\pi))\).

\begin{established}[Ordered corner and block evaluation]\label{hyp:theta-marked-corner}
The auxiliary parameter set \(X'\) records an ordering of
the eigenvalues of a marked \(t\)-dimensional corner. Its
coefficient algebra \(\mathcal T'\subset O^{X'}\) has a
basis \(b_j\) over the pullback of \(\mathcal T\) along
\(p:X'\to X\). For each \(x\in X\setminus Z\), evaluation gives an injection
\[
 O^{t!}\lhook\joinrel\longrightarrow O^{p^{-1}(x)},\qquad
 (c_j)_j\longmapsto\bigg(\sum_j c_j b_j(x')\bigg)_{x'\in p^{-1}(x)} .
\] At the chosen
point of \(X'\) above \(\pi\) the matrix algebra
\(\mathcal A\) has a specialization into matrices over
\(O\). Its image is block triangular with one block of
size \(r\) and \(t\) blocks of size one, and each block
map \(\rho_k:\mathcal A\to M_{d_k}(O)\) is surjective.
The determinants of square matrices over \(\mathcal A\),
computed on the full matrix spaces at all points of
\(X'\), belong to \(\mathcal T'\).
\end{established}
\begin{proof}
\proofstep{Integral coefficients of the matrix determinants.}
These integral properties can be seen from the marked
corner construction. Let \(\mathcal A_0\) be the
\(\mathcal T\)-span of the Galois matrices \(r(g)\),
which is an algebra since products of Galois matrices
are Galois matrices. The coefficients of the
characteristic polynomials of \(r(g)\) are in
\(\mathcal T\). For any finite number of block rows,
expand a matrix over \(\mathcal A_0\) as a scalar linear
combination of \(e_{ij}\otimes r(g)\), with \(e_{ij}\)
the elementary matrices on the block rows. To see that
its characteristic coefficients are in \(\mathcal T\),
give each summand \(A_i=e_{ij}\otimes r(g)\) an
indeterminate \(X_i\). For a primitive cyclic word \(w\), let
\(A_w\) be its ordered matrix product and \(X_w\) its monomial.

In the completed polynomial ring in characteristic zero,
\[
 \log\det(1-\textstyle\sum_iX_iA_i)
 =-\sum_{a\ge1}\frac{\operatorname{tr}((\sum_iX_iA_i)^a)}a
 =-\sum_{w\ {\rm primitive}}\sum_{b\ge1}
                    \frac{X_w^b\operatorname{tr}(A_w^b)}b .
\]
Indeed \(w^b\) contributes \(|w|\) rotations at length \(b|w|\).
Exponentiating gives \(\prod_w\det(1-X_w A_w)\), with only
finitely many factors in each degree.
A product \(A_w\) is zero, has one off-diagonal
block, or has a single diagonal block carrying a
Galois matrix. The individual factors are therefore
\[
 \det(1-X_w A_w)\in\mathcal T[X_w].
\]
For the zero and off-diagonal cases the determinant equals one. After the identity has been
proved in the fraction fields, its coefficient
expression is integral and uses no division in
\(\mathcal T\). Substitution of the scalar coefficients
proves the assertion, also over any \(\mathcal T\)-algebra
of scalar functions.

\proofstep{The ordered corner roots and their splitting algebra.}
The corner operator in \(\mathcal A_0\) is
\(e_Br(g_0)e_B+s\), where \(e_B\) is the marked
rank-\(t\) idempotent and \(s=e_Bse_B\) is the
perturbation with \(s(\pi)=0\). It is zero on the
complement of the corner. The previous argument
places its corner characteristic polynomial in
\(\mathcal T[X]\), with \(X\) here a polynomial
indeterminate. The marked perturbation gives distinct
roots at every fibre in the tied splitting field.
The successive monic quotients adjoining the ordered roots have degrees \(t,t-1,\ldots,1\). Freeness at each step gives
\[
 \operatorname{rank}_{\mathcal T}(\text{splitting algebra})=\prod_{a=1}^{t}a=t!.
\]

Over each rational fibre this algebra is the
product of \(t!\) copies of the coefficient field,
by the Chinese remainder theorem. The free basis
therefore stays independent when evaluated at all
orderings in each fibre. This also gives an embedding
of the splitting algebra into \(O^{X'}\), since the
roots of integral monic polynomials are integral.
Take \(\mathcal A\) to be the span of the Galois
matrices over this splitting algebra. The determinant
argument applies with these coefficients. All
matrices are integral, so \(\mathcal A\) is finite
over \(O\). At the chosen fibre the marked
decomposition gives the triangular blocks, and
the residual full matrix span of the rank-\(r\)
constituent gives surjectivity on that block by
Nakayama's lemma. The scalar blocks have surjectivity
from the coefficient ring. These algebraic steps
use the marked idempotents and the perturbation
on the given family, constructed in Proposition~\ref{prop:retained-integral-projections}.
\end{proof}

\begin{proposition}[The congruence ideal on the theta locus]
\label{prop:theta-determinant-length}
Take the ordered corner and block maps of Lemma~\ref{hyp:theta-marked-corner}. Let \(M\) be the finite
presentation with the column relations in the proof. If these
relations have the prescribed finite column dimensions away from
\(Z\), then
\[
 \operatorname{len}(O/I_Z(\pi))\le
       \sum_k\operatorname{len}\operatorname{Hom}_{\mathcal A}
                          (M,(F/O)^{d_k}).
\]
The \(k\)-th action is via \(\rho_k\), with the same coefficient
ring and chosen point in all the summands.
\end{proposition}

\begin{proof}
Lemma~\ref{hyp:theta-marked-corner} supplies the splitting algebra, its fibre basis and the full block evaluations.

Write
\[
 B=e_Br(g_0)e_B+s
\]
on the corner, and let its ordered
roots be \(\alpha_0,\ldots,\alpha_{t-1}\). The operator separating
the first eigenline is
\[
 E=\prod_{i>0}(B-\alpha_i e_B),\qquad
 E|_{V_{\alpha_i}}=
 \begin{cases}\Delta\,{\rm id},&i=0,\\0,&i>0,\end{cases}
 \qquad\Delta=\prod_{i>0}(\alpha_0-\alpha_i).
\]
Distinctness of the roots gives \(E^2=\Delta E\) on every fibre.
At \(\pi\) the roots are \(\chi(g_0)\varepsilon(g_0)^{-i}\), where
\(\varepsilon\) is the cyclotomic character. The
marked value is \(\varepsilon(g_0)=1+\ell^s v\)
with \(v\in\mathbb Z_\ell^\times\), at \(s\ge2\) for
\(\ell=2\) and \(s\ge1\) for \(\ell=3\).
The binomial calculation gives
\[
     v_\ell(\Delta_\pi)=(t-1)s+v_\ell((t-1)!).
\]
\label{cor:congruence-determinant-1}

\proofstep{The column relations in a finite presentation.}
Here is the module used in the proof. Write \(U_g\) for the
Galois operator after dividing by the first character.
Use \(E,E_1\in\mathcal A\) with \(E^2=\Delta E\) and
\(E_1^2=E_1\), with \(\Delta\in\mathcal T'\) central.
The two operators have rank one over \(F\) at each
point of \(X'\) above \(X\setminus Z\), and
\(\Delta\) is then nonzero. Put
\(Y=\sum_g\mathcal A (U_g-1)E\subset\mathcal A\).

At each column place \(u\) with decomposition subgroup
\(D_u\), adjoin a generator \(n_u\), \(j_u\) generators
\(s_{u,a}\) and a generator \(s_{u,0}'\) to \(Y\).
Impose the relations
\[
 \begin{gathered}
 E n_u=0,\qquad (E_1-1)s_{u,a}=0,\qquad
 (E-\Delta)s_{u,0}'=0,\\
 (U_g-1)E+(U_g-1)n_u
     =\Delta\left(\sum_{a=1}^{j_u}K_{u,a}(g)s_{u,a}
                  +K_{u,0}(g)s_{u,0}'\right)\qquad(g\in D_u).
 \end{gathered}
\]
Here the first term in the last relation is the
element of \(Y\). The functions \(K_{u,a},K_{u,0}\)
are the matrix cocycles in the column data. Denote the
quotient by \(M\). Since \(\mathcal A\) is left
Noetherian, \(Y\) is finitely generated and \(M\)
has a finite presentation.

\proofstep{A kernel vector on each complementary fibre.}
The candidate determinant has integral coefficients. To lie in the theta congruence ideal it must vanish when evaluated on a complementary fibre, which can be tested by producing a vector in its kernel there.

Consider a rational evaluation space \(V\) at \(x'\)
over \(X\setminus Z\). Fix a nonzero vector \(\eta=Ew\)
in the image of \(E\) and a nonzero \(s\) in the image
of \(E_1\). The two conditions on a vector in the first image are incompatible:
\[
 \operatorname{im}E\ \cap\ \bigcap_g\ker(U_g-1)=0.
\]
At \(u\), the
column condition says that the functions
\(K_{u,a}(\cdot)s\) and \(K_{u,0}(\cdot)\eta\)
lie in a space containing the coboundaries and of
dimension at most \(j_u+\dim V^{D_u}\) modulo
coboundaries. Here invariants and coboundaries use \(U\).
If \(V^{D_u}\ne0\), take \(v_u\ne0\) invariant there
and all scalars zero. Otherwise \(V^{D_u}=0\). Let \(L_u\) be the span of the classes of the \(j_u+1\) functions in the quotient by coboundaries. The column bound gives
\[
 \dim L_u\le j_u<j_u+1.
\]

A nontrivial linear dependence in this quotient gives a nonzero tuple \((v_u,(a_{u,a})_a,a_{u,0})\)
such that
\[
 (U_g-1)v_u
      =\sum_a a_{u,a}K_{u,a}(g)s+a_{u,0}K_{u,0}(g)\eta
           \qquad(g\in D_u).
\]
If at some \(u\) we can take \(Ev_u=0\), map \(Y\)
and all other places to zero and map the generators
at \(u\) to \(\Delta v_u,a_{u,a}s,a_{u,0}\eta\).
This satisfies the relations and gives a nonzero map
\(M\to V\). In the other case rescale each tuple so
that \(Ev_u=\eta\). Map \(y\in Y\) to \(yw\) and
the local generators to
\(\Delta v_u-\eta,a_{u,a}s,a_{u,0}\eta\).
The relations follow from \(E\eta=\Delta\eta\)
and the equation at \(D_u\). The map is nonzero
on \(Y\) because \(\eta\) is moved by some \(U_g\).
Thus
\[
 \operatorname{Hom}_{\mathcal A}(M,V)\ne0
\]
at every evaluation point over \(X\setminus Z\).

\proofstep{Blockwise Smith lengths at the chosen theta point.}
Present \(M\) with \(g\) generators and a finite
matrix of relations \(P\), written so that applying
\(\operatorname{Hom}_{\mathcal A}(-,V)\) identifies it
with the kernel of \(P\) acting on \(V^g\). At the
chosen point above \(\pi\), use the block targets
\[
 H_k=(F/O)^{d_k}
\]
with action via \(\rho_k\).
If their Hom lengths are finite, the expanded relation
matrix \(P_k\) in each block has full column rank.
Let \(\varpi^{a_{k,i}}\) be its elementary divisors, with
\(1\le i\le gd_k\). Since
\[
 (F/O)[\varpi^a]\simeq O/(\varpi^a),
\]
the kernel on the divisible target and its determinantal ideal satisfy
\[
 \begin{aligned}
 \operatorname{len}\operatorname{Hom}_{\mathcal A}(M,H_k)
     &=\sum_{i=1}^{gd_k}a_{k,i}=:s_k,\\
 I_{gd_k}(P_k)&=(\varpi^{s_k}),
 \end{aligned}
\]
where \(I_{gd_k}\) is generated by the maximal minors. Thus one such
minor attains valuation \(s_k\).

For each \(k\), there exists a contraction of the rows
giving a square \(gd_k\)-matrix with determinant
valuation \(s_k\). Surjectivity of \(\rho_k\) lifts
such a contraction to a rectangular matrix \(C_k\)
over \(\mathcal A\). Use the matrix
\(C(\mathbf u)=\sum_j u_j C_j\) over a polynomial
ring with indeterminates \(u_j\).
In block \(k\), the determinant of \(CP\), divided
by a uniformizer to the power \(s_k\), is a polynomial
over \(O\). Indeed Cauchy--Binet expresses the numerator as
\[
 \det(C(\mathbf u)_kP_k)=
   \sum_{|J|=gd_k}\det(C(\mathbf u)_{k,*,J})\det(P_{k,J,*}).
\]
Each term is divisible by \(\varpi^{s_k}\) in \(O[\mathbf u]\).
Its reduction is nonzero since setting \(u_k=1\)
and the other \(u_j=0\) gives a unit. Therefore the
product of the reductions is a nonzero polynomial,
and some coefficient is a unit.

\proofstep{A determinant coefficient in the original congruence ideal.}
The same characteristic coefficient argument applied
over \(\mathcal T'[\mathbf u]\) puts the determinant of
\(C(\mathbf u)P\) in \(\mathcal T'[\mathbf u]\). At every
point above \(X\setminus Z\) the nonzero maps \(M\to V\)
give a kernel vector, so the determinant vanishes
as a polynomial at those points. At the data at
\(\pi\) the determinant is the product of the block
determinants by triangularity. Take a coefficient
\(f'\in\mathcal T'\) corresponding to a unit coefficient
of the reduced product. Its valuation at this point
is \(\sum_k s_k\), and it vanishes off the theta subset.
Write
\[
 f'=\sum_j(f_j\circ p)b_j.
\]

For every
\(x\notin Z\), its evaluations give a vanishing vector:
\[
 0=(f'(\xi))_{\xi\in p^{-1}(x)}
   =\sum_j f_j(x)(b_j(\xi))_{\xi\in p^{-1}(x)} .
\]
Independence on that fibre forces \(f_j(x)=0\) for
all \(j\). Thus \(f_j\in I_Z\), and the values \(b_j(\pi')\)
are integral in the common coefficient ring, so \(f'(\pi')\)
lies in \(I_Z(\pi)\). Its valuation gives
\[
 \operatorname{len}(O/I_Z(\pi))
 \le \sum_k\operatorname{len}\operatorname{Hom}_{\mathcal A}(M,H_k).
\]
If a Hom length is infinite this inequality already
holds in extended lengths. The empty generator case
has determinant one. Division by the ramification
index of \(O/\mathbb Z_\ell\) gives the normalized
inequality.  Choose a residue field
of size greater than \(g\sum_k d_k\) to specialize the
contraction simultaneously. Coefficient extraction gives
the same ideal bound over any residue field. \label{cor:congruence-determinant-2}

\end{proof}

\subsubsection*{Choosing the column cocycles}

The determinant argument requires actual columns in the integral presentation. A finite quotient specifies their evaluations only up to torsion, so we choose lifts and bound the exponent of a relation among them.

\begin{lemma}[Column lifts and their torsion tests]
\label{lem:congruence-column-cocycles}
Under the finite-class lift bound \(e'\) on the designated local
lattices, write \(c_a\) for a lift of the \(a\)-th finite quotient basis class and \(c_*\) for the extra column. Their evaluations are
\[
 c_a\longmapsto\ell^{e'}\beta_a\ (1\le a\le j_u),\qquad c_*\longmapsto\ell^{e'}\Delta_\pi\quad\text{in the unramified trivial block}.
\] These columns
satisfy the finite-class dimension bound away from \(Z\) and their
torsion dependences in the \(W\)-block are killed by \(\ell^{e'}\).
\end{lemma}

\begin{proof}
\proofstep{Integral finite bases at a column place.}
At the selected point let \(T\) be the integral lattice
in the \(W\)-block. Choose a basis \((\beta_a)_a\) of
the free quotient of \(H^1_f(E_u,T)\). This quotient
is the lattice of integral finite classes in
\(H^1_f(E_u,W)\). It spans the rational finite group
since a continuous cocycle on the compact local Galois
group has bounded image and can be scaled into \(T\).
Thus its rank is determined on the rational finite part:
\[
 \operatorname{rank}_O\bigl(H^1_f(E_u,T)/\mathrm{tors}\bigr)=\dim H^1_f(E_u,W)=j_u.
\]


Use the left ideals \(\mathcal A E_1,\mathcal A E\)
with local action by \(U_g\). Rationally these are
subrepresentations of finite direct sums of the full
twisted family, by embedding matrices as columns in
the evaluation spaces. In block \(0\), choose a
generator \(v_1\) of the image of \(E_1\) on \(T\).
Evaluation on \(v_1\) maps \(\mathcal A E_1\) onto
\(T\), because the block image is the full matrix
algebra. Evaluation of \(\mathcal A E\) on the
trivial block has image \(\Delta_\pi O\), on which
the twisted action is trivial. The two block evaluations are thus surjections of stable lattices
\[
 \mathcal A E_1\twoheadrightarrow T,\qquad \mathcal A E\twoheadrightarrow\Delta_\pi O.
\]

The twisted Galois action on the second target is trivial.
The first column lift produces finite cocycles in
\(\mathcal A E_1\) with images representing
\(\ell^{e'}\beta_a\), using lifts of the \(\beta_a\)
to integral finite classes. Use these as \(K_{u,a}\).
The second lift uses the unramified class in
\(H^1_f(E_u,\Delta_\pi O)\) with value \(\Delta_\pi\)
on arithmetic Frobenius. This class is finite since
the trivial representation's rational finite part
is the unramified line. Its lift \(K_{u,0}\) has image
with Frobenius value \(\ell^{e'}\Delta_\pi\). This
value holds on every lift of that Frobenius because
coboundaries of the trivial module vanish.

\proofstep{The local dimension at every evaluation.}
At an evaluation point \(x'\), a vector fixed by \(E_1\)
and an eigenvector \(E\eta=\Delta\eta\) give the equivariant
evaluation maps \(a\mapsto as\) on \(\mathcal A E_1\) and
\[
 a\mapsto a\eta
\]
on \(\mathcal A E\), respectively
(\(E_1s=s\)). Functoriality sends their finite cohomology
classes to finite classes. The negative Hodge jumps of the full
twisted family are constant at the column points. Removing the
cyclotomic characters at \(\pi\) removes no negative jumps.
As \(H^0(E_u,W)=0\), their common total over embeddings is
\(j_u\). The finite-part dimension formula at \(x'\) is therefore
\[
        \dim H^1_f(E_u,V)=j_u+\dim V^{D_u}.
\]
This is the local dimension used for the finite column cocycles.

\proofstep{Torsion relations in the chosen basis.}
The rational dimension leaves the exponent of integral torsion relations undetermined.

Test a torsion relation on the \(W\)-block. Write a linear combination
of the evaluated column cocycles with coefficients
in \(F/O\), and lift those coefficients to \(F\).
If the combination is a coboundary modulo \(T\),
the rational finite class \(\sum_a y_a\ell^{e'}\beta_a\)
comes from \(H^1(E_u,T)\) by exactness for
\(0\to T\to W\to W/T\to0\).
It therefore comes from \(H^1_f(E_u,T)\) by the
definition of the integral finite part. Because
the \(\beta_a\) form its free basis, each
\(y_a\ell^{e'}\) belongs to \(O\). This proves the
killing exponent for such a dependence.   \label{cor:congruence-determinant-3}

\end{proof}

\subsubsection*{The Selmer length comparison for this presentation}
\label{subsec:congruence-selmer}

We can now evaluate the presentation at the chosen theta point. Each block contributes a Selmer length, while the valuation of the specialized determinant records the cost of imposing the integral conditions.

At the chosen point let \(\Delta_\pi\ne0\) be the
specialization of \(\Delta\) and put \(v_\Delta
=v_\ell(\Delta_\pi)\). From now on lengths are
normalized by \(v_\ell(\ell)=1\).
The block targets are \(H_0=(F/O)^r\) with the
action of \(W\), \(H_1=F/O\) with trivial action,
and \(H_k=(F/O)(1-k)\) for \(2\le k\le t\).
Write \(P_k=\operatorname{Hom}_{\mathcal A}(M,H_k)\).
\begin{proposition}[From block lengths to Selmer kernels]
\label{prop:congruence-block-selmer-length}
With the integral presentation and column lifts of
Lemma~\ref{lem:congruence-column-cocycles}, the normalized sum of
block Hom lengths in Proposition~\ref{prop:theta-determinant-length}
is bounded by the global kernel lengths and the local and
presentation costs below. The three global kernels are
the rank-\(r\) finite kernel, the unramified first character and
the strict classes in twists \(1-k\) for \(2\le k\le t\).
If \(c_k\) is the normalized contribution charged to coboundaries in block \(k\), then
\[
 c_1=0,\qquad c_k\le2d_kv_\Delta\quad(k\ne1).
\]
\end{proposition}

\begin{proof}
\proofstep{The map to continuous cocycles.}
A homomorphism \(\varphi\) supplies the cocycle
\[
    c_\varphi(g)=\varphi((U_g-1)E).
\]
The cocycle identity follows by multiplying \(U_g U_h\).
In the integral presentation all the matrix functions
are continuous, and the orders and their modules are
finite over \(O\). A map to \(H_k\) kills a common
power of the uniformizer times the module because the
images of a finite set of generators are torsion.
Thus \(c_\varphi\) is continuous.

\proofstep{Its kernel on the local generators.}
First consider the map from \(P_k\) to cocycles before
taking cohomology. Its kernel consists of maps vanishing
on \(Y\). They are determined by the images of the local
generators, subject to the relations of \(M\).
The specializations on the target blocks satisfy
\[
 E|_{H_0}=0,\qquad E|_{H_1}=\Delta_\pi,\qquad E|_{H_k}=0\quad(k\ge2).
\]
The idempotent \(E_1\) has rank one
in block \(0\) and is zero on the character blocks.
The extra column cocycle \(K_{u,0}(g)\) lies in the
left ideal generated by \(E\).

\begin{itemize}
\item In block \(0\) the column lift input says that
a linear combination \(\sum_a K_{u,a}(g)s_a\) with
\(s_a\in E_1H_0\) which is a coboundary on \(D_u\)
forces \(\ell^{e'}s_a=0\) for every \(a\).
For a map vanishing on \(Y\) the combination with
coefficients \(\Delta_\pi\varphi(s_{u,a})\) is a
coboundary by the relation for \(D_u\). The tuple of images of the \(s_{u,a}\) lies in
\[
   (E_1 H_0)[\ell^{e'}\Delta_\pi]^{\,j_u},
\]
of normalized length at most \(j_u(e'+v_\Delta)\).

For a fixed tuple, differences of choices of \(n_u\) lie in
\(H_0^{D_u}\), of length \(h_W(u)\), and the remaining
image of \(s'_{u,0}\) lies in \(H_0[\Delta_\pi]\), of length
at most \(rv_\Delta\). Kernel plus image for projection onto
the tuple therefore gives
\(j_u(e'+v_\Delta)+h_W(u)+rv_\Delta\).

\item In block \(1\) the matrix \(K_{u,0}\) has some
value on \(D_u\) acting by
\(\ell^{e'}\Delta_\pi\) times a unit, the first-character
column input. The relation for that element gives
\[
 \ell^{e'}\Delta_\pi^2\varphi(s_{u,0}')=0
\]
for a
map vanishing on \(Y\). The image of \(n_u\) is killed
by \(\Delta_\pi\) and that of \(s_{u,a}\) is zero.
The resulting cost is at most \(e'+3v_\Delta\).

\item In block \(k\ge2\) the \(s_{u,a}\) vanish and
\(\Delta_\pi\varphi(s_{u,0}')=0\). The relation for
\(D_u\) forces \(\varphi(n_u)\in H_k^{D_u}\), since
\(K_{u,0}\) acts by zero. With the scalar invariant
bound \(h_0\) the cost is \(h_0+v_\Delta\).
\end{itemize}

\proofstep{The cost of quotienting by coboundaries.}
Coboundaries change the length differently on each block.

Pass next to cohomology. In block \(1\) the coboundaries
are zero. In the other blocks the central scalars \(a_g\) of the marked corner satisfy
\[ E=\sum_g a_g U_g,\qquad \sum_g a_g=\Delta_\pi\quad\text{at the chosen point}. \]
If \(c_\varphi(g)=(U_g-1)v\), the element
\[
    y=\sum_g a_g (U_g-1)E=(E-\textstyle\sum_g a_g)E
\]
has image \(-\Delta_\pi v\) in \(H_k\), since \(E\)
acts by zero on \(H_k\). The algebra identity
\(Ey=\Delta y\) then forces \(\Delta_\pi^2 v=0\).
The space of image cocycles that are coboundaries
has length at most \(2d_k v_\Delta\).

\proofstep{Relative Selmer kernels and local correction terms.}
Finally let \(C_k\) be the image of \(P_k\) in
\(H^1(E,H_k)\). At each local place impose a target
subgroup \(F_{k,u}\) of local cohomology. Suppose
the restrictions of \(C_k\) are in subgroups
\(B_{k,u}\supset F_{k,u}\) and that the classes also
lie in a global subgroup \(Q_k\). The kernel of
\[
   C_k\longrightarrow\bigoplus_u B_{k,u}/F_{k,u}
\]
lies in the subgroup of \(Q_k\) satisfying all the
local conditions. Thus its global length bound and
the sum of the local colength bounds give a length
bound for \(C_k\). In the split comparison, the
column images already satisfy \(F_{k,u}\).

In
block \(0\) the row quotient has the colength
bound proved above, and the global kernel has length
at most \(s_\pi\). In the character blocks the row and
tame quotients have bounds \(([E_w:\mathbb Q_\ell]+2)c_B\)
and \(v_\ell(\zeta_v-1)\). If \(K_1\) is the global kernel in block one, unramified class field theory gives
\[
 K_1\hookrightarrow\operatorname{Hom}(\operatorname{Cl}(E),F/O),\qquad \operatorname{len}_{\rm norm} K_1\le v_\ell(h_E).
\]
In the higher blocks the local
conditions defining this kernel are strict at \(\ell\)
and unramified elsewhere. Its length is bounded by the
strict group \(\operatorname{Sha}^1(E,(F/O)(1-k))\).
The normalization and the row and tame restriction arguments
follow below.

\proofstep{The sum of the normalized block costs.}
Let \(k_c\) be the number of columns. Summing these
bounds over all blocks gives the expression used in
the split-row estimate. The terms from the rank-\(r\)
and character blocks contribute
\[
 \sum_u\{j_u(e'+v_\Delta)+h_W(u)+e'\},
\]
and the \(t-1\) higher blocks contribute at most
\(tk_ch_0\). The remaining denominator terms from
the presentation and the coboundaries are bounded by
\(2k_cv_\Delta+4R(k_c+1)v_\Delta\). Indeed the dimensions \(d_k\) of the blocks sum to \(R\), so their contribution is
\[
 \sum_k d_k k_c v_\Delta=R k_c v_\Delta,
\]
and the coboundaries cost at most \(2Rv_\Delta\).

The row and tame costs of the
character blocks occur at most \(t\) times. Adding the
global kernels accounts for \(s_\pi\), the class
number and the strict cohomology sum. Pairing the row
cost for \(W\) at \(w\) with its conjugate column
\(u\) and using \(j_u\le [E_u:\mathbb Q_\ell]m\)
gives the formula below.
\label{cor:congruence-determinant-4}
\end{proof}

\subsubsection*{Changing the frame at the selected point}
\label{subsec:split-selmer}
\hypertarget{source-split-selmer}{}

To identify those abstract block lengths with the prescribed local conditions, choose integral coordinates adapted to the decomposition at the good point. The surrounding family still includes the points off the theta locus that enter its congruence ideal.

Use the integral family \(R_x\) at the fixed Hecke
localization, including the points outside the theta subset.
At the good theta point \(\pi\) the given rational decomposition
is an isomorphism with \(\rho_\pi\oplus\bigoplus_{k=1}^t\chi_k\).
Here \(\rho_\pi\) is the specified integral rank-\(r\) representation
and \(\chi_k/\chi_1=\epsilon^{1-k}\) for the cyclotomic character
\(\epsilon\). Replace \(R_\pi\) by the direct sum in these
integral bases and keep all other matrices at their own points.
\begin{lemma}[Change of frame in the family algebra]
The rational change of basis identifies the two integral matrix
spans, their marked corners and congruence ideals. At the chosen point \(\pi\), evaluation on the complement gives
\[
 \operatorname{ev}_{\pi,W}(\mathcal A)=M_r(O),
\]
where \(\mathcal A\) is the normalized integral span.
\end{lemma}
\begin{proof}
If \(C_x\) changes the frame at \(x\), conjugation sends each generator
\(R_x(g)\) to \(R'_x(g)\). Thus it carries
\[
 \left(\sum_j a_j(x)R_x(g_j)\right)_x
 \longmapsto\left(\sum_j a_j(x)R'_x(g_j)\right)_x .
\]
The coefficients are the very same Hecke functions. Applying
inverse conjugation proves equality of the integral spans.
Characteristic coefficients are unchanged. Row embeddings and
column quotients transport by rational conjugation, with the
specified lattices of \(\rho_\pi\) and its twisted dual as targets.
The resulting finite family is continuous at each point.
\label{cor:split-selmer-1}
The marked data before the change of frame consist of
orthogonal summand projections at \(\pi\), extended as integral
idempotents in the family span, and a corner matrix \(X_0\) with
distinct roots in every fibre. The fibrewise ranks and the two specializations are
\[
 \operatorname{rank}e_B=t,\qquad \operatorname{rank}e_1=1.
\]

At \(\pi\) the first projects onto the character block and the second onto a line of \(\rho_\pi\). These data can be
transported from an integral marked model on the same coefficient
localization. Write \(R_x^{\rm m}\) for the marked model at \(x\).
The shared Hecke character gives
\[
 \det(1-TR_x(\operatorname{Frob}_v))
   =\det(1-TR_x^{\rm m}(\operatorname{Frob}_v))
\]
at the unramified Frobenius classes. Taking traces, then using
continuity and Chebotarev density, gives equality of the trace
functions on the entire Galois group. Brauer--Nesbitt applied
to the two semisimple rational representations gives an
intertwiner at every \(x\). At \(\pi\), cancelling the specified
character summands identifies the rank-\(r\) complements.
Write \(C\) for the resulting tuple of changes of frame and
\(\mathcal A^{\rm m},\mathcal A'\) for the two integral spans.
The generator identity above gives
\begin{equation}\label{eq:selmer-transport-corner}
 C\mathcal A^{\rm m} C^{-1}=\mathcal A',\qquad
 C(e\mathcal A^{\rm m}e)C^{-1}
      =(CeC^{-1})\mathcal A'(CeC^{-1})
      \quad(e=e_B,e_1).
\end{equation}
By \eqref{eq:selmer-transport-corner} the corner matrix is conjugated in every fibre. Its
characteristic roots, the two idempotent ranks and their
orthogonality are preserved.  \label{cor:split-selmer-2}

The cyclotomic test
\(g_0\) lies in the prescribed inertia and has
\[
 \epsilon(g_0)=1+\ell^s u,
\]
with \(u\) a unit. In the new basis
\(X_0(\pi)\) is \(e_B R_\pi(g_0)e_B\). This identification follows
because both models of the direct sum transport the action of
every sum \(\sum_j a_j(\pi)R_\pi(g_j)\) to the same block action.
Thus the perturbation \(X_0-e_B R(g_0)e_B\) lies in the new span,
is supported on the corner, and vanishes at \(\pi\). The ranks,
corner characteristic polynomials and orthogonality relations
are unchanged by conjugation. The matrix-span and ordered-root
construction consequently applies with the same Hecke localization
and the same marked point.

Let
\[
 B_\pi\subseteq M_r(O)
\]
be the image of the normalized algebra
on the specified rank-\(r\) lattice. Absolute irreducibility of its
reduction gives the full matrix algebra over the residue field by
the density theorem. Writing \(\varpi\) for a uniformizer, this says
\[
 B_\pi+\varpi M_r(O)=M_r(O),\qquad
 M_r(O)/B_\pi=\varpi\bigl(M_r(O)/B_\pi\bigr).
\]
Nakayama's lemma applies to this finite \(O\)-module. Thus \(B_\pi=M_r(O)\): the localized order contains every integral endomorphism of the complementary lattice used in the row calculation.
Section~\ref{subsec:retained-galois} constructs the marked model.
Proposition~\ref{prop:retained-rank-complement} proves residual
absolute irreducibility at this localization.

\end{proof}
\subsubsection*{The lattice used in the rank-\(r\) row quotient}

The new coordinates identify the residual complementary block. Computing its row quotient requires the particular lattice on which its normalized Galois action operates, together with the dual action at the conjugate place.

Write \(U_g=\chi_1(g)^{-1}R(g)\), \(Y=\sum_g\mathcal A(U_g-1)E\),
and \(y(g)=(U_g-1)E\) in the normalized algebra. Let \(T\) be
the specified lattice of the \(U\)-action on the rank-\(r\) block,
so \(W=T[1/\ell]\). Write \(T^*=\operatorname{Hom}_O(T,O)(1)\)
and \(W^*=T^*[1/\ell]\). At a row place the dual finite-class map uses the source lattice and right action
\[
 P=e_1\mathcal A,\qquad g\cdot b=\epsilon(g)bU_g^{-1}.
\]
To see its target, write
\(e_1(\pi)=v\lambda\) with \(\lambda(v)=1\) over \(O\) in the
rank-\(r\) block. An element \(b\in P\) has evaluation \(v\lambda_b\)
on this block. The map \(b\mapsto\lambda_b\) onto \(T^*\) is
surjective because the normalized algebra maps onto \(M_r(O)\).
\begin{lemma}[The local row quotient]
Evaluation \(p:P\to T^*\), \(b\mapsto\lambda_b\), is an integral
surjection for the twisted dual action. With
\(A=W/T,\ L_w=H^1_f(E_w,T^*)\) and
\(J_w=p_*H^1_f(E_w,P)\), put
\(F_w=\operatorname{im}(H^1_f(E_w,W)\to H^1(E_w,A))\).
Local Tate duality induces
\[
 J_w^\perp/F_w\simeq\operatorname{Hom}_O(L_w/J_w,F/O).
\]
\end{lemma}
\begin{proof}
Since the block maps onto \(M_r(O)\), every row \(\lambda_b\) occurs.
If \(\lambda_b\) is a row functional, transposition converts its right action into the column-coordinate action
\[
 \lambda_b^{\mathsf T}\longmapsto\epsilon(g)W(g)^{-\mathsf T}\lambda_b^{\mathsf T}.
\]
This is precisely the action on \(W^*\). For the standard basis \(e_i\), integrality is the coordinate test
\[
 u(O^r)\subseteq O\quad\Longleftrightarrow\quad
 (u(e_1),\ldots,u(e_r))\in O^r
 \qquad(u\in\operatorname{Hom}_F(F^r,F)).
\] Restriction through the
chosen completion identifies the decomposition group with the
local absolute Galois group. The coefficient isomorphism there
carries the fraction field and DVR simultaneously, so it transports
continuous cocycles, coboundaries and the finite conditions with
the lattice. Recording matrix rows at the ordered points embeds the rational source in a finite sum of the evaluated twisted duals. This is the row-colength map.  \label{cor:split-selmer-3}

Use local Tate duality in its \(F/O\)-valued form, obtained from
the usual \(\mathbb Q_\ell/\mathbb Z_\ell\)-valued pairing by the
trace and a generator of the inverse different. An integral class
annihilates \(F_w\) precisely when its rational image pairs to zero
with \(H^1_f(E_w,W)\): any nonzero rational value can be scaled to
a nonintegral one. By Bloch--Kato orthogonality this condition
cuts out \(L_w\). Double annihilation gives
\[
 F_w=L_w^\perp.
\]

The restriction
\[
 J_w^\perp\to\operatorname{Hom}_O(L_w/J_w,F/O)
\]
has this kernel and is onto by injectivity of \(F/O\) over the DVR.
The extension is continuous on the finitely generated integral
cohomology group. Thus the normalized length equals that of
\(L_w/J_w\), bounded by the row index at this same place.
\cite{bloch-kato}
\end{proof}

\begin{lemma}[Annihilation of the integral row classes]
For an \(\mathcal A\)-linear \(\varphi:Y\to A\), put \(c(g)=\varphi(y(g))\). Local Tate duality gives
\[
 \operatorname{res}_w[c]\in J_w^\perp.
\]
\end{lemma}
\begin{proof}
To prove annihilation of \(J_w\), let \(a\) be a continuous
row cocycle with values in \(P\), and let \(c(g)=\varphi(y(g))\)
with values in \(A\). The multiplication cochain
\[ z(g,h)=a(g)U_g y(h) \]
takes values in \(e_1Y\), with the action on this lattice
given by the scalar \(\epsilon\).

The contraction \(z\) is a continuous two-cocycle. This follows by the cocycle identity for \(a\)
with action \(b\mapsto\epsilon(g)bU_g^{-1}\) and the
identity \(y(gh)=y(g)+U_g y(h)\). Its image in
\(e_1\mathcal A E(1)\) is a coboundary. Indeed the
one-cochain \(d(g)=a(g)U_g E\) satisfies
\[
 \epsilon(g)d(h)-d(gh)+d(g)=-a(g)U_g(U_h-1)E=-z(g,h).
\]
The inclusion \(e_1Y\subseteq e_1\mathcal A E\) induces an
injection on continuous \(H^2\) after the Tate twist. Both
lattices are finite free over \(O\), and the local invariant
map
\[
 H^2(E_w,\mathbb Z_\ell(1))\simeq\mathbb Z_\ell,
\]
extended to finite free coefficient modules, identifies the
induced map with the original lattice inclusion. This works
also for a nonsaturated inclusion. Therefore the cochain
\(z\) is already a coboundary in \(e_1Y(1)\).

Apply to that smaller lattice the map \(m\mapsto\lambda(\varphi(m))\).
The pairing is well defined modulo \(O\), and is equivariant
for the scalar Tate action. Since \(\varphi\) is
\(\mathcal A\)-linear, the image of \(z(g,h)\) is
\(\langle p(a(g)),U_g c(h)\rangle\), with
\[
 p:P\to T^*
\]
the row quotient.

This is the cup pairing in the reversed
order, which differs from the desired order by the sign
for two degree-one classes. It follows that \([c]|_{E_w}\)
annihilates the image of every such \([a]\), in particular
\(J_w\). The local class of \(c\) is continuous, as the
universal cocycle is continuous in its finite integral
lattice and a linear map from that lattice into \(A\) has
finite image and is continuous. Consequently the integral index \(\operatorname{len}(L_w/J_w)\) bounds the excess over the finite condition on the complementary row. The character rows require their own local conditions. \label{cor:split-selmer-4}
\end{proof}
\subsubsection*{Character row and column conditions}

The complementary row has now been accounted for. On the remaining rank-one blocks the cyclotomic twists distinguish the finite conditions, and the row and column impose different restrictions even at the same coefficient prime.

\begin{proposition}[Local conditions on the character blocks]
At a column the specialization is finite on block \(0\),
unramified on block \(1\), and zero on blocks \(k\ge2\).
At a row above \(\ell\), for \(1\le k\le t\) let \(L_k\) be the normalized length of the excess over the finite subgroup on block \(k\). Then
\[
 L_k\le([E_w:\mathbb Q_\ell]+2)c_B,\qquad c_B=40(t+1)(n+1)\lambda_N(h).
\]
\end{proposition}
\begin{proof}
At a column, the last defining relation of \(M\) reads in local
cohomology as a combination of the classes \(K_{u,a}\) and
\(K_{u,0}\) after evaluation. In block \(0\) the first classes
are finite and the last term vanishes because \(K_{u,0}\) has a
right factor \(E\). In block \(1\) the first terms vanish because
\(E_1\) acts by zero, and the last class is unramified. In
blocks \(k\ge2\) both sets of terms vanish. The generator
\(n_u\) contributes a coboundary in each case. These conclusions
hold at the normalized point and use the specific column lifts
of the determinant construction.

For the character rows, form the stable lattice \(Y\oplus O\)
with action
\[
  (z,a)\longmapsto (U_g z+a\,y(g),a).
\]
With \(e=(0,1)\), the quotient is the trivial line, and
\(ge-e=(y(g),0)\). Put \(i(z,a)=(z+aE,a)\).
The cocycle is \(y(g)=(U_g-1)E\), so
\[
    i(U_g z+a\,y(g),a)=(U_g(z+aE),a)=g\,i(z,a).
\]
The map \(i\) is injective, with the left \(U_g\)-action
on matrices and the trivial action on the second factor.

Evaluation
on the matrix columns embeds the rational space of this lattice
into a finite sum of the twisted family \(R_x^\circ\) and one
trivial line. The trivial character occurs in \(R_\pi^\circ\). The target of
the embedding is therefore crystalline, and so is its rational
subrepresentation \(V=(Y\oplus O)\otimes_O F\). For each
embedding \(\tau\), strictness of the crystalline filtration gives
\[
 \operatorname{Jumps}_\tau(V)
     \subseteq \bigcup_x \operatorname{Jumps}_\tau(R_x^\circ)
     \subseteq [-h,h].
\]
The union on the right over all embeddings has cardinality at
most \(n\) by the common row bound, with \(t\le h\).

The character lifting bound is used with that lattice, the
surjection onto the trivial line and the lift \(e\) just given.
It applies to an equivariant map of the kernel \(Y\) into the
divisible quotient of a rank-one lattice acted on by
\(\epsilon^{1-k}\), for \(1\le k\le t\). Extend this integral
map \(O\)-linearly to the full rational source using injectivity
of \(F/O\) over \(O\). The comparison only requires equivariance
on the kernel lattice.

Indeed, if \(\phi\) is the equivariant map on \(Y\) and
\(\widetilde\phi\) any such extension, then
\begin{equation}\label{eq:selmer-row-kernel-cocycle}
 c(g)=\widetilde\phi(ge-e)=\phi(y(g)),\qquad
 c(gh)=c(g)+\epsilon^{1-k}(g)c(h).
\end{equation}
Both identities in \eqref{eq:selmer-row-kernel-cocycle} use only the restriction to \(Y\). The completion lies in
the local level-\(\operatorname{lcm}(\ell,N)\) cyclotomic field,
the coefficients contain those roots of unity, and \(4\mid N\)
at \(\ell=2\). With
\[
 c_B=40(t+1)(n+1)\lambda_N(h),
\]
the
comparison places \(\ell^{c_B}[c]\) in the finite subgroup of
the target. For \(k=1\) this subgroup is unramified. For \(k>1\)
it is zero, by the rank-one jump \(k-1>0\) and the absence of
rational invariants in the Bloch--Kato dimension formula.
\cite{bloch-kato}\label{cor:split-selmer-5}

For the length allowance, let \(F_{k,w}\) be this finite
subgroup. It is divisible, as the image of a vector space.
Choose \(b\in F_{k,w}\) with \(\ell^{c_B}b=\ell^{c_B}[c]\).
Then \([c]-b\) represents the same class modulo \(F_{k,w}\) and
belongs to
\[
 H^1(E_w,A_k)[\ell^{c_B}]
   =\operatorname{im}\bigl(H^1(E_w,A_k[\ell^{c_B}])
                          \longrightarrow H^1(E_w,A_k)\bigr),
\]
by the cohomology sequence for multiplication by \(\ell^{c_B}\). For \(M_k=A_k[\ell^{c_B}]\), local Tate duality and the bound on
the coefficient module give \(\operatorname{len}H^i(E_w,M_k)\le c_B\)
for \(i=0,2\). The local Euler characteristic gives
\[
 \operatorname{len}H^1(E_w,M_k)
 = [E_w:\mathbb Q_\ell]c_B+
   \operatorname{len}H^0(E_w,M_k)+\operatorname{len}H^2(E_w,M_k).
\]
This gives an allowance for each character row even when a sharper row estimate is available. It leaves the inertia conditions away from the coefficient prime to be imposed.
\end{proof}
\subsubsection*{Tame tests on the character blocks}

It remains to control ramification away from the split row and column places. The tame inertia relations on each character block bound the quotient by unramified classes at the auxiliary level primes.

For a character block let \(A_k=(F/O)(1-k)\) and let
\(c(g)=\varphi(y(g))\) for an \(\mathcal A\)-linear map
from the column presentation into \(A_k\).

\begin{lemma}[Tame inertia and the remaining global condition]
Away from the split row and column places the character class is
unramified outside \(Q\); at \(v\in Q\) the ramified quotient \(D_v\) satisfies
\[
 \operatorname{len}(D_v)\le v_\ell(\zeta_v-1).
\] After imposing
the row, column and unramified conditions the residual length is at
most \(v_\ell(h_E)\) for \(k=1\); for \(k\ge2\) it is bounded by
the strict length with exponent \(k-1\) over the same field.
\end{lemma}
\begin{proof}
At a place
outside the split row and column places and the auxiliary
sets \(S,Q\), the family inertia acts trivially on every
fibre and \(\chi_1\) is trivial there. The last assertion
can be checked on its summand at \(\pi\). Thus \(y(g)=0\)
on this inertia.

At \(v\in Q\), the chosen inertial element
\(\gamma_v\) maps every fibre into the two eigenspaces for
\(1\) and \(\zeta_v=\chi_1(\gamma_v)\). The defining conditions on the tame generator give \(\zeta_v\ne1\), equality of all
\(\chi_k\) on inertia, and a common power \(n(g)\) of
\(R(\gamma_v)\) and of \(\chi_1(\gamma_v)\) for the actions
of each \(g\) on this inertia. The common power comes from
the finite image of the chosen tame character. In the family
algebra we have
\[
 (R(\gamma_v)-1)y(\gamma_v)
 =\zeta_v^{-1}(R(\gamma_v)-1)(R(\gamma_v)-\zeta_v)E=0.
\]

Evaluating under \(\varphi\) gives
\((\zeta_v-1)c(\gamma_v)=0\). On the target \(U_{\gamma_v}\) acts trivially. On ordinary powers the
cocycle therefore satisfies
\[
 c(\gamma_v^a)=\sum_{i=0}^{a-1}U_{\gamma_v}^ic(\gamma_v)=a\,c(\gamma_v)
 \qquad(a\ge1).
\]
Continuity on the procyclic quotient gives \(c(g)=n(g)c(\gamma_v)\).
Coboundaries vanish on inertia here. Therefore the quotient
by the unramified condition injects, by evaluation at
\(\gamma_v\), into \(A_k[\zeta_v-1]\), of normalized length
\(v_\ell(\zeta_v-1)\). This bound uses only \(\zeta_v-1\ne0\).

At \(v\in S\), there is an integral family operator \(a\)
such that
\[
 a(R(g)-\chi_1(g))E=0
\]
on this inertia and
\(a\) acts invertibly on \(A_k\). At the two type primes,
\(a=P_v(R(\gamma_v))\), where the coefficients of \(P_v\)
are in the localized Hecke algebra. The tame polynomial
condition asserts that on each fibre this polynomial kills
\(R(g)-\chi_1(g)\) and that \(P_v(\chi_1(\gamma_v))\) is
a unit upon evaluation at \(\pi\). Since all the scalar
characters agree on the chosen inertia, the second property
is exactly invertibility on \(A_k\).

At the remaining
\(S\)-places the family is simultaneously diagonal on inertia
with values \(1\) or \(\chi_1\) on the corresponding
summands. The factors kill the trivial and the \(\chi_1\) summands respectively. Hence on this inertia
\[
 (R(\gamma_v)-1)(R(g)-\chi_1(g))=0.
\]
At these places
\(\zeta_v-1\) is a unit, so the former operator is invertible
on the character target. Consequently \(c\) vanishes on
inertia at every \(S\)-place. These restriction bounds apply under their respective tame hypotheses.
\label{cor:split-selmer-6}

Subtract the row and \(Q\)-local allowance by taking the
kernel of restriction to the corresponding quotients. The
strict conditions for \(k\ge2\) at the coefficient prime
are zero on row and column decomposition groups. For \(k=1\)
the local finite condition for the trivial character at those
places is unramified. The declared rows and columns partition
the places over \(\ell\), and the remaining inertia groups
are among \(S,Q\) or the places with trivial family inertia.
In block \(k=1\), global class field theory therefore gives
\[
 \ker(\text{restriction to the local quotients})\hookrightarrow\operatorname{Hom}(\operatorname{Cl}(E),F/O).
\]
Its normalized length is at
most \(v_\ell(h_E)\).

For \(k\ge2\), let \(S_\ell\) consist of
the places over \(\ell\) and the infinite places. Put
\(A=(\mathbb Q_\ell/\mathbb Z_\ell)(1-k)\). The strict group is
\[
 S_{k-1}=\ker\!\left(H^1(G_{E,S_\ell},A)
        \longrightarrow\prod_{w\mid\ell}H^1(E_w,A)\right).
\]
The restriction kernel maps into \(S_{k-1}\otimes O\).
Indeed the local class vanishes at \(\ell\), all other places
are unramified, and finite freeness over \(\mathbb Z_\ell\) lets
tensoring by \(O\) commute with cochains and localization kernels.
Each \(\mathbb Z_\ell\)-length-one factor becomes \(O/\ell O\),
so the normalized length is \(\operatorname{length}_{\mathbb Z_\ell} S_{k-1}\). The strict
cohomology bound below applies with this exponent and on
this same auxiliary field. \label{cor:split-selmer-7}

\end{proof}
For block \(0\), the same vanishing at the places with trivial
family inertia follows directly from \(y(g)=0\). The row and
column conditions in the restriction kernel lie in the finite
condition on \(W/T\) at all places over \(\ell\). These are
the unramified and finite conditions of the relaxed Selmer
group at \(\pi\) used for \(s_\pi\). The length at this fixed point can therefore be compared with the signature Selmer length, provided the auxiliary character is the one used in its theta lattices.

\subsubsection*{Choosing one twist for the Selmer group and the theta lattices}
\label{subsec:retained-twist}
\hypertarget{source-retained-twist}{}

The same localization supplies a Selmer length and theta denominators. A single character must therefore control every signature and theta value occurring there.

Fix a twisted coefficient place of residue characteristic \(\ell\).
Write \(v(\ell)=1\) in every coefficient field. If \(O\) is a
coefficient DVR, normalized length means \(\operatorname{length}_O/e_O\),
where \(e_O=\operatorname{ord}_O(\ell)\). There are finitely many
good signatures \(s\) for the lower cusp at this place, finitely
many abelian theta values indexed by \(\lambda\), and two sides
for each of the \(m\) theta rows. The sign of the twist on each
of these objects is fixed by its side. A choice for each object
separately would leave the congruence comparison on different
Hecke modules. We describe the simultaneous specialization and
its passage to the retained module.

\begin{established}[Common character data and coefficient tests]\label{hyp:simultaneous-retained-twist}
At each nonreference coefficient place of the initial signed row, both theta comparisons use the same ray coordinates. For each occurring good signature the integral conjugate-duality isomorphism and the isotypic and valuation tests hold. Its jumps and type primes satisfy
\[
 \operatorname{HT}_\tau(T_s)\cap[-1,t]=\varnothing
 \quad(\tau\mid\ell),\qquad
 r\ge100,\quad\min(g_0,g_1)\ge2r .
\] The abelian values have nonzero untwisted specialization and an allowed critical integer. The two Hecke-stable integral rows admit commuting projected theta bases with uniformly bounded denominators and the required residual eigensystems.
\end{established}
\begin{proof}
Fix the row and both signs before selecting a ray character. Lemma~\ref{lem:initial-row-signature-tests} applies to the signatures in the images of the two initial type projectors. It proves the assertions on their lattices, jumps and residual valuations. The simultaneous numerical choices are those of Lemma~\ref{lem:window-parameter-choice} and Proposition~\ref{prop:two-type-primes}. The nonzero abelian specialization is the value in \ref{cor:toric-twist-5}, with the critical range and common denominator of Lemma~\ref{lem:toric-common-denominator}. On either signed row Proposition~\ref{prop:signed-theta-data} gives the commuting action on the integral class sums, their universal denominator and the residual character. Proposition~\ref{prop:retwisting-row-basis} extracts a basis from their projected span before specialization. Both signs are functions of the same class coordinates: they apply their respective coefficient maps to the same finite translations. These are precisely the spaces over which the subsequent character is chosen.
\end{proof}

\begin{auxlemma}[A bound for spherical parameters]
\label{lem:initial-spherical-bound}
Let \(q\) be the residue cardinality of a nonarchimedean local field
and \(\beta_1,\ldots,\beta_r\) the Satake parameters of an unramified
unitarizable representation of \(\mathrm{GL}_r\), in the normalized
Satake convention. For \(r\ge1\),
\[
             \max_i|\beta_i|\le4q^{(r-1)/2}\le q^{(r+3)/2}.
\]
\end{auxlemma}
\begin{proof}
Normalize Haar measure so that the maximal compact subgroup has volume one.
The Hecke operator of \(\operatorname{diag}(\varpi I_j,I_{r-j})\)
is a sum over \(\genfrac{[}{]}{0pt}{}{r}{j}_{\!q}\) residue subspaces of dimension \(j\), each giving
one unitary translate of a spherical vector of norm one. The Gaussian
binomial coefficient counts these subspaces by the number of ordered
independent \(j\)-tuples divided by \(|\mathrm{GL}_j(\mathbb F_q)|\). The normalized Satake formula gives
\[
 |e_j(\beta)|\le q^{-j(r-j)/2}\genfrac{[}{]}{0pt}{}{r}{j}_{\!q}
 \le q^{j(r-j)/2}\prod_{h=1}^j(1-q^{-h})^{-1}.
\]
The last product is at most \(2\) for \(j=1\) and less than \(4\)
for \(j\ge2\), since \(\prod_{h\ge3}(1-2^{-h})\ge3/4\).
Thus \(|e_j(\beta)|^{1/j}\le2q^{(r-1)/2}\).

If a root \(z\) of
the Satake polynomial satisfied \(|z|>4q^{(r-1)/2}\), division by
\(z^r\) would give
\[
      1\le\sum_{j=1}^r|e_j(\beta)/z^j|
                   < \sum_{j=1}^r2^{-j}<1 .
\]
Consequently \(|\beta|\le4q^{(r-1)/2}\le q^{(r+3)/2}\).
\end{proof}

\begin{auxlemma}[The signatures on the initial signed rows]
\label{lem:initial-row-signature-tests}
Fix a coefficient index \(i\) at which the marked character differs
from the reference character. Write \(O_i\) for its coefficient DVR and \(\ell\) for its residue characteristic. Use the initial rows and the two type markings of Proposition~\ref{prop:signed-theta-data}. A candidate signature \(s\) records a lift \(j\), a side \(b\) and
split Hecke eigenvalues \(a_{s,k,P}\) of the lower cuspidal row. Write
\(e_b\) for its assigned finite type projector and \(V_{i,j,b}\) for
the rational space of that lower row. Use precisely the signatures with
\[
 \begin{gathered}
 \exists v_s\in V_{i,j,b}\setminus\{0\},\qquad e_bv_s=v_s,\\
 T_{k,P}v_s=a_{s,k,P}v_s\quad(1\le k\le r),\qquad
 P\text{ split outside the excluded set}.
 \end{gathered}
\]

Let \(\rho_s\) denote the integral Galois representation of this
lower cusp, \(\xi_s\) the first character of its prescribed string,
and \(W=\rho_s\otimes\xi_s^{-1}\) on its chosen coefficient lattice.
Write \(c_E\) for complex conjugation. Then
\[
 (W^\vee(1))^{c_E}\simeq W(t),\qquad
 \operatorname{gr}^a D_{\rm dR}(W)_\tau=0
       \quad(-1\le a\le t,\ \tau\mid\ell).
\]
The isomorphism is integral. Both \(W\) and its Tate dual have
zero decomposition invariants over the coefficient field at every
finite place of the host.
\end{auxlemma}
\begin{proof}
The finite-dimensional projected Hecke spaces are sums of their
simultaneous eigenspaces, by Proposition~\ref{prop:rational-row-isotypic-span}.
Taking the nonzero summands thus gives a finite list on both sides and
at all lifts. Each summand has the cohomological cuspidal attachment
of Proposition~\ref{prop:signed-theta-data}. The two selected types
satisfy the divisor-degree and separation tests of
Definition~\ref{hyp:retained-type-tests}. The inertia argument of
Proposition~\ref{prop:retained-rank-complement} gives absolute
residual irreducibility for \(\bar\rho_s\). These assertions concern
the marked type polynomials on the lower row. On the upper theta
locus Proposition~\ref{prop:cut-selector} identifies the matching
points by
\[
                    Z_i\cap S_i=Z_i^+.
\]

Let \(k_i\) be the residue field of the signature. Absolute
irreducibility and the matrix algebra theorem give the first equality
below after extension to \(\overline{k_i}\), hence also over \(k_i\).
The span is already an algebra because \(\bar\rho_s\) is a
group representation. For \(v\mid g_d\) with \(g_d\ne\ell\), the
depth-zero action is trivial on wild inertia \(P_v\). The tame labels
are separated from the first string character by the tests in
Definition~\ref{hyp:retained-type-tests}. Reading them on a marked
tame inertia generator \(u_v\) gives
\[
 \begin{gathered}
 \operatorname{span}_{k_i}\bar\rho_s(G_E)=M_r(k_i),\qquad
 \bar\rho_s(P_v)=1,\\
 \det\bigl(\bar\rho_s(u_v)-\bar\xi_s(u_v)I_r\bigr)\ne0 .
 \end{gathered}
\]
These are the residual tests on the coefficient ring. Multiplication
by the scalar character \(\bar\xi_s^{-1}\) preserves absolute
irreducibility.

Put \(h=(t-1)/2\). By the weight calculation of
Lemma~\ref{lem:marked-signed-jumps}, the labelled jumps in the
complement have the form \(a=h\pm z\), where the interlacing
inequalities give \(|z|\ge h+2\). Consequently
\[
 a\le h-(h+2)=-2\quad\text{or}\quad
 a\ge h+(h+2)=t+1.
\]
At the conjugate embedding the corresponding jump is \(t-1-a\).
The displayed ranges therefore give exactly one negative jump for
each pair of corresponding entries. With multiplicities this yields
\[
 \dim(D_{\rm dR}(W)_\tau/\operatorname{Fil}^{0})+
 \dim(D_{\rm dR}(W)_{\bar\tau}/\operatorname{Fil}^{0})=r.
\]
The untwisted lower cusp is unramified outside \(\ell,g_0,g_1\).
At \(v\mid\ell\) the two attachments \(\rho_s,W\) are de Rham.
They are crystalline if \(\ell\notin\{g_0,g_1\}\), and potentially
crystalline in the marked type case because their local monodromy
vanishes, as in Proposition~\ref{prop:signed-theta-data}.

At a split unramified Frobenius the polarized attachment, with the
same character normalization on both sides, identifies the
characteristic polynomials of \((W^\vee(1))^{c_E}\) and \(W(t)\).
Chebotarev and Brauer--Nesbitt give an intertwiner over the fraction
field. Multiply it by a scalar to make its entries integral with at
least one unit. For its reduction \(\bar U\), residual
irreducibility gives
\[
 \bar U\ne0\quad\Longrightarrow\quad\ker\bar U=0
 \quad\Longrightarrow\quad \det U\in O_i^\times .
\]
This proves
the required isomorphism on the given lattices.

At \(v\mid\ell\), a nonzero fixed vector would give a jump zero.
For the Tate dual the jumps are \(-1-a\). Thus
\[
        H^0(E_v,W\otimes_{O_i}\operatorname{Frac}(O_i))=
        H^0(E_v,W^\vee(1)\otimes_{O_i}\operatorname{Frac}(O_i))=0 .
\] 
At a type prime of residue characteristic different from
\(\ell\), the separated tame labels give a generator with no
eigenvalue one on the normalized complement. Its inverse transpose
has the same property. At the distinguished auxiliary places the
normalized inertia action includes a nonidentity scalar, with its
inverse on the Tate dual. 
At any remaining place \(v\), put \(q_v=\mathrm Nv\). An arithmetic Frobenius
root is \(\alpha=\beta q_v^{(1-t)/2}\), where \(\beta\) is an
unramified Satake parameter of the unitary lower cusp. Lemma~\ref{lem:initial-spherical-bound} gives
\(|\beta|\le q_v^{(r+3)/2}\). Since \(t\ge r+5\), this gives
\[
 |\alpha|\le q_v^{(r+4-t)/2}\le q_v^{-1/2}<1<q_v .
\]

A fixed vector in the representation or its Tate dual would force
\(\alpha=1\) or \(\alpha=\mathrm Nv\), respectively. This covers
type primes above \(\ell\) by the first case.
\cite{blggt-attachment,caraiani-monodromy}
\end{proof}

\begin{proposition}[One character for the Selmer and theta comparisons]\label{prop:one-retained-character}
For the signatures, abelian values and rows of
Lemma~\ref{hyp:simultaneous-retained-twist}, one finite-order anticyclotomic character
with the prescribed coordinate readings gives all the
following bounds on the original local models, with \(\nu_\ell\) as in
\eqref{eq:signature-cyclotomic-minimum}:
\begin{itemize}
\item for the normalized length \(L_s(\chi)\) at signature \(s\) after applying this character \(\chi\), there are errors \(e_s\) with
\[
 L_s(\chi)\le\nu_\ell+e_s<\infty,\qquad e_s\ge0,\qquad\sum_s e_s\le\ell^{-2};
\]
\item the abelian theta values are nonzero and have valuation
  between \(-N_0\) and \(v(\theta_\lambda(1))+\ell^{-2}\),
  with \(N_0\) the common integral denominator bound;
\item the nonnegative row-index excesses over the untwisted
  indices have sum at most \(\ell^{-2}\);
\item the sum of the valuations of the nontrivial marked
  inertia readings, multiplied by \(t\), is at most \(\ell^{-2}\).
\end{itemize}
The twist may avoid a prescribed finite set and can be
transported at conjugate coefficient places with the
prescribed inversion.
\end{proposition}

\begin{proof}
\proofstep{A common parameter space for the three tests.}
Proposition~\ref{prop:ray-coordinate-pairs} supplies the two
groups of pairs. For the fixed finite family there exist integers \(R_*,D_*,M_0\).
Use \(2R_*\) variables, grouped as \(X_i,Y_i\). At every
\(M\ge M_0\), enlarged to be at least one, there are \(2R_*\)
distinct conjugate pairs of split degree-one places, unramified
over \(\mathbb Q\), of norms \(1\bmod 2\ell^M\). Their union
\(Q\) avoids the level and any prescribed finite set.

There are
continuous anticyclotomic coordinate characters \(d_i\) to
\(\mathbb Z/\ell^M\), each surjective on inertia at both places
of its pair, unramified away from the pair and trivial on all
decomposition groups above \(\ell\). The first half of the
selected Frobenius elements are conjugate to a prescribed tame element
modulo an open normal subgroup. If \(v_j\) denotes the selected representative of the
\(j\)-th second-half pair, the first-half characters satisfy
\[
 \bigl(d_i(\operatorname{Frob}_{v_j})\bigr)_{1\le i,j\le R_*}
       =I_{R_*}\quad\text{over }\mathbb Z/\ell^M\mathbb Z .
\] The
open subgroup may be shrunk and the avoidance set enlarged to
satisfy the finitely many abelian and row conditions.

\begin{enumerate}
\item For each good \(s\) there is an integral polynomial
\(F_{s,M}\) in these variables over a fixed coefficient DVR.
One coefficient of total degree at most \(D_*\) has valuation
at most \(\nu_\ell\). For coordinates \(u_i^{\ell^M}=1\), put
\[
 \psi(g)=\prod_i u_i^{d_i(g)},
\]
over both halves of the
coordinates. The signature Selmer length for \(\psi\) equals
\(v(F_{s,M}(u-1))\), with value infinity at zero. To see
the determinant step in this description, take the square
matrix over \(O[(\mathbb Z/\ell^M)^{2R_*}]\) whose kernel on
discrete \(\ell\)-primary modules computes the Selmer functor
with inverse universal twist. Lift its determinant by sending
each variable to the corresponding group element minus one.

Let \(A(u)\) be the specialized matrix, \(O_u\) the
valuation ring at this specialization and \(K_u\) its fraction
field. Write \(a\) for the size of this square presentation. Its Smith entries
\(c_j\) give the kernel on \((K_u/O_u)^a\), one summand at a time:
\[
 \ker(c_j:K_u/O_u\to K_u/O_u)=
 \begin{cases}
    c_j^{-1}O_u/O_u,&c_j\ne0,\\
    K_u/O_u,&c_j=0 .
 \end{cases}
\]
The normalized length of the first summand is \(v(c_j)\).
Hence, with \(v(0)=+\infty\), specialization yields
\begin{equation}\label{eq:selmer-twist-smith-specialization}
 \operatorname{len}_{\rm norm}\ker A(u)
       =\sum_j v(c_j)=v(\det A(u))=v(F_{s,M}(u-1)).
\end{equation}
For the opposite side, evaluate \eqref{eq:selmer-twist-smith-specialization} at \(u^{-1}\).
\item For each abelian value and each row there is an integral
polynomial \(f_M\) satisfying the following condition, over a
fixed DVR for that object. There are integral polynomials
\(B_M,g_M\), an integer \(b\ge0\) independent of \(M\), and a
bound \(n_*\) independent of \(M\) such that
\[
 f_M\equiv B_M g_M\pmod{(Y_i,(1+X_i)^{\ell^M}-1)_i},
 \qquad \operatorname{ord}_O(g_M(0))\le b.
\]
Modulo the uniformizer and the powers \(X_i^{\ell^M},Y_i^{\ell^M}\),
\(B_M\) is a monomial \(\prod_i X_i^{n_i}\), \(n_i\le n_*\),
times a polynomial with nonzero constant term.

For an abelian
label the value of \(f_M(u-1)\) is \(\ell^{N_0}\theta_\lambda(\psi)\),
where
\[
 b/e_O=N_0+v(\theta_\lambda(1))
\]
and the last value
is taken at empty auxiliary level. For a row, at any specialization
with nonzero value, the relative index of any lattice containing
the specified projected twisted generators, in their prescribed
rational span, is at most the untwisted index plus
\(v(f_M(u-1))-b/e_O\). These two formulas are used
with the sign appropriate to the label or row.
\end{enumerate}

We prove that for any \(\delta>0\) the valuations of all these
polynomials at one signed tuple are bounded by their respective
baselines \(\nu_\ell,b/e_O\) plus \(\delta\). The tuple can
also be required to have
\[
 2t\sum_i v(u_i-1)\le\delta,
\]
where \(i\) runs over both
halves and all orders are nontrivial. We give a limiting argument using discrete valuations.
The theta coefficient DVR has algebraically closed residue field,
so compactness of its coefficients is unavailable.

Choose data at a sequence of levels tending to infinity. There
are finitely many choices of the distinguished monomials of the bounded-degree signature and theta coefficient tests in the three families of polynomials. Fix them on a subsequence. In each fixed
coefficient DVR take the sequences of all coefficients of the
polynomials, including \(B_M,g_M\), and the constant sequence
of a uniformizer. The subring of the ring of sequences that
they generate over \(\mathbb Z\) is countable. Enumerate the countably many coefficient sequences as \(a_1,a_2,\ldots\). A diagonal subsequence \(M_j\) then satisfies
\[
 \lim_{j\to\infty}\operatorname{ord}(a_i(M_j))\in\mathbb N\cup\{\infty\}\qquad(\text{each }i).
\]

This is possible because any sequence of nonnegative orders
has a constant subsequence or one tending to infinity.
For such a sequence ring \(R\), put
\[
       I_\infty=\{a\in R:\operatorname{ord}(a_M)\to+\infty\},
       \qquad v_\infty(\bar a)=\lim_M\operatorname{ord}(a_M).
\]
The valuation inequality and nonnegative orders show that
\(I_\infty\) is an ideal. If \(\bar a\ne0\), its order is
finite (hence eventually constant), and adding an element
of \(I_\infty\) does not change it. For nonzero classes
\[
 v_\infty(\bar a\bar b)=v_\infty(\bar a)+v_\infty(\bar b),
 \qquad v_\infty(\bar a+\bar b)\ge
                   \min(v_\infty(\bar a),v_\infty(\bar b)).
\]
In particular \(R/I_\infty\) is a domain. Extend the order
to its fraction field by subtracting orders. The constant
uniformizer sequence has order one, and the order
of \(\ell\) is unchanged. Use the valuation ring of this field
for the limit series. Every fixed expression in finitely many
coefficients has its limiting order by the definition of the
quotient.

In particular each limit series \(F_s\) has a coefficient
of valuation at most \(\nu_\ell\).

For the second condition, set all \(Y_i=0\). In each fixed
total degree the coefficients of \((1+X_i)^{\ell^M}-1\) tend
to zero, including its zero constant term, by the binomial
formula. Multiplication by arbitrary integral ideal multipliers
preserves this convergence in each bounded degree. The limits
thus satisfy
\[
 f(X,0)=B(X,0)g(X,0).
\]
The selected monomial
coefficient of \(B(X,0)\) is a unit, and
\(\operatorname{ord}_O(g(0))\le b\). For an integral nonzero
formal series let its content valuation be the minimum of the
valuations of its coefficients. This minimum is attained by
discreteness. Write \(c(A)\) for this content valuation. After scaling each nonzero series to have content zero, the residue series are nonzero in a domain. Consequently
\[
 c(AB)=c(A)+c(B).
\]

Consequently the content valuation of \(f\) is at most
\(b/e_O\). This proof permits every \(B_M(0)\) to vanish.
For a negative side first substitute \((1+Z_i)^{-1}-1\) for
each variable \(Z_i\) of the limit and its approximants.
This integral formal automorphism preserves content and is
continuous coefficientwise.

Normalize each resulting limit series by its content so that
its reduction is nonzero. For all of these reductions at once,
choose nonnegative integers \(h_i\) so that substitution
\[
 Z_i=T^{\ell^{h_i}}
\]
leaves a nonzero series. Indeed, for
each series take the lexicographically first monomial of its
support. Starting with the last coordinate, choose each earlier
weight \(\ell^{h_i}\) greater than the total contribution of
the later coordinates in every one of the finitely many chosen
monomials. At the first coordinate where a competing monomial
differs, its increase costs more than can be saved at later
coordinates. Thus the chosen term remains uniquely lowest.
In characteristic zero now substitute
\[
          Z_i=(1+T)^{\ell^{h_i}}-1 .
\]
This reduces to the substitution just considered. Let \(d\)
be the degree of the first unit coefficient in a normalized
one-variable series.

For a primitive \(\ell^a\)-th root
\(\xi\), \(w=v(\xi-1)=1/(\ell^{a-1}(\ell-1))\).
Write \(\sum_k a_k T^k\) for the normalized series.
At \(T=\xi-1\) the difference from the valuation \(dw\)
of its degree-\(d\) term satisfies
\[
 v(a_kT^k)-dw\ge
 \begin{cases}
       1/e_O-(d-k)w,&k<d,\\
       (k-d)w,&k>d.
 \end{cases}
\]
Both bounds are positive when \(dw<1/e_O\). The unique
least term gives valuation \(dw\) for the normalized
series. Take \(a>\max h_i\) so large that
each \(dw<\delta\), all preceding strict inequalities hold,
and \(2t\sum_i\ell^{h_i}w\le\delta\).
Set \(u_i=\xi^{\ell^{h_i}}\), for which
\[
 v(u_i-1)=\ell^{h_i}w.
\]
For zero variables there is simply a
finite list of constants and the empty tuple has zero tame cost.

Fix this \(a\). The preceding single-variable computation uses
only orders of coefficients and the value \(w\). It transfers to
the original polynomials as follows. Write \(m\) for the
normalized content valuation of a limit before dividing by its
uniformizer power. Take \(N\) so large that
\[
 (N+1)w>m+dw
\]
for all series, and \(N\ge d\) for all of
them.

The coefficients of the original one-variable substitutions
through degree \(N\) are fixed integral polynomial expressions
in finitely many input coefficients, including for inverse
coordinates. Their orders therefore converge to the orders in
the limit. Write \(b_j\) for the coefficients at one such level and take
\(T=\xi-1\), so \(v(T)=w\).

For all sufficiently large levels on
the subsequence, finite coefficient convergence and integrality
of the tail give
\[
 \begin{aligned}
 v(b_dT^d)&=m+dw,\\
 \min_{\substack{0\le j\le N\\j\ne d}}v(b_jT^j)&>m+dw,\\
 \inf_{j>N}v(b_jT^j)&\ge (N+1)w>m+dw .
 \end{aligned}
\]
An empty minimum has value \(+\infty\). The ultrametric inequality is
therefore an equality on the full series:
\[
 v(\sum_{j\ge0}b_jT^j)=m+dw.
\]
Inverse-coordinate series may be evaluated in the
completion of the coefficient field, where the substitution
agrees with inversion at the root of unity. The test works
for every primitive root of that order because it uses only
the valuation \(w\) and a unique minimal term. Thus the same
level gives the bounds for every power of the character prime
to \(\ell\).

At such a level \(M\ge a\) use its selected primes and its
characters \(d_i\). The ray class character at the chosen level is read on the same coordinates by
\[
 \psi(g)=\prod_i u_i^{d_i(g)}.
\]
Each coordinate is abelian, unramified off its selected pair and tamely ramified there, so the Artin map gives the desired ray character. The universal theta coordinates at the chosen primes agree with the characters \(d_i\) by the restriction property for the chosen coordinate characters. Ramification is exact because each \(u_i\ne1\)
and no other coordinate ramifies at that pair. The image on either inertia group at pair \(i\) is
\(\langle u_i\rangle\), with \(u_i\) or \(u_i^{-1}\) realized
on a generator.

For any integer \(n\),
\(u_i^n-1\in(u_i-1)O'\) in the coefficient valuation ring
\(O'\); for negative \(n\), use
\[
 u_i^{-n}-1=-u_i^{-n}(u_i^n-1).
\]
Thus the two places in
each pair have the same minimum, and the weighted cost is
\[
  t\sum_i 2\min_{\psi(g)\ne1}v(\psi(g)-1)
      =2t\sum_i v(u_i-1)
      =2t\sum_i\ell^{h_i}w \le \delta ,
\]
where the minimum at each \(i\) is restricted to its inertia
group.

\proofstep{Choose one specialization that fits all four allowances.}
All evaluations now use a common auxiliary parameter. Distributing the permitted error among the finitely many signatures and abelian values allows one specialization to satisfy all of the length and index bounds.

Take \(\delta\) smaller than \(\ell^{-2}/(1+2m+\#\{s\})\).
Writing \(L_s\) for each specialized signature length, the common choice gives
\[
 L_s\le\nu_\ell+e_s,\qquad e_s\ge0,\qquad \sum_s e_s\le\ell^{-2}.
\]
In particular every \(L_s\) is finite. The abelian
theta values are nonzero with valuations between \(-N_0\) and
\(v(\theta_\lambda(1))+\ell^{-2}\), the lower bound following
from integrality of \(f_M\). The nonnegative excesses over the
untwisted row indices sum to at most \(\ell^{-2}\).

The tame
cost multiplied by \(t\) is at most \(\ell^{-2}\) as well.
The number of primes, the depth, the coefficient extension and
the coordinate orders may depend on the host and the given
finite family. Only the four displayed allowances are used at
subsequent places. At conjugate coefficient places transport
the character with the prescribed inversion. For finitely many
other choices add all rational primes already used to the
avoidance set, taking all places over them. This gives the
disjoint auxiliary levels of the norm data. \label{cor:retained-twist-1}
\end{proof}

\subsubsection*{The relaxed global length at the chosen point}

Enlarging the coefficient field can realize the chosen character and its values. The Selmer cost in the determinant must still be measured in the original units of normalized length.

\begin{auxlemma}[Normalized torsion length under a DVR extension]
\label{lem:retained-dvr-normalized-length}
Let \(A\to C\) be a local inclusion of discrete valuation rings of mixed
characteristic \((0,\ell)\), let \(v_A,v_C\) give value one to the respective
uniformizers, and let \(D\) be a torsion \(A\)-module. Then
\[
 \frac{\operatorname{len}_C(C\otimes_A D)}{v_C(\ell)}
   =\frac{\operatorname{len}_A D}{v_A(\ell)}
\]
with infinite length allowed.
\end{auxlemma}

\begin{proof}
For a uniformizer \(\pi_A\), put \(e=v_C(\pi_A)\). The extension is flat
because \(C\) is torsion-free over the DVR. Tensoring a composition series
for a finite-length module replaces each simple factor by
\(C/\pi_A C\), of length \(e\). Also \(v_C(\ell)=e v_A(\ell)\).
An infinite-length torsion module has finite-length submodules of arbitrarily
large length. Flatness preserves their inclusions, proving the assertion
in that case.
\end{proof}

\begin{lemma}[Specialization on the actual retained lattice]\label{lem:retained-signature-specialization}
At a good retained point of the marked localization, the
selected lower complement on the prescribed lattice identifies
after coefficient extension with its good signature. The
relaxed Selmer group injects into that of the specialized
signature in the normalized comparison. In particular its
finite length satisfies
\[
 s_\pi\le\nu_\ell+\ell^{-2},\qquad
 t\sum_{v\in Q}v(\zeta_v-1)\le \ell^{-2}.
\]
\end{lemma}

\begin{proof}
\proofstep{Descend the rational intertwiner and use its nonzero integral reduction.}
At a good retained point \(\pi\), take the signature \(s\)
with the same lower cusp, tuple and side. Cancelling the
prescribed character factors in the row parameter
identifies the traces at split unramified Frobenius elements.
Chebotarev and continuity give
\[
 \operatorname{tr}\rho_\pi(g)=\operatorname{tr}\rho_s(g)
\]
for every \(g\), after scalar extension. Absolute irreducibility
of the reduction implies absolute irreducibility in characteristic
zero: a proper invariant subspace would have a proper saturated
intersection with a stable lattice after finite coefficient
extension, and reduction would give a proper nonzero invariant
subspace of the residual lattice. Brauer--Nesbitt gives
a fraction-field intertwiner, which descends because the
intertwiner equations are linear. Multiply it by a scalar so
that its entries are integral with at least one unit.

If \(A\) is the scaled intertwiner, its nonzero reduction has proper stable kernel. Residual irreducibility and equality of ranks give
\[
 \ker\bar A=0\quad\Longrightarrow\quad\det\bar A\ne0\quad\Longrightarrow\quad\det A\in O^\times.
\]
Thus \(A\) is an integral isomorphism. Twisting by
\(\chi_1^{-1}\) identifies \(W=\rho_\pi\chi_1^{-1}\)
with the signature lattice twisted by the inverse auxiliary
character of the chosen side. The signature lattice already
contains the inverse untwisted rank-one factor.

The two character evaluations identify their generated
orders by equality of kernels on the base group algebra.
Take the integral closure in their common fraction field.
Let \(A\to C\) denote either of the two local torsion-free
DVR extensions from this common base and \(\pi_A\) a
uniformizer. Fix lifts \((c_i)_i\) of a basis of \(C/\pi_A C\)
over \(A/\pi_A A\). They give, compatibly for every \(n\ge1\),
\[
 (A/\pi_A^n A)^{(i)}
     \xrightarrow{\ (x_i)\mapsto\sum_i c_i x_i\ }
                         C/\pi_A^n C ,
\]
where the superscript denotes the direct sum on the index set.
For surjectivity lift residues successively through order \(n-1\),
using finitely many terms at each order.

For injectivity, suppose \(\sum c_i x_i\in\pi_A^n C\) and some
\(x_i\notin\pi_A^n A\). For \(m=\min_i v_A(x_i)<n\), torsion-freeness
permits division by \(\pi_A^m\) in \(C\). Reduction gives
\[
 \sum_i\overline c_i\,\overline{x_i/\pi_A^m}=0\quad\text{in }C/\pi_A C.
\]
At least one coefficient is nonzero in \(A/\pi_A A\), contrary to the
choice of the residue basis. The same lifts work on all bounded torsion
submodules. Also \(C[1/\pi_A]=K_C\) since \(\pi_A\) has positive
\(C\)-valuation, so flatness gives
\[
 C\otimes_A(K_A/A)=K_C/C.
\]


If \(D\) is a discrete torsion coefficient module with this
action, continuous cochains commute with the indicated direct
sum. Each cochain has finite image on a finite power of the
compact Galois group, and consequently uses finitely many
copies. Differentials and restrictions act coordinatewise.
The auxiliary character is trivial on each decomposition group
over \(\ell\). The signature finite condition is the scalar span
of the base images, which is the direct sum of the base
conditions under this decomposition. For the finite local
coefficient fields the crystalline condition commutes with
extension by the finite basis decomposition of the rational
period comparison, and taking the image in the torsion quotient
commutes with it as well. Conjugate embeddings over the same
place transport this condition. Writing \(S_A\) for the Selmer module over the common base and \(S_C\) for either extension, compatibility of all local kernels gives
\[
 C\otimes_A S_A\simeq S_C.
\]

Lemma~\ref{lem:retained-dvr-normalized-length}, applied to the common
base Selmer group, identifies both normalized lengths, including the
criterion for finiteness.

Finally identify \(K^r/O^r\) with \((K/O)^r\) coordinatewise.
This carries the global Selmer group defining \(s_\pi\) into
the signature group defining the specialized determinant.
Unramified restrictions and unrestricted conditions at
\(S\cup Q\) away from \(\ell\) are preserved. At
\(v\mid\ell\) a generator of the local finite condition
comes from the reduction of a rational crystalline finite
cocycle. Write \(c\) for the chosen torsion cocycle and \(z\) for a rational finite cocycle. If \(c-\bar z=d a\), lift \(a\) to \(\tilde a\in K^r\); then
\[
 \overline{z+d\tilde a}=\bar z+d a=c .
\] Its coordinates then give the lifts in
the definition of the signature local condition, and taking
the integral span preserves membership. Coordinatewise
identification induces an isomorphism on ambient cohomology,
so the map on these Selmer subgroups is injective. In
particular finiteness of the specialized signature length
gives finiteness at \(\pi\) and
\[
          s_\pi\le \nu_\ell+\ell^{-2},
       \qquad t\sum_{v\in Q} v(\zeta_v-1)\le\ell^{-2}.
\]

\proofstep{Read the norm and signature characters on the same generator of inertia.}
Passing to a larger coefficient field preserves normalized length. We must further check that the character used to define the relaxed condition there is the one selected by the local norm data.

For the second inequality, the generator selected in the
marked model has full inertia image. The character of the
norm data and the twisted signature have the same character
values up to coefficient embedding and inversion, both of
which preserve the valuation of the difference from one.
These are the two global allowances in the local congruence
table. The pair-removal calculation derives the two theta polynomial assertions from their integral level and Euler relations. \label{cor:retained-twist-2}
\end{proof}

\subsubsection*{The square Selmer presentation}

The character choice used a polynomial coming from the Selmer group. Its construction begins with a square presentation over the group ring, so specialization gives a determinant on the prescribed auxiliary coordinates.

\begin{established}[Coefficient algebra and Selmer functor]\label{hyp:signature-presentation}
Put \(\Lambda=O_0[(\mathbb Z/\ell^M)^{2R_*}]\), with \(O_0\) the coefficient DVR of the signature, and let \(k\) denote the residue field. This algebra is finite free, local and complete. For the signature Selmer functor \(S\), with inverse universal character and relaxed conditions at the selected pairs, \(S\) is additive, \(\Lambda\)-linear and left exact. Every class is defined over a finitely generated torsion submodule, and \(S(k)\) is finite. For \(0\to k\to B\to C\to0\) one has
\[
 \#\operatorname{coker}(S(B)\longrightarrow S(C))\le\# S(k).
\]
\end{established}
\begin{proof}
The group elements form an \(O_0\)-basis. Indeed every maximal ideal lies
over the maximal ideal of \(O_0\), and after reduction one has
\[
 k[(\mathbb Z/\ell^M)^{2R_*}]\simeq k[X_1,\ldots,X_{2R_*}]/(X_1^{\ell^M},\ldots,X_{2R_*}^{\ell^M}).
\]
Here \(X_i\) is the \(i\)-th generator minus one, so the augmentation ideal is nilpotent.
Thus the maximal-ideal and base uniformizer topologies agree.
Its residue field \(k\) is finite.

 Additivity and linearity follow from the cocycle maps. Lemma~\ref{lem:signature-left-exact} proves left exactness and finite generation for individual classes. Finiteness over \(k\) follows from global cohomological finiteness for the finite module \(T(k)\). In \eqref{eq:sig-obstruction} a global connecting class and a tuple of local corrections give obstruction groups of orders \(|A|\) and \(|C|\), with product \(|S(k)|\). The fibre count following that identity gives the required bound, before the construction of the square presentation.
\end{proof}

\begin{lemma}[Square presentation of the signature functor]
\label{lem:signature-square-presentation}
For the ring and functor of Lemma~\ref{hyp:signature-presentation}, put
\(a=\dim_k S(k)\). There is one matrix \(D\) over the maximal ideal such that for every discrete torsion coefficient module \(B\) one has, naturally in \(B\),
\[
 D\in M_a(\mathfrak m),\qquad S(B)\simeq\ker(D:B^a\longrightarrow B^a).
\]
\end{lemma}

\begin{proof}
On finite modules use the exact duality
\(D(-)=\operatorname{Hom}_{\mathbb Z_\ell}(-,
\mathbb Q_\ell/\mathbb Z_\ell)\). The modules \(S(B)\) are
finite for finite \(B\) by a composition series and left
exactness. Hence \(J(B)=D S(D B)\) is a right exact linear
functor on finite modules.

For \(A_n=\Lambda/\mathfrak m^n\)
put \(P_n=J(A_n)\). If \(B\) is an \(A_n\)-module of finite
length then \(J(B)\simeq B\otimes_{A_n}P_n\) naturally.
For \(b\in B\) let \(\lambda_b:A_n\to B\) send \(1\) to
\(b\). The map \(b\otimes p\mapsto J(\lambda_b)(p)\)
is balanced and is the identity on \(A_n\). Additivity makes
it an isomorphism on finite free modules.

For a presentation \(A_n^q\xrightarrow{U}A_n^d\to B\to0\),
applying \(J\) gives
\[
       J(B)\cong \operatorname{coker}
                    (P_n^q\xrightarrow{U}P_n^d)
                \cong B\otimes_{A_n}P_n .
\]
The maps are induced by \(b\otimes p\mapsto J(\lambda_b)p\),
so the construction is natural in \(B\). In particular, view \(A_n\)
as an \(A_m\)-module when \(m\ge n\). This gives
\[
      P_n=J(A_n)\simeq A_n\otimes_{A_m}P_m
                       =P_m/\mathfrak m^nP_m .
\]
Write \(a=\dim_k S(k)=\dim_k P_1\), using \(Dk\simeq k\).
Lift a basis of \(P_1\) compatibly in the surjections \(P_m\to P_n\).

By Nakayama this gives surjections \(\Lambda^a\to P_n\) with
descending kernels \(K_n\). For \(m\ge n\),
\(K_n=K_m+\mathfrak m^n\Lambda^a\).
These kernels are compact in the finite base coefficient
topology.

For each \(x\in K_n\) the compact sets
\(K_m\cap(x+\mathfrak m^n\Lambda^a)\) are nested and
nonempty for \(m\ge n\). Consequently for \(K=\bigcap K_n\)
one has
\[
 K_n=K+\mathfrak m^n\Lambda^a.
\]
It follows that
\(P=\Lambda^a/K\) represents the original functor:
on a finite module killed by \(\mathfrak m^n\), dualizing the
tensor description gives
\[
 S(B)=\operatorname{Hom}_\Lambda(P,B)
\]
naturally. Every finite torsion module is killed by such a
power. The finite-generation condition on classes and left
exactness extend this equality to arbitrary discrete torsion
modules, since \(P\) is finitely generated and every finitely
generated torsion submodule is finite.

For the relations set
\(I=\operatorname{Hom}_{\mathbb Z_\ell}(\Lambda,
\mathbb Q_\ell/\mathbb Z_\ell)\). It is injective over
\(\Lambda\) by adjunction and divisibility over \(\mathbb Z_\ell\).
Finiteness of \(\Lambda\) over \(\mathbb Z_\ell\) makes \(I\)
torsion, and dualizing \(\Lambda\to k\) embeds \(Dk\simeq k\)
in it. Apply \(\operatorname{Hom}_\Lambda(P,-)\) to
\(0\to k\to I\to I/k\to0\). The connecting map is surjective
because \(I\) is injective, so the arithmetic cokernel is
\[
 \operatorname{coker}\bigl(S(I)\longrightarrow S(I/k)\bigr)
          =\operatorname{Ext}^1_\Lambda(P,k).
\]
Applying \(\operatorname{Hom}_\Lambda(-,k)\) to
\(0\to K\to\Lambda^a\to P\to0\) identifies this group
with \(\operatorname{Hom}_k(K/\mathfrak mK,k)\) because restriction
from \(\Lambda^a\) is zero when
\(K\subseteq\mathfrak m\Lambda^a\).
Its size bound implies \(\dim_k K/\mathfrak m K\le a\).

The ring is Noetherian, so at most \(a\) generators suffice
for \(K\). Use them as relations and pad by zero relations
to get an \(a\)-square matrix with entries in \(\mathfrak m\).
The kernel of its action on \(B^a\), taking the relations
as rows, is \(S(B)\) naturally in \(B\). The square matrix now represents the Selmer functor. Detecting its determinant in bounded degree is an arithmetic step: the auxiliary Frobenius elements must be realized by primes with prescribed ray coordinates. \label{cor:retained-twist-3}
\end{proof}


\subsubsection*{Ray characters and the two Frobenius prescriptions}
\label{subsec:aux-coordinates}
\hypertarget{source-auxiliary-coordinates}{}

The variables of the group-ring determinant must be realized by ray characters. A new prime must satisfy two Frobenius prescriptions, and both prescriptions have to survive the restrictions imposed at the distinguished twist prime.

Fix the auxiliary CM field \(E\), its conjugation \(c\) (with the chosen
extension of order two), the coefficient prime \(\ell\), and
\(L=\ell^M\), \(M\ge1\). Write \(T\) for the places over the
distinguished twist prime. Its residue characteristic is prime
to \(2L\). We first construct an open normal subgroup \(H\)
of \(G_E\), containing inertia above \(T\), and a finite set
\(S\) such that each split place \(P\notin S\) whose
decomposition group lies in \(H\) admits the coordinate
character of Section~\ref{subsec:retained-twist}.

\begin{proposition}[Prescribed Frobenius elements for the ray coordinates]\label{prop:ray-coordinate-pairs}
For the prescribed finite list of Frobenius conjugacy
conditions, the two groups of split auxiliary pairs can
be chosen outside any enlarged finite avoidance set.
Their anticyclotomic coordinates are surjective on the local
inertia of their own pair, unramified off that pair, and
trivial at decomposition groups above \(\ell\). Writing \(\eta_i\) for the additive first-group coordinates and \(P_j\) for the marked primes of the second group, one has
\[
 \eta_i(\operatorname{Frob}_{P_j})=\delta_{ij}\quad\text{in }\mathbb Z/L.
\]
\end{proposition}

\begin{proof}
\proofstep{Kill the global unit obstruction before prescribing a ray coordinate.}
Choose ideal representatives for the finite class group of \(E\)
prime to \(T\). This is possible by weak approximation at \(T\)
applied to the valuations of each representative. Put their prime
supports and the primes above \(\ell\) in \(S\), disjoint from \(T\).
Every fractional ideal has a generator away from \(S\). Choose
generators \(u_1,\ldots,u_a\) for the \(S\)-units, including
generators for their torsion. Take for \(H\) the subgroup fixing
\(E(\mu_L,u_1^{1/L},\ldots,u_a^{1/L})\). This is a finite
Galois extension. At \(v\in T\) every radicand is a unit
and \(L\) is invertible. The roots of each defining polynomial
lift its simple roots over the maximal unramified extension, so
inertia at \(v\) lies in \(H\).

For \(P\) as above, \(L\mid (\mathrm NP-1)\) and every \(u_i\)
reduces to an \(L\)-th power. Choose a surjection
\[
 \omega:k(P)^\times\to\mathbb Z/L
\]
from the cyclic residue
group. For an ideal \(J\) prime to \(P\), take a generator
\(\alpha\) away from \(S\). Its order at \(P\) is zero, so
\[
 J\mapsto\omega(\bar\alpha)
\]
is defined. Two generators
differ by an \(S\)-unit, and multiplication of generators proves
additivity. A principal ideal generated by \(1\bmod P\) has
value zero. Ray class reciprocity now gives a continuous
character \(\xi_P:G_E\to\mathbb Z/L\), unramified away from
\(P\). Its image on inertia at \(P\) is the image of \(k(P)^\times\):
on principal ideals prime to \(P\) the character evaluates the
residue character, and every residue unit lifts to an element
prime to \(P\).

At every \(v\mid\ell\) take \(1\) for a
generator of \(v\) away from \(S\). Both Frobenius and inertia
therefore act trivially there. Set
\[
       d_P(g)=\xi_P(g)-\xi_P(cgc^{-1}).
\]
Since \(P\ne cP\), only one summand ramifies at each member of
the pair. For the resulting difference character \(d_P\), the two inertia images are
\[
 d_P(I_P)=d_P(I_{cP})=\mathbb Z/L,
\]
including at \(\ell=2\).
This gives the desired anticyclotomic character, trivial on
decomposition groups above \(\ell\). \cite{milne-cft}

The universal coordinates in the theta data can be fixed once
per unordered pair. For each pair take a character with these
four properties (anti-invariance, triviality above \(\ell\),
ramification confined to the pair and surjectivity on its two
inertia groups) if one exists, and the trivial character otherwise.
Use this choice for the coordinate at \(P\) in every
conjugation-stable set containing it. Thus the ray construction
certifies existence at every admitted pair and the actual universal
coordinate there has the required properties. No uniqueness of a
ray character is required.
Restrict the chosen character to the local level of each selected series.

Here is the form of prime selection needed for the Galois
determinant. Let \(U\) be any open normal subgroup and let \(z\)
be the prescribed tame element at \(T\). Let
\[
       W=U\cap H\cap G_{E(\mu_{2L})}.
\]
We can prescribe the Frobenius of each first-half prime \(P_b\)
to be conjugate to \(z\) modulo \(W\), and the Frobenius of
each second-half prime \(P'_a\) to be conjugate modulo \(U\)
to any specified \(h_a\in H\cap G_{E(\mu_{2L})}\). In addition
we can require
\begin{equation}\label{eq:aux-coordinate}
       d_{P_b}(\mathrm{Frob}_{P'_a})=\mathbf1_{a=b}
                         \quad\text{in }\mathbb Z/L .
\end{equation}

To prove this, apply Chebotarev over \(E\) in the finite
quotient by \(W\), avoiding any given finite set, the ramification
loci, their conjugates and the primes above \(2L\). Degree-one
primes suffice: in the prime count over a number field primes of
higher absolute residue degree up to norm \(x\) contribute
\(O_E(\sqrt{x}\log x)\). A degree-one prime unramified over
\(\mathbb Q\) splits over its restriction to the maximal real
subfield because an inert quadratic place would have even residue
degree. The element \(z\) lies in \(H\) and fixes \(\mu_{2L}\)
by the inertia condition on \(T\). The primes selected have
\(\mathrm NP_b\equiv1\bmod 2L\) and are admitted. Their
inertia groups lie in \(W\). After each choice exclude the
whole pair.

\proofstep{Multiply each later Frobenius by inertia before applying Chebotarev.}
The new Frobenius must realize the prescribed values on the earlier ray coordinates as well.

The map
\[
 W\to(\mathbb Z/L)^R
\]
given by the \(R\) first-half
characters is surjective. Inertia at \(P_b\) realizes every
multiple of the \(b\)-th basis vector, and the target is abelian.
Choose \(e_a\in W\) so that the coordinates of \(h_a e_a\)
are the \(a\)-th basis vector. Apply Chebotarev again, now
modulo
\[
 W\cap\bigcap_b\ker d_{P_b},
\]
to \(h_a e_a\).
These primes are admitted by the same norm and splitting tests,
satisfy the prescribed conjugacies modulo \(U\), and satisfy
\eqref{eq:aux-coordinate} for every Frobenius lift. All exclusions
can be enlarged before either choice. This proves the prime and
character part of the construction at every depth. Bounds on
the norms of the auxiliary primes do not enter the length
estimate. The numbers of pairs and the degree tests uniform in
\(M\) come from the residual Selmer argument below.
\end{proof}


\subsubsection*{The Selmer group on a signature lattice}
\label{subsec:signature-selmer}

With ray coordinates available, the remaining inputs to the square presentation concern the Selmer functor. We verify them on the lattice for one good signature, retaining its fixed twist and its local conditions.

\begin{established}[Signature coefficients and local conditions]\label{hyp:signature-module-local}
Let \(T\) denote the rank-\(r\) signature lattice including its
prescribed scalar twist, \(O_0\) its coefficient DVR, and \(V\)
its fraction-field representation. Write \(k\) for the residue
field of the coordinate group ring \(\Lambda=O_0[(\mathbb Z/\ell^M)^{2R_*}]\) and \(\kappa\) for
the inverse universal character. Put
\[
 T(B)=(T\otimes_{O_0}B)\otimes\kappa
\]
for an
\(\mathfrak m\)-power torsion module \(B\). Let
\(\Sigma\) consist of the places above \(\ell\), the specified
type and distinguished places away from \(\ell\), and the
selected pairs \(Q\). Define \(G_\Sigma=\operatorname{Gal}(E_\Sigma/E)\), where \(E_\Sigma\) is the maximal extension unramified outside \(\Sigma\). The global signature classes unramified away from \(\Sigma\) are exactly the classes on \(G_\Sigma\). At the finite non-\(\ell\) places in \(\Sigma\) the conditions are relaxed, so the remaining local tests are at \(v\mid\ell\).
\end{established}
\begin{proof}
The underlying representation is unramified off its type places, the distinguished places and the places above \(\ell\), by Proposition~\ref{prop:signed-theta-data}. The universal character is unramified outside the chosen pairs and the prescribed ray places. Thus inertia outside \(\Sigma\) acts trivially on \(T(B)\). A cocycle representing an unramified class restricts to zero on each such inertia group because every coboundary there is zero. It vanishes on their conjugates and hence by continuity on the closed normal subgroup they generate. It therefore factors through \(G_\Sigma\); inflation gives the inverse, including for discrete torsion \(B\).
\end{proof}

\begin{lemma}[Variation in the signature coefficient module]\label{lem:signature-left-exact}
For \(T(B)\) and \(S(B)\) in Lemma~\ref{hyp:signature-module-local},
the finite local condition at \(v\mid\ell\) lifts under
surjections of \(\mathfrak m\)-power torsion coefficient
modules \(B\) and is reflected under injections. Every short exact coefficient sequence \(0\to B_1\to B_2\to B_3\to0\) induces an exact sequence
\[
 0\longrightarrow S(B_1)\longrightarrow S(B_2)\longrightarrow S(B_3).
\] Each of its classes comes from a
finitely generated coefficient submodule.
\end{lemma}

\begin{proof}
\proofstep{Describe the finite local condition as a span of integral cocycles.}
For \(v\mid\ell\) we have \(\kappa|_{G_v}=1\). Let
\(Z_v\) be the module of \(T\)-valued continuous cocycles on
\(G_v\) whose rational classes belong to \(H^1_f(E_v,V)\).
The local condition \(L_v(B)\) is the span in cohomology of
\[
 g\mapsto z(g)b,
\]
for \(z\in Z_v\) and \(b\in B\),
using a basis of \(T\) to read this formula coordinatewise.
All coboundaries lie in the corresponding span of cocycles:
use the integral coboundaries of the coordinate basis vectors
of \(T\). Thus a class belongs to \(L_v(B)\) precisely when
any representing cocycle is such a finite sum. Maps of
coefficient modules preserve this description. For a coefficient surjection \(q:B_1\twoheadrightarrow B\), take any lift \(b_1\) of \(b\). Termwise lifting follows from
\[
 q(z(g)b_1)=z(g)q(b_1)=z(g)b.
\]


\proofstep{Use saturation to read each local relation before enlarging coefficients.}
The finite local condition was described after passing to the rational representation. Saturation identifies the integral relations that must be preserved on extending the coefficient module.

Injections of coefficient modules reflect this condition.
Here the integral assertion is at the level of functions on
\(G_v\times\{1,\ldots,r\}\). The module \(Z_v\) is saturated
in all integral functions: if \(az\in Z_v\) with \(a\ne0\),
division in the rational vector space preserves the cocycle
law, continuity and the finite local condition. Suppose a
\(B'\)-valued function has image \(\sum_{j=1}^n z_j b_j\)
under an inclusion \(B'\hookrightarrow B\). Discard zero
terms.

Take an entry \(a=z_1(i_0)\) of least valuation and put
\[
 u=z_1/a\in Z_v,
\]
with \(u(i_0)=1\). For the original function \(f\) with values in \(B'\), let
\(\iota:B'\hookrightarrow B\) be the coefficient inclusion.
With \(f(i_0)\in B'\), the subtraction gives
\[
 \iota\!\left(f-u f(i_0)\right)
  =\sum_{j>1}\bigl(z_j-u z_j(i_0)\bigr)b_j .
\]
Each new cocycle belongs to \(Z_v\), and the representation
has one fewer term. Induction puts \(f-u f(i_0)\) in the
required span over \(B'\), which also contains \(u f(i_0)\). This argument works when
\(B'\subset B\) has torsion quotient.

\proofstep{Kill global invariants by reduction to the coefficient socle.}
The absolute residual complement is irreducible of dimension
\(r>1\), by the retained type tests. On a simple coefficient
module the residual scalar twist preserves irreducibility and
\(\kappa\) reduces to one. Hence
\[
 H^0(E,T(k))=0,
\]
also for
the Tate dual. For arbitrary \(B\), a nonzero invariant vector
would generate a cyclic module killed by a power of
\(\mathfrak m\). Its last nonzero maximal-ideal power would
give a nonzero invariant in the socle, where the coefficient
space is a sum of copies of \(k\). This proves the same
vanishing on all torsion modules.

The long exact sequence of
global cohomology now gives injectivity on \(H^1\) for
inclusions and exactness at the middle for a short exact
coefficient sequence. The local reflection just proved gives
left exactness of \(S\). A continuous cocycle on a compact
group with discrete target has finite image, so its coordinates
generate a finitely generated torsion submodule. Local reflection
puts the class into the Selmer group of that submodule. This
proves the finite-generation assertion used in constructing
the representing module.
\end{proof}

\paragraph{Cardinalities of the finite local conditions.}
The good-signature isotypic condition gives \(H^0(E_v,V)=0\)
at \(v\mid\ell\), and \(V\) is de Rham there. Set
\(d_v=\dim H^1_f(E_v,V)\), the dimension over the coefficient
field. The Bloch--Kato dimension formula gives
\[
 d_v=\sum_{\tau\mid v}\dim\big(D_{\rm dR}(V)_\tau/
                                      \operatorname{Fil}^0\big).
\]
Here the coefficients split the embeddings, and in the formula
each summand is the dimension of the indicated quotient.
The signature filtration at conjugate embeddings has complementary
quotient dimensions, with sum \(r\). As \([E:\mathbb Q]=2f\),
let \(h_\tau\) be the quotient dimension at \(\tau\). Summation gives
\[
 \sum_{v\mid\ell}d_v
   =\sum_{\{\tau,\bar\tau\}}(h_\tau+h_{\bar\tau})
   =\frac{[E:\mathbb Q]}2\,r=fr .
\]
\cite{bloch-kato}

The fraction-field space of finite cocycles has dimension
\(r+d_v\), since coboundaries have dimension \(r\). Its
intersection \(Z_v\) with integral functions is a full free
lattice. Compactness bounds each of a finite basis of rational
cocycles and hence clears its denominators. For finite generation,
choose finitely many evaluations whose coordinate functionals
span the dual of the rational space, giving an injection of
the integral functions into a finite free module. Reduction of
\(Z_v\) modulo its uniformizer injects into the functions
over the residue field by saturation. Consequently the span
of these cocycles with coefficients in \(k\) has cardinality
\(|k|^{r+d_v}\). Dividing by the residual coboundaries, of
size \(|k|^r/|H^0(E_v,T(k))|\), gives
\begin{equation}\label{eq:signature-local-size}
  \prod_{v\mid\ell}|L_v(k)|
   =|k|^{fr}\prod_{v\mid\ell}|H^0(E_v,T(k))|.
\end{equation}
This calculation allows the reduction to have local invariants.

\paragraph{Counting the obstruction to lifting a Selmer class.}
Put \(D=T(k)\), \(m=|D|^f\), and evaluate cohomology on
\(G_\Sigma\). Define
\[
 \begin{aligned}
 A&=\ker\Big(H^2(E_\Sigma/E,D)\longrightarrow
                   \prod_{v\mid\ell} H^2(E_v,D)\Big),\\
 C&=\operatorname{coker}\Big(H^1(E_\Sigma/E,D)
       \longrightarrow\prod_{v\mid\ell} H^1(E_v,D)/L_v(k)\Big).
 \end{aligned}
\]
These groups are finite by local and global finiteness for a
finite module, a finite ramification set containing all primes
over \(\ell\), and a totally imaginary field. The last
localization in degree two is surjective by Poitou--Tate:
the cokernel of localization at all places of \(\Sigma\)
is dual to the invariants of the Tate dual, which vanish.
Projection to the places over \(\ell\) preserves surjectivity.
The global Euler characteristic gives \(|H^1|=m|H^2|\)
because there are \(f\) complex places and \(H^0=0\).
The product of the local Euler characteristics at \(\ell\)
gives
\[
 \prod_{v\mid\ell}|H^1(E_v,D)|=
 m^2\prod_{v\mid\ell}|H^0(E_v,D)|\,|H^2(E_v,D)| ,
\]
where \(\sum_v[E_v:\mathbb Q_\ell]=2f\). Taking kernels
and cokernels of the two displayed localization maps and using
\eqref{eq:signature-local-size} yields \cite{milne-duality}
\begin{equation}\label{eq:sig-obstruction}
                        |S(k)|=|A|\,|C|.
\end{equation}

For
\[
 0\to k\to B\to B''\to0,
\]
the connecting map on
\(S(B'')\) has image in \(A\), since every allowed local class
lifts to an allowed local class. On its kernel choose a global
lift to \(H^1(G_\Sigma,T(B))\). Subtract an allowed local
lift at each \(v\mid\ell\): the difference comes from
\(H^1(E_v,D)\). Its tuple of preimages, modulo the local
conditions and global image, records the remaining obstruction.
For the cardinality bound one can allow choices: for each
class in the kernel choose one global lift, local lifts and
preimages.

If two of the resulting classes in \(C\) agree,
subtract a global \(D\)-class and use the local \(L_v(k)\)
to put the difference of the global lifts in \(S(B)\).
There are at most \(|C|\) cosets in the kernel modulo the Selmer image and at most \(|A|\) fibres. With \eqref{eq:sig-obstruction} this gives
\[
 \#\operatorname{coker}\bigl(S(B)\to S(B'')\bigr)\le |A|\,|C|\le |S(k)|.
\]
Together with left exactness and passage
to finitely generated submodules this supplies the hypotheses
of the square-presentation proof for every selected coordinate
group ring. Its size may depend on the signature and the
auxiliary level.

\subsubsection*{The residual coefficient of the signature determinant}
\label{subsec:residual-twist}

The signature lattice now gives a square matrix. To control its determinant at a specialization, we need a coefficient detectable from the mod-\(\ell\) Selmer classes and the Frobenius elements chosen for them.

Keep the notation of Section~\ref{subsec:signature-selmer} and write
\(W=T\otimes k\), \(W'=W^\vee(1)\). Conjugation of a group
element by complex conjugation is denoted by \(h\mapsto h^c\).
The good-signature duality gives an integral isomorphism of
\((T^\vee(1))^c\) with \(T\otimes\epsilon^t\), where \(\epsilon\)
is cyclotomic. The tame marking gives an inertia element \(z\)
acting on \(W\) by \(d\), with \(d^2\ne1\), and on \(W'\) by
\(d^{-1}\). At \(z^c\) the two scalars are reversed.
Both modules are absolutely irreducible. Let \(S,S'\) be the
residual Selmer groups for these modules at empty auxiliary
level, with finite local conditions at \(\ell\) defined by
the integral lattices.

\begin{lemma}[Residual Selmer dimensions]
For the residual primal and dual Selmer spaces of a good signature over its residue field \(k\),
\[
 \dim_k S=\dim_k S'.
\]
\end{lemma}
\begin{proof}
At \(v\mid\ell\), the cup of two integral finite cocycles has
zero rational invariant by Bloch--Kato orthogonality. The local
invariant map embeds integral \(H^2\) in rational \(H^2\) on the
scalar Tate twist, hence the reductions are orthogonal. For the
lattice of \(W\), Section~\ref{subsec:signature-selmer} gives
\[
 \dim L_v(k)-\dim H^0(E_v,W)
   =\dim H^1_f(E_v,V)-\dim H^0(E_v,V).
\]
The right side is the dimension of the de Rham quotient by
\(\operatorname{Fil}^0\). Add this equality to its analogue for
the dual lattice. At each embedding the two quotient dimensions
sum to \(r\) by the dual Tate filtration. Local duality identifies the invariants of the Tate dual with the
dual of \(H^2(E_v,W)\). Thus the local Euler characteristic reads
\[
 \dim L_v(k)+\dim L'_v(k)
 = [E_v:\mathbb Q_\ell]r+\dim H^0(E_v,W)+\dim H^2(E_v,W)
 =\dim H^1(E_v,W).
\]
The orthogonal subspaces therefore have complementary dimensions.

At the type and distinguished places
away from \(\ell\) the inertia test of the signature excludes
invariants on the residual modules: an
inertia element acts with no eigenvalue one. The dual has
the inverse eigenvalues there. The local Euler characteristic
and duality give \(H^1=0\). Away from these places and
\(\ell\), the unramified subgroups are exact annihilators.
Global duality (the finite-module Selmer order formula) now
gives \(|S|/|S'|=1\): the global invariant factors are one,
the complex factors contribute \(|W|^{-f}\), and the
\(\ell\)-adic factors contribute \(|W|^f\) by
\eqref{eq:signature-local-size}. \cite{bloch-kato,milne-duality}

\end{proof}
\paragraph{Forms specified by Frobenius elements.}
There is an open normal subgroup \(J\subset H\), stable under
\(c\), on which the residual representations and the action
on \(\mu_{2L}\) are trivial and which contains all commutators
\([g,z]\). Intersect \(H\) with the cyclotomic and residual
kernels and with the conjugate of the intersection. The
commutators lie in these groups because \(H\) contains
inertia at the distinguished pair, cyclotomic is abelian,
and the marked elements act by scalars. The same construction
works simultaneously for finitely many signatures. If needed
one may replace each residual kernel here by the kernel modulo
\((\ell^{\nu_\ell})\mathfrak m_{O_0}\): the distinguished
inertia acts by scalars integrally on \(T\).

For \(h\in J\) put
\[
       B_h(x,y)=x(h)\cdot y(h)-x(h^c)\cdot y(h^c)
            \qquad (x\in S,\ y\in S').
\]
\begin{lemma}[Corank of the Frobenius pairing]
\label{lem:residual-frobenius-forms}
For every good signature there is a linear combination \(B\) of the second-half Frobenius forms for which
\[
 \dim\ker\bigl(S\longrightarrow(S')^*,\ x\longmapsto B(x,\cdot)\bigr)\le \mathbf1_{\nu_\ell>0}.
\] The same second-half primes can be
used for the finitely many signatures. Before combining their
lists, each requires at most \((\dim S+\dim S')^2\) elements of \(J\).
\end{lemma}
\begin{proof}
The values are independent of the cocycle representatives.
Restriction to \(J\) is injective on cohomology. In fact a
cocycle \(x\) vanishing there factors through a quotient in
which \(z\) is central. The two expressions for \(x(gz)=x(zg)\)
give
\[
 (d-1)x(g)=(g-1)x(z),\qquad
 x(g)=(g-1)\bigl((d-1)^{-1}x(z)\bigr).
\]
Thus \(x\) is a coboundary. On the dual use \(d^{-1}\) in the
same identity; both denominators are nonzero. For independent classes
\(x_1,\ldots,x_a\) with values in \(W\), the vectors
\((x_j(h))_j\) span \(W^a\) over \(k\) as \(h\) varies in
\(J\). Let \(U\) be the span of the evaluation tuples in \(W^a\). It is stable by conjugation. If \(U\) were proper, semisimplicity and absolute irreducibility would give scalars, not all zero, with
\[
 \sum_{j=1}^a c_j x_j(h)=0\qquad(h\in J).
\]
Indeed any map \(W^a\to W\) on this isotypic space is a linear combination of coordinate projections. Restriction injectivity rules it out.

If \(W\not\simeq (W')^c\), then
\(W,W',W^c,(W')^c\) are pairwise inequivalent: the scalars
at \(z\) separate the two groups and conjugation carries the
remaining inequivalence to its partner. For bases of \(S,S'\)
the simultaneous evaluations of these four lists on \(h\)
span the direct sum of their four full isotypic spaces. This
follows from the preceding projection test on each simple
type. To show that the \(B_h\) span all bilinear forms, suppose
\[
 \sum_{i,j}a_{ij}B_h(x_i,y_j)=0
\]
for every \(h\in J\).

Write the four evaluation blocks as \(v_i,w_j,\tilde v_i,\tilde w_j\),
so the relation says that every evaluation tuple is a zero of
\[
 Q(v,w,\tilde v,\tilde w)=\sum_{i,j}a_{ij}
      (v_i\cdot w_j-\tilde v_i\cdot\tilde w_j).
\]
Evaluation on \(J\) is additive, so the bilinear form
\(B_Q(s,t)=Q(s+t)-Q(s)-Q(t)\) vanishes on all pairs of evaluation
tuples and hence on their full linear span. Fix \(i,j\); give
\(s\) just the entry \(v\) in the \(i\)-th first block and
\(t\) just the entry \(w\) in the \(j\)-th second block. Then
\begin{equation}\label{eq:frobenius-polarization-block-test}
 0=B_Q(s,t)=a_{ij}\,v\cdot w .
\end{equation}
Choosing a pair with nonzero pairing in \eqref{eq:frobenius-polarization-block-test} forces \(a_{ij}=0\), including
in characteristic two.

If \(W\simeq(W')^c\), choose an invertible intertwining
matrix \(A\), so that
\(W(g^c)^t A W(g)=\bar\epsilon(g) A\), with superscript
\(t\) on a matrix denoting transpose. Transposing the same
identity at \(g^c\) shows \(A^t=\lambda A\) by Schur's lemma.
Here \(\lambda^2=1\) and \(\lambda^r=1\) by determinants, so
\(A^t=A\) since \(r\) is odd. Transport \(S'\) to a subspace
\(S_2\) of \(H^1(E,W)\) by sending \(y\) to \(A^{-1}y^c\).
Choose a basis of \(S+S_2\) and evaluate it at \(h,h^c\).
The two evaluations span the two full isotypic blocks together
because \(z\) separates \(W\) and \(W^c\).

In this basis the
form has matrix entries
\[
     x_i(h)^t A x_j(h^c)-x_i(h^c)^t A x_j(h).
\]
A relation with coefficients \(a_{ij}\) depends on
\(a_{ij}-a_{ji}\). Polarization and testing on vectors in
the two separate blocks forces \(a_{ij}=a_{ji}\). Conversely
these symmetric relations annihilate the displayed alternating
matrices, including their diagonals. Thus the forms span the restrictions of alternating forms on
\(S+S_2\). Set \(J_0=S\cap S_2\) and take complements
\(S=J_0\oplus U\), \(S_2=J_0\oplus U'\). These complements
have the same dimension. With dual bases for the pairing
of \(U\) with \(U'\), choose the restriction \(S\times S_2\)
with block matrix
\[
 \begin{pmatrix}J_{\dim J_0}&0\\0&I_{\dim U}\end{pmatrix},
 \qquad
 J_d=\bigoplus_{i=1}^{\lfloor d/2\rfloor}
       \begin{pmatrix}0&1\\-1&0\end{pmatrix}
       \quad\text{with a final zero block if \(d\) is odd}.
\]
The opposite off-diagonal pairing on \(U'\times U\) is the
negative transpose, so this matrix extends to an alternating
form on \(S+S_2\). Its left corank is \(d\bmod 2\), including
in characteristic two.

Consequently there is a linear combination of the forms
with left corank at most \(e\), where \(e=0\) if
\(\nu_\ell=0\) and \(e=1\) otherwise. Indeed an equivalence
\(W\simeq(W')^c\) would imply
\[
 \bar\epsilon^{tr}=1
\]
by taking determinants in the signature duality, which
implies \(\nu_\ell>0\) by \eqref{eq:signature-cyclotomic-minimum}.
The number of elements \(h\) needed is at most
\((\dim S+\dim S')^2\), independently of \(J\) and \(M\).
Choose such elements for each good signature and concatenate
the lists, padding with identities if necessary.

For the finite list \(\mathcal Z\) of restrictions of the basis cocycles and their conjugates, use the common kernel
\[
 J_0=\bigcap_{z\in\mathcal Z}\ker(z|_J).
\]
The conjugation formula gives an open normal subgroup of \(J\).
Apply Section~\ref{subsec:aux-coordinates} prescribing these
elements for the second half. On replacing \(h\) by a
Frobenius conjugate to it modulo this subgroup the matrix
\(B_h\) changes only by the nonzero cyclotomic scalar of the
conjugating element. Absorb it into the coefficients. Thus at
the second-half Frobenius elements \(\phi_i\) we have, for each signature,
coefficients \(b_i\) such that \(\sum_i b_i B_{\phi_i}\) has
the stated corank. The same primes work for all signatures.

\end{proof}
\begin{lemma}[A residual coefficient from the Selmer pencil]\label{lem:residual-selmer-coefficient}
Use the two halves of auxiliary places and the Frobenius forms
of Lemma~\ref{lem:residual-frobenius-forms}, with \(R\) places
in each half and \(L>4rR+2\). When \(\nu_\ell=0\), every polynomial determinant lift \(D\) satisfies
\[
 \exists\alpha:\quad |\alpha|\le4rR,\qquad [X^\alpha]D\in O^\times.
\] For positive \(\nu_\ell\) the corresponding
unit coefficient is obtained at this stage in an adjugate
entry, for a nonempty square presentation. An empty presentation
already has determinant one.
\end{lemma}

\begin{proof}
\proofstep{Use the inertia and Frobenius evaluations in the relaxed residual space.}
Let \(R\) denote the common length of each half. At the first
half the untwisted residual module is unramified with Frobenius
scalar \(d\) or \(d^{-1}\); both differ from one. The
cyclotomic character is trivial there modulo \(\ell\).
Thus \(H^0,H^2\), and by the local Euler characteristic also
\(H^1\), vanish. At each second-half prime the Frobenius is
in \(J\), as is inertia. Fix a tame \(\ell\)-generator
\(\gamma_i\) with \(d_{P'_i}(\gamma_i)=1\). This is possible
since the character on inertia is onto. In the following
evaluation pairs, use at the conjugate place the transports
\(\gamma_i^c,\phi_i^c\). Superscript \(c\) on group elements
continues to denote conjugation.

Write \(S_Q\) for the relaxed residual Selmer space and
\[
 D=(k^r\oplus k^r)^R.
\]
Define \(I:S_Q\to D\) by evaluation
at \((\gamma_i,\gamma_i^c)_i\). Its kernel is \(S\), since
on the trivial residual module a homomorphism on inertia
factors through its pro-\(\ell\) tame quotient. The first-half
places contribute nothing. For \(x\in S_Q\), \(y\in S'\),
reciprocity gives units \(a_i\in k^\times\), independent of
\(x,y\), such that
\begin{equation}\label{eq:residual-reciprocity}
 \sum_i a_i\{x(\gamma_i)\cdot y(\phi_i)
             -x(\gamma_i^c)\cdot y(\phi_i^c)\}=0 .
\end{equation}
Indeed the cup invariant at an unramified pair of classes is
zero, and at the primes over \(\ell\) it vanishes by the local
orthogonality above.

At a second-half place the second class
is unramified. The invariant on a tame character of
\(\gamma_i\) paired with an unramified character of \(\phi_i\)
is a nonzero scalar: the two unramified classes cup to zero
and local Tate duality is nondegenerate on the two-dimensional
space for the trivial residual line. Bilinearity gives the
formula for vectors. The transport to the conjugate place
reverses the chosen identification of \(\mu_\ell\) with the
prime field, since \(c\) acts by \(-1\) there. This explains
the common scalar and the minus sign. The sum of local
invariants vanishes by global reciprocity.

Define
\[
 F:S_Q\to D
\]
by taking the pairs of Frobenius
values of each class multiplied by \(b_i/a_i\). The two necessary linear conditions on a lift are
\[
 Ix=0,\qquad Fx=Ix_1\quad(x_1\in S_Q).
\]
The first gives \(x\in S\) and \eqref{eq:residual-reciprocity} applied
to \(x_1\) puts \(x\) in the left radical of
\(\sum b_i B_{\phi_i}\). Consequently
\begin{equation}\label{eq:residual-pencil}
  \dim\{x\in\ker I: Fx\in\operatorname{im}I\}\le e,
  \qquad \dim S_Q\le 2rR+e .
\end{equation}
For the second inequality apply rank-nullity to
\(\ker I\to D/\operatorname{im}I\) induced by \(F\).

We give the specialization producing a coefficient of bounded
degree; setting all group variables to the same variable would
lose the coefficients \(b_i/a_i\). Put \(K=k((s))\).
In the polynomial presentation matrix reduce modulo the
uniformizer and set \(X_j=\alpha_j Z\), \(Y_i=\beta_i Z\),
where
\[
  \beta_i=s\, b_i/a_i,\qquad
  \alpha_j=1-\sum_i\beta_i\,\overline{d_{P'_i}(\phi_j)} .
\]
Write \(P(Z)\) for this square matrix over \(K[Z]\).
Modulo \(Z^2\) the substitution factors through the group ring,
since \(L\) is a positive power of \(\ell\). On the second-half
decomposition group the coefficient of \(Z\) in the inverse
universal character is \(-1\) at \(\phi_i\) and \(-\beta_i\)
at \(\gamma_i\), by \eqref{eq:aux-coordinate}. At the
conjugate place both signs reverse.

A first-order lift of \(v\in\ker P(0)\) is a vector \(v+Zv_1\) satisfying
\[
 P(Z)(v+Zv_1)\equiv0\pmod{Z^2}.
\]
Naturality of the square presentation gives
a Selmer cocycle in the dual numbers lifting its residual
class \(x\in K\otimes_k S_Q\). This identification commutes
with taking its constant term. Indeed additive functoriality
gives the assertion for finite sums of \(k\), and each class
over a discrete extension field is supported on a finite
coefficient subspace.

At each second-half place write the
cocycle as \(x+Zx_1\), and the character as \(1+Z\lambda\).
The tame relation
\[
 \phi\gamma\phi^{-1}=\gamma^q,
\]
where
\(q=\mathrm NP'_i\), gives
\[
             \lambda(\phi)x(\gamma)=\lambda(\gamma)x(\phi).
\]
For this calculation the crossed term of the first-order
cocycle at \(g,h\) is \(\lambda(g)x(h)\). In its value at
\(\gamma^q\) the coefficient \(\binom q2\) vanishes in \(k\)
because \(q\equiv1\bmod 2L\); the coefficient \(q\) is one.
The prime-to-\(\ell\) inertia kernel has zero image in the
successive additive groups used in this calculation.
Thus \(Ix=sFx\), including at the conjugate places.

The subspace of \(K\otimes S_Q\) satisfying this equation
has dimension at most \(e\). Intersect it with the lattice
\(S_Q[[s]]\). This saturated sublattice has reduction of
dimension equal to its rank and may be given a basis with
independent constant terms by elimination over \(k[[s]]\).
Coefficient comparison in \(Ix=sFx\), with \(x=x_0+sx_1+\cdots\), gives
\[
 Ix_0=0,\qquad Ix_1=Fx_0.
\]
Now apply
\eqref{eq:residual-pencil}. This argument can be used over
any finite residue field, including \(\mathbb F_2\).

\proofstep{Use the first-order kernel to bound the number of invariant factors of order at least two.}
The evaluations have bounded the kernel of the first-order matrix. Reading the same matrix through its invariant factors converts that kernel bound into the restriction on the determinant's low-degree coefficients.

Take Smith form over \(K[[Z]]\), with orders
\(0\le a_1\le\cdots\le a_d\le\infty\), where a zero
diagonal entry has order \(\infty\). Reduce the kernel of the \(i\)-th diagonal entry on \(K[[Z]]/(Z^2)\)
modulo \(Z\). Its image in \(K\) is
\[
 \begin{cases}
 0,&a_i=0\text{ or }a_i=1,\\
 K,&2\le a_i\le\infty.
 \end{cases}
\]
For \(a_i=1\) its kernel is precisely the ideal \((Z)\), whose
residue in \(K\) is zero.
The dimensions just computed give
\[
 \#\{i:a_i>0\}=h=\dim S_Q\le2rR,
 \qquad \#\{i:a_i\ge2\}\le e.
\]
If \(e=0\), all Smith exponents are zero or one. Therefore
\[
 \operatorname{ord}_Z\det P=\sum_i a_i=\#\{i:a_i=1\}=h .
\]
If \(e=1\) and \(d>0\), the ideal generated by the submaximal
minors has order
\[
 \sum_{i=1}^{d-1}a_i\le\max(0,h-1)\le2rR.
\]

At least one adjugate entry has that order. For an empty matrix
the determinant is already one. Both conclusions suffice with
the degree bound \(4rR\). For \(L>4rR+2\) the
choice of polynomial lift cannot change any reduced coefficient
in these degrees: the coordinate relations reduce to
\(X_j^L=Y_i^L=0\). A nonzero coefficient of the line
specialization requires a nonzero coefficient of the same total
degree before substitution. Hence every lift of the determinant
has a unit coefficient of bounded degree when \(\nu_\ell=0\).
For \(\nu_\ell>0\) the conclusion at this stage concerns one
adjugate entry. Its conversion to a determinant with valuation
at most \(\nu_\ell\) requires one further arithmetic step.
\end{proof}


\subsubsection*{The positive cyclotomic valuation}
\label{subsec:positive-twist}

Suppose the cyclotomic minimum \(\nu=\nu_\ell\) of
\eqref{eq:signature-cyclotomic-minimum} is positive. Write \(O=O_0\), choose a
uniformizer \(\pi\), and put
\[ e=\operatorname{ord}_O(\ell^\nu),\qquad I=\ell^\nu\mathfrak m_O=(\pi^{e+1}). \]
The residual calculation detects an entry of the adjugate. One elementary divisor can still carry the entire remaining loss, so the determinant itself must now be detected modulo \(I\).

Let \(b\) be the number of good signatures at this coefficient
place, and \(R_0\) the sum of the bounds for the number of Frobenius elements in their
residual tests, padding each list with identities. Use \(R=R_0+b\) first pairs, and initially
\(n=R_0\) second pairs. Prescribe the first \(R_0\) Frobenius elements of
the second list as in Section~\ref{subsec:residual-twist}.
Future second pairs will use the remaining first coordinates
one at a time. Write \(p=2r\) and set
\[
 w_0=4rR_0,\qquad w_{j+1}=2(e+1)w_j+2+p .
\]
Only the finite signature data enter these constants. Throughout
take \(M\) so large that \(\ell^{M-\nu-1}>w_b+2\) and
\(M\ge\nu+2\).

\begin{lemma}[Independence of low-degree coefficients]
Let \(F,G\) lift the same group-ring element. After reduction to \(O/I\), or to the residue field, their coefficients agree in the range
\[
 [X^\alpha]F=[X^\alpha]G\qquad(|\alpha|\le w_b+1).
\]
\end{lemma}
\begin{proof}
For \(0<i<\ell^M\) one has
\[
 i\binom{\ell^M}{i}=\ell^M\binom{\ell^M-1}{i-1},\qquad
 v_\ell\binom{\ell^M}{i}=M-v_\ell(i),
\]
since \(\binom{\ell^M-1}{i-1}\) is a unit by expansion modulo \(\ell\).
For \(1\le i\le w_b+1<\ell^{M-\nu-1}\) the coefficient belongs
to \(I\). Each coordinate relation \((1+U)^{\ell^M}-1\), and hence
each polynomial multiple, has zero coefficients in these degrees.
\end{proof}

The induction keeps a unit coefficient in a cofactor through
degree \(w_j\) for every signature with a nonempty presentation,
after setting to zero the unused first variables. For each of
the first \(j\) signatures it also keeps a determinant coefficient
outside \(I\) through degree \(w_j\), with all variables allowed.
To start, set the extra first variables to zero. At these places
the residual Frobenius is scalar different from one on either
side, and its norm is one modulo \(\ell\). Writing \(D\) for that residual representation, local duality and the local Euler characteristic give
\[
 H^j(E_v,D)=0\qquad(j=0,1,2).
\]

Successive extensions give the same vanishing on finite
coefficient modules after this specialization. Deleting the
extra pairs therefore gives the old Selmer functor. Its representing
module after scalar change is the old module by naturality.
The first Fitting ideal (the ideal generated by the submaximal
minors of a square presentation) commutes with this base change.
The unit cofactor from Section~\ref{subsec:residual-twist}
therefore supplies the initial assertion for any presentation.
An empty presentation has determinant one; it can also be
stabilized by a unit diagonal entry.

\paragraph{Preserving a coefficient at a larger level.}
Consider a new second pair satisfying the coordinate conditions
and lying in the subgroup \(J\) for every signature.
\begin{lemma}[Fitting ideals at an added pair]
After setting the new second variable to zero, let \(P_+\) and
\(P_0\) be the new and old modules reduced modulo \(I\).
For the Frobenius \(\sigma\) at the new prime put \(u=\eta(\sigma)-1\).
Then
\[
 u^p\operatorname{Fitt}_i(P_0)\subseteq\operatorname{Fitt}_i(P_+)\qquad(i=0,1).
\] A nonzero determinant coefficient modulo \(I\), or
a unit cofactor coefficient modulo \(\pi\) with unused variables
zero, can be retained with degree increased by at most \(p\)
in the asserted low-degree range.
\end{lemma}
\begin{proof}
First set its own variable to zero and reduce modulo \(I\).
Write \(P_+\) for the new representing module after this
specialization and \(P_0\) for the old module. There is a
surjection
\[
 P_+\to P_0
\]
with kernel generated by \(p\)
inertia coordinates. In fact in the new Selmer functor the
unramified classes at the new pair are exactly the old ones.
The old action is unramified there; the only possible
inertia homomorphisms on finite primary modules factor through
the pro-\(\ell\) tame quotient. Evaluating a representative at
one generator on each side gives the claimed rows in \(P_+\),
independent of coboundaries. The kernel description now follows
by testing on all finite modules. Here the specialized rings
are finite.

Put \(u=\eta(\sigma)-1\) for the old universal character on
a Frobenius at the new prime, choosing the direct character
\(\eta\). For the inertia generators \(t_1,\ldots,t_p\) of the kernel this gives the scalar relations
\[
 (\eta(\sigma)-1)t_i=u t_i=0\qquad(1\le i\le p).
\]

Indeed the untwisted Frobenius matrix is one modulo \(I\) on
\(J\) and the tame relation multiplies an inertia value by
the residue norm, which is one in \(O/I\). Inversion on the
conjugate side or in the twist changes this annihilator by
a unit. If \(A_0\) is a presentation matrix for \(P_0\), lift its
relations to the enlarged list of generators for \(P_+\).
Together with the inertia relations these give a block of the
new presentation of the form
\[
        \begin{pmatrix} \widetilde A_0 & *\\ 0 & u I_p\end{pmatrix}
\]
when rows write relations. Adjoining all \(p\) inertia rows
and columns to a minor for \(\operatorname{Fitt}_i(P_0)\)
multiplies it by \(u^p\).

For \(i=0,1\) this proves
\[
        u^p\operatorname{Fitt}_i(P_0)
                \ \subseteq\ \operatorname{Fitt}_i(P_+).
\]
It suffices that these are relations; further relations enlarge
the ideal. One may stabilize either presentation for \(i=1\).
The polynomial for \(u\) has constant term zero and its
linear coefficient at the newly used first variable is one
by the Frobenius coordinate prescription. Even after deleting
the still unused first variables its linear form is monic in
that variable. Write \(U_1\) for its linear part. Over \(O/I\), a monic
polynomial in the newly used first variable is a non-zero-divisor.
If \(F_d\ne0\) is the lowest nonzero homogeneous component of an
element \(F\) in the old ideal, then
\[
 (u^pF)_{d+p}=U_1^p F_d\ne0.
\]

All smaller degrees are zero. The degree test for the coordinate
relations above then shows that any old detection through
degree \(w\) gives some new generator detected through
\(w+p\). Apply this modulo \(I\) to the determinant, and
modulo \(\pi\) to the adjugate with unused first variables
zero. In the latter test the old cofactor survives when the
newly used first variable is restored, since it survived even
with that variable zero.

\end{proof}
\paragraph{Producing a dual class.}
Process one remaining signature, with current bound \(w\).
Put \(D=(e+1)w\). If its old determinant is already detected
through degree \(2D+1\), take a new pair with trivial prescribed
Frobenius modulo the common subgroup and use the preceding
comparison. Otherwise choose an old square presentation with
matrix \(\Phi\) of positive size \(a\) and a polynomial lift
of its determinant with all coefficients through degree
\(2D+1\) in \(I\).

Choose a cofactor, written as an adjugate
entry \(H\), with a unit coefficient through degree \(w\)
after the unused first variables are set to zero. Choose a nonzero homogeneous part \(\bar H_j\) of the residual cofactor in the used variables, with \(j\le w\). The slopes \(\bar a_i\) must satisfy
\[
 \bar a_i\ne0\quad (i\text{ used}),\qquad \bar H_j((\bar a_i)_i)\ne0.
\]
Such a vector exists on the torus over
the algebraic closure and has finitely many algebraic coordinates.
Lift the slopes to units \(a_i\) in a finite unramified DVR
extension \(O'/O\), putting \(a_i=0\) at unused variables.
Work in
\[
          A=(O'/I O')[\tau]/(\tau^{2D+2}),\qquad
          U_i=a_i \tau,\qquad c_0=\ell^\nu \tau^D .
\]
\begin{lemma}[A dual Selmer class in the truncated ring]
For the resulting specialization \(\Psi\) with zero determinant
and unit cofactor coefficient in degree at most \(w\), let
\(H=(\operatorname{adj}\Psi)_{jl}\) denote the selected cofactor.
Column \(l\) gives a primal Selmer class. On the dual side,
\[
 \exists\,y\in S^\perp(A):\qquad c_0y\ne0.
\]
\end{lemma}
\begin{proof}
The coordinate relations hold by the binomial estimate. The
specialized matrix \(\Psi\) has determinant zero. Suppose the
first unit coefficient of the specialized \(H\) occurs at
degree \(d\le w\). Then \(H=\pi B+\tau^d V\) with \(V\) a
unit. Modulo \(H\), the relation \(\tau^d=-\pi B/V\) gives
\(\tau^{(e+1)d}=0\). Thus \(H\mid\tau^D\), since \(D=(e+1)w\).
Because \(\Psi\,\operatorname{adj}\Psi=(\det\Psi)I_a=0\),
column \(l\) represents \(x\in S(A)\).
Writing \(\tau^D=H H'\), one has
\[
 (c_0\tau H)H'=\ell^\nu\tau^{2D+1}\ne0
 \quad\text{in }(O'/IO')[\tau]/(\tau^{2D+2}).
\]

We use the natural identification of kernel vectors with
classes for every coefficient module. Representatives can be
chosen from one cocycle with coefficients in the representing
module, by applying the functor isomorphism to its identity.
A chosen set of lifts of its entries to the free presentation
then gives matrices whose contraction with a kernel vector
represents the corresponding cocycle.

Let \(Y=\operatorname{Hom}(A,\mathbb Q_\ell/\mathbb Z_\ell)\)
be the additive dual with the balanced evaluation pairing.
Write \(S^\perp\) for the Selmer condition on the Tate dual
with inverse character, orthogonal to the signature condition
on \(Y\). The finite-module Selmer order formula gives
\begin{equation}\label{eq:pos-order}
 |\ker\Psi(Y)|=|S^\perp(A)|
             \prod_{v\in Q}|H^0(E_v,T^\vee(1)\otimes\kappa^{-1} A)|.
\end{equation}
Here \(\kappa\) is the specialized inverse universal
character, and the product includes both places of each pair.

To check the local factors for arbitrary finite coefficients,
use an integral basis of the saturated finite-cocycle module
at \(v\mid\ell\). Its residual functions are independent:
a dependency would divide a primitive integral combination
by \(\pi\) in the same lattice. Finitely many evaluations
give an invertible minor. The span on a finite module \(Y\)
thus has order \(|Y|^{r+d_v}\), where \(d_v\) is the de Rham
quotient dimension in Section~\ref{subsec:signature-selmer}.

Let \(Z_v(Y)\) be this span of cocycles, \(B^1_v(Y)\)
its coboundaries, and \(L_v(Y)=Z_v(Y)/B^1_v(Y)\).
The coboundary map has source of order \(|Y|^r\) and
kernel \(H^0_v\), so
\[
 |L_v(Y)|=\frac{|Z_v(Y)|}{|B^1_v(Y)|}
          =\frac{|Y|^{r+d_v}}{|Y|^r/|H^0_v|}
          =|Y|^{d_v}|H^0_v|.
\]
The product of the factors \(|Y|^{d_v}\) cancels the
complex-place factor.
At the full local conditions away from \(\ell\) the local
Euler characteristic leaves the dual invariants. They vanish
at the exceptional places by the inertia test. Outside the
ramification set unramified conditions are paired. Both global
invariant groups vanish by the scalar marking. This proves
\eqref{eq:pos-order}, and the same proof applies to the
perfectly paired modules \(Y/c_0\tau Y\) and \(A[c_0\tau]\).
\cite{milne-duality}

There exists \(y\in S^\perp(A)\) with \(c_0[y]\ne0\).
Indeed at a first pair the dual invariants vanish by the
residual Frobenius test and finite dévissage. At each current
second pair the inertia of its own coordinate has value
\(1+a_i \tau\) with \(a_i\) a unit. For any invariant \(v\) at such a place,
\[
 (1+a_i\tau)v=v\quad\Longrightarrow\quad \tau v=0 .
\] Suppose all the dual classes were killed
by \(c_0\). For a representative \(y\), write
\[
 c_0y=dm.
\]

The image of \(m\) in the coefficient quotient by \(c_0\)
is invariant globally, hence zero. Subtracting a coboundary
from \(y\) gives values in \(A[c_0]\). At a strict place a
local expression \(y=dm_v\) then has \(c_0m_v\) invariant,
so \(c_0\tau m_v=0\). At \(\ell\) every finite local class
on \(Y/c_0\tau Y\) lifts from \(Y\) by the local description;
naturality of the pairing gives the required orthogonality.
The class thus lifts to the dual Selmer group on \(A[c_0\tau]\).
The invariant products in \eqref{eq:pos-order} coincide,
so
\[
 |\ker\Psi(Y)|\le |\ker\Psi(Y/c_0\tau Y)|.
\]

For an endomorphism \(U\) of a finite module \(M\),
\[
 |\ker U|=|M|/|U(M)|=|\operatorname{coker}U|.
\] Accordingly, the kernel inequality implies
that the quotient surjection
\[
 \frac{Y^a}{\Psi Y^a}
 \longrightarrow\frac{Y^a}{\Psi Y^a+c_0\tau Y^a}
\]
is bijective. Hence \(c_0\tau Y^a\subseteq\Psi Y^a\). Applying
row \(j\) of \(\operatorname{adj}\Psi\), which annihilates this
image, to \(c_0\tau Y\) in coordinate \(l\) shows
\((c_0\tau H)Y=0\). The evaluation pairing of \(A\) with \(Y\)
is perfect, implying \(c_0\tau H=0\), a contradiction.

\end{proof}
\paragraph{A Frobenius element with the prescribed evaluations.}
Intersect \(J\) with the kernels of the used coordinates to
obtain \(J'\). This subgroup is normal and conjugation-stable.
The two twisted actions are trivial on it, the commutators
\([g,z]\) lie in it, and \(\kappa(z)=1\) by the ramification
condition. Both restrictions of the cocycles to \(J'\) are
therefore additive, independent of representatives and
equivariant under conjugation.

\begin{lemma}[A second-half Frobenius with nonzero pairing]
For the primal and dual classes just chosen, there is a new prime
pair, avoiding the previous exceptional set, with the prescribed first-coordinate values, satisfying at the chosen prime \(P\), with Frobenius \(\sigma\),
\[
 \mathrm NP\equiv1\pmod{2\ell^M},\qquad B_\sigma\ne0.
\] Its untwisted Frobenius and that
of its conjugate act trivially modulo \(I\).
\end{lemma}
\begin{proof}
Some \(h\in J'\) has
\[
 B_h=x(h)\cdot y(h)-x(h^c)\cdot y(h^c)\ne0\quad\text{in } A .
\]
For suppose all these differences vanished. Polarizing at
\(h,k\in J'\) gives equality of the two cross sums. Replacing
\(h\) by \(zhz^{-1}\) multiplies its \(x\)-value by \(d\)
and its \(y\)-value by \(d^{-1}\), with the inverse multipliers
on the conjugate pair. Since \(d-d^{-1}\) is a unit, subtraction
separates the two terms and yields
\[
            x(h)\cdot y(k)=x(k^c)\cdot y(h^c).
\]

Let \(R\) and \(R'\) be the twisted primal and dual actions.
Conjugate duality and anti-invariance of \(\kappa\) identify
\((R')^c\) with \(\chi R\), \(\chi=\epsilon^{\text{tw}}\) with \(\text{tw}\) the
signature's fixed twist exponent.
Replace \(h\) by \(ghg^{-1}\) in the last identity.
Equivariance on \(J'\) gives, for every \(g\),
\[
 \operatorname{tr}(R(g)U)=\operatorname{tr}(\chi(g)R(g)V),
 \qquad U=x(h)y(k)^{\mathsf T},
\]
where \(V\) is the conjugate of \(y(h^c)x(k^c)^{\mathsf T}\)
by the fixed intertwiner. The span of \(R(g)\) is the full
matrix algebra over \(A\), by residual absolute irreducibility
and Nakayama. The same holds for \(\chi R\).

Consider the subalgebra spanned by the pairs \((R(g),\chi(g)R(g))\) in
\(M_r(A)^2\). It surjects onto both factors. Its intersections
with the separate factors are two-sided ideals, of the form
\(M_r(I_1)\) and \(M_r(I_2)\) by multiplication with matrix
units. If \(\mathcal A\) is this subalgebra, a pair \((B,C)\in\mathcal A\)
associates the class of \(B\) modulo \(M_r(I_1)\) to the class
of \(C\) modulo \(M_r(I_2)\). Surjectivity and the definitions
of \(I_1,I_2\) give a well-defined inverse in the other direction.
It is an isomorphism
\[
      M_r(A/I_1)\xrightarrow{\sim} M_r(A/I_2)
\]
of \(A\)-algebras, so their scalar annihilators coincide and
\(I_1=I_2=I'\).

If the quotient is nonzero its matrix algebra
automorphism is inner. Indeed the images of the diagonal matrix
units cut the standard module over the local ring \(A/I'\)
into isomorphic projective summands, necessarily free of rank
one. Transport a generator of the first by the images of the
off-diagonal matrix units; the resulting basis conjugates
all matrix units at once. Taking determinants in the quotient
gives \(\chi(g)^r-1\in I'\) for every \(g\), also true for the
zero quotient. By definition of \(\nu\), these scalars generate
\((\ell^\nu)\) after coefficient extension. The trace identity
on \((\ell^\nu E_{ab},0)\) gives
\[
 \ell^\nu x(h)_i y(k)_j=0
\]
for all coordinates.

Let \(J_x\) be the ideal generated by the entries of all
\(x(h)\), \(h\in J'\). Modulo \(J_x\) the cocycle factors
through \(G_E/J'\), on which \(z\) is central and acts by
a scalar differing from one by a unit. Comparing the cocycle law at \(zg\) and \(gz\) gives
\[
        (d-1)x(g)=(g-1)x(z)
\]
over \(A/J_x\). As \(d-1\) is a unit, this expresses the
cocycle as a coboundary with constant \((d-1)^{-1}x(z)\). Naturality and injectivity
of the kernel description put every entry of column \(l\)
of the adjugate in \(J_x\).

In particular \(\tau^D\in J_x\)
since \(H\mid\tau^D\). As \(\ell^\nu J_x y(k)=0\) for \(k\in J'\),
this gives
\[
 c_0y|_{J'}=0.
\]
In the quotient
\(G_E/J'\), the central element \(z\) acts on the dual by \(d^{-1}\)
with \(d^{-1}-1\) invertible. For \(f=c_0y\) the two evaluations of \(f(zg)=f(gz)\) give
\[
 (d^{-1}-1)f(g)=(g-1)f(z).
\]
Dividing by the unit exhibits \(f\) as a coboundary, contradicting the preceding lemma.

Choose an open normal subgroup in \(J'\) on which both
cocycles and their conjugates vanish. This is possible since
their restrictions are continuous homomorphisms to finite
modules; intersect the kernels and take the normal core.
This subgroup can contain all inertia at the unused first
pairs. There the used characters and signature actions are
unramified and the residual Frobenius test gives vanishing
of local \(H^1\) for both modules by finite dévissage.
Local coboundaries vanish on inertia, also at every conjugate
place and for every conjugate of its inertia group.
By Chebotarev choose a new second pair with Frobenius
\(\sigma\) conjugate to \(h\) modulo this subgroup, prescribing the first-coordinate vector
\[
 (0,\ldots,0,1,0,\ldots,0)
\]
with the single \(1\) in the next unused first coordinate (additive notation). Indeed \(h\) is in
the kernels of the used coordinates and the unused first
inertia groups can correct all the remaining coordinates.

All old used characters then vanish on the new decomposition
group. The Frobenius can be adjusted by its own inertia to
give value zero also in the new coordinate. Equivariance and
the dual pairing multiply \(B_h\) under conjugation by a
cyclotomic unit, equal on the two conjugate terms. Thus
\(B_\sigma\ne0\). The new prime avoids all exceptional
places, has norm one modulo \(2\ell^M\), and its untwisted
action and that at its conjugate are trivial modulo \(I\).
At each stage the open subgroup and avoidance set may include
the common earlier Frobenius requirements of the other signatures.

\end{proof}
\begin{proposition}[A determinant coefficient after all positive-valuation steps]\label{prop:positive-signature-determinant}
Starting from Lemma~\ref{lem:residual-selmer-coefficient},
under the hypotheses in this construction, after choosing the
reserve pairs, there is a common degree bound \(D\) such that every signature determinant \(F_s\) has a coefficient satisfying
\[
 \exists\alpha:\quad |\alpha|\le D,\qquad v_\ell([X^\alpha]F_s)\le\nu_\ell.
\] The number of pairs, degree and depth threshold depend
only on the original finite data. For each sufficiently large
\(M\), one choice of pairs serves all signatures. If
\(\nu_\ell=0\), no reserve pair is needed.
\end{proposition}

\begin{proof}
\proofstep{Make the new coordinate substitutions before taking the Fitting ideal.}
Use the finite ring
\[
 B=A[X,Y]/(X^2,Y^p),
\]
over \(A\) by
constant inclusion and with projection \(B\to A\).
Keep the old used variables \(a_i \tau\), set the reserve first
variable corresponding to this pair to \(Y\), and set its
own second variable to \(X\). These assignments satisfy all
coordinate relations by the binomial bound. At the new
prime take a tame generator \(\gamma\) mapped to one in the
new coordinate. The inverse universal character has values
\((1+X)^{-1}\) and \((1+Y)^{-1}\) on \(\gamma,\sigma\);
the conjugate pair has inverse values. On projection to \(A\)
all these values become one.

Let \(P'\) be the new representing module after this base
change. For its universal cocycle write \(u_i,t_i\in P'\)
for the Frobenius and inertia entries at the two places,
\(1\le i\le p\). The tame relation gives
\[
          X u_i=Y\eta_i t_i,\qquad
             \eta_i\in B^\times,\quad \eta_i\bmod(X,Y)=1 .
\]
To verify it write \(q\) for the norm. On the pro-\(\ell\)
tame quotient
\[
 \sigma\gamma=\gamma^q\sigma.
\]

For an inertia scalar \(1+v\), \(v^2=0\), the cocycle on
\(\gamma^q\) is \((q+\binom q2 v)\) times its value on
\(\gamma\), hence just its value on \(\gamma\), since
\(q\equiv1\bmod 2\ell^M\). Wild and prime-to-\(\ell\) tame
inertia in the kernel of the action have zero cocycle values.
Expanding the relation now gives the displayed identity,
with
\[
 \eta_i=(1+Y)^{-1}
\]
or one.

In \(P'/(X,Y)P'\), quotienting by the inertia entries gives
the old module with presentation \(\Psi\), by the unramified
condition. Choose lifts of its \(a\) generators. The projected
Frobenius rows modulo inertia are the old rows \(U_i\)
evaluated at \(\sigma,\sigma^c\); an ambiguity of the
representative is a coboundary and vanishes at these elements
of trivial action. Define \(k_i\) to be the entries of
\(y(\sigma)\) on the first side and \(-y(\sigma^c)\) on the
second. Global reciprocity gives a further relation
\[
                \sum k_i t_i=0 \pmod{(X,Y)P'} .
\]
Indeed test against the additive dual of \(A\). A new relaxed
class on that dual with the old character pairs with the old
dual class \(y\). The local cup invariants vanish at the old
places by their conditions, and outside the level by
unramifiedness.

At each of the two new places the action is
trivial and \(y\) is unramified. Only the new inertia value pairs with \(y\)'s Frobenius value. To determine the two signs, identify the scalar Tate coefficients by a primitive \(\ell^{\nu+1}\)-st root. By local reciprocity the unramified character taking Frobenius to one pairs to a generator with the Kummer class of a uniformizer, whose tame pro-\(\ell\) inertia value is a unit. Unramified parts pair to zero.
Conjugation takes the chosen primitive root to its inverse. Thus, in the additive identification of the scalar Tate coefficients,
\[
 \zeta^a\longmapsto\zeta^{-a},\qquad a\longmapsto-a.
\]
This gives the opposite sign at the conjugate place. Applying the
sum of invariants and varying the homomorphism to the additive
dual proves the relation, since that dual is a cogenerator
for finite \(A\)-modules. \cite{milne-duality,milne-cft}

The lifted old generators and \(\eta_i t_i\) generate \(P'\)
by nilpotence of \((X,Y)\). Use them in this order. In the
reduction to \(A\), write the Frobenius rows as \((U_i,*)\),
and lift to rows over \(B\). The top \(a\) relations may be
chosen to reduce to \((\Psi,C)\) for some \(C\), except that
row \(l\) is replaced by \((0,k)\) by reciprocity. Each
reduced relation lifts: its sum lies in \((X,Y)P'\), so
subtract an expression in the same generators with
coefficients in \((X,Y)\).

The additional \(p\) exact tame rows are
\[
 Y(0,E_i)-X(\widetilde U_i,*)\qquad(1\le i\le p),
\]
where \(E_i\) is the \(i\)-th coordinate row. The determinant of this square
array lies in \(\operatorname{Fitt}_0(P')\). Expand in the
last \(p\) rows. The terms with no \(X\) or with at least two copies vanish, since \(Y^p=X^2=0\). For the surviving terms
\[
 X\,(XY^{p-1})=0,\qquad Y\,(XY^{p-1})=0,
\]
so all their other coefficients
may be reduced to \(A\). Expanding now on row \(l\) and the
\(p-1\) unchanged coordinate rows gives, up to one common sign,
\[
       XY^{p-1}\sum_i k_i(U_i\operatorname{adj}\Psi)_l
                      =XY^{p-1}B_\sigma\ne0.
\]
The last inequality uses the monomial basis over \(A\).

\proofstep{Use the bounded nilpotence degree to locate a coefficient.}
Consequently the specialized new determinant, which generates
this Fitting ideal, is nonzero. A monomial of total degree
greater than \(2D+1+p\) vanishes on this substitution:
the old variables are multiples of \(\tau\), and the two others
have exponents at most one and \(p-1\). Thus every lift of
the determinant has a coefficient outside \(I\) in the
asserted degree range. This reasoning applies to any new
square presentation for the same functor.

\proofstep{Use the common degree bound while preserving earlier signature detections.}
The final degree bound must retain the coefficients already detected at earlier signatures.

Repeat this construction for the \(b\) signatures in turn.
The Fitting comparison at each new pair preserves all earlier
detections and the available cofactors within the next bound
\(w_{j+1}\). For each signature determinant \(F_s\), the final common degree bound reads
\[
 \min_{|\alpha|\le w_b} v_\ell([X^\alpha]F_s)\le\nu_\ell.
\]
Here the multi-index ranges over all the variables. For each sufficiently large
\(M\), the same choice of pairs serves all signatures, and the
pairs may depend on \(M\). For \(\nu_\ell=0\)
the residual argument finishes without the reserve pairs.
The number of pairs, the degree and the depth threshold may
depend on the original finite data at the coefficient place;
they are fixed before increasing \(M\). The subsequent
specialization of Section~\ref{subsec:retained-twist} retains
the valuation allowance \(\nu_\ell+\delta\).
\end{proof}


\subsubsection*{The abelian theta coefficients}
\label{subsec:toric-twist}
\hypertarget{source-toric-twist}{}

The Selmer determinant controls the signature contribution to the simultaneous twist. The abelian theta values require a separate argument: changing the toric order changes both the coefficients and their period normalization. Let \(E/F\) be the CM extension of the
selected host, with \(F\) totally real, and let \(\varphi\)
be its chosen CM type. Fix a toric label, its integer
\(s\ge1\), the nonnegative harmonic exponents \(\kappa\),
its base conductor \(\mathfrak f\), and its unitary Hecke
character \(\chi\), trivial on the ideles of \(F\).
Its infinity exponents are \(\kappa_i+s\) at the places
of \(\varphi\) in the unitary normalization. At the base
conductor its local conductor is prescribed and nontrivial.
Enlarge the order conductor to
\[
 \mathfrak c=\mathfrak f\prod_{v\in Q}v
\]
in the real
field, at split places prime to the base conductor and to
the coefficient prime \(\ell\).
Let \(\operatorname{Pic}(\mathfrak c)\) denote the
ideal class group of the resulting order.

For a finite idele \(y\) representing a class, let
\(S_y(z)\) be the continuation of the lattice sum on
\(E\cap y\widehat{\mathcal O}_{\mathfrak c}\), with the
harmonic numerator and shifted exponent of the fixed label,
over nonzero vectors modulo \(\mathcal O_F^\times\).
The toric term is
\[
       T_y=\frac{S_y(0)}
          {\chi_f(y)\,|y|_f^{s}},
\]
where the finite modulus \(|y|_f\) is the one in the
unitary Hecke normalization. In the convergent region the
weights are homogeneous under
multiplication of the lattice by an element of \(E^\times\).
The character at infinity and the norm exponent give
the inverse of this change in the finite prefactor.
The completed order units preserve the lattice and
the finite character is trivial on them. The resulting
normalized term therefore depends on the class. These
identities pass to the continued value.

\paragraph{The denominator after changing the order.}
Choose the fixed CM period lattice with basis
\(1,(1+\sqrt{-d})/2\) at the imaginary quadratic subfield,
and choose a twelfth root of its discriminant. Incorporate the
fixed power of this root and the normalization constant
of the harmonic numerator into a nonzero scalar
\(\Omega_*\). In this normalization
\[
 \beta(T)=(T/\Omega_*)^{12}.
\]

The local class-value estimate for lattices with
\(y_w=1\) at every \(w\mid\ell\) has the following form.
Put \(W=\operatorname{wtM}(2s,\kappa)\), the degree of
the period used in the twelfth power. Let \(a_0\) be the
maximum conductor exponent of \(\mathfrak f\) over \(\ell\).
Let \(c_E\) be the largest denominator exponent at
\(\ell\) of the partial zeta values of \(F\) at \(1-2s\),
divided by \(2^{[F:\mathbb Q]}\) before taking the
denominator. Put \(f_0=|I|=[F:\mathbb Q]\) for the size
of the chosen CM type, \(K_\kappa=\sum_{i\in I}\kappa_i\),
and \(W=2sf_0+2K_\kappa\).
Write \(c_\ell=a_0\) if \(a_0>0\) and
\(c_\ell=v_\ell(D_E)/[E:\mathbb Q]\) otherwise. The
ramified allowance is
\[
 n_\ell =
 \begin{cases}
 0,&c_\ell=0,\\
 2sf_0c_\ell+K_\kappa(2c_\ell+1+
       \mathbf 1_{\ell=2}+\tfrac12\mathbf 1_{\ell=d}),&c_\ell>0 .
 \end{cases}
\]
\begin{proposition}[Uniform denominator of the CM class values]\label{prop:toric-class-denominator}
At every embedding into \(\overline{\Q}_\ell\), for a representative
with \(y_w=1\) at all \(w\mid\ell\), the normalized twelfth
power is algebraic and satisfies
\[
    v_\ell\bigl(\beta(S_y(0))\bigr)\ge-12(c_E+n_\ell+Wa_0).
\]
Auxiliary split places disjoint from \(\ell\) may vary.
\end{proposition}
The proof follows in the local jet calculation below. First we pass
from one lattice to the coefficients of the theta element.

\begin{lemma}[Clearing the abelian theta coefficients]\label{lem:toric-common-denominator}
With the CM type, base conductor, critical weight and coefficient
prime fixed, there is an integer \(c\ge0\) such that every coefficient \(a\) of a normalized abelian theta element, at any allowed auxiliary level and after any allowed finite-order ring class twist, satisfies
\[
 v_\ell(a)\ge-c.
\]
\end{lemma}

\begin{proof}
Every ideal class admits a representative that is one
above \(\ell\). Indeed, use weak approximation in \(E\)
to put the components above \(\ell\) in the local
completed order-unit subgroups and then divide by those
units. On such a representative \(\chi_f(y)\) is a unit
at \(\ell\) in the algebraic normalization. This follows
from the formula for its infinity type on principal
ideles, applied after taking a power to clear the ideal
class, and from the roots of unity on the compact unit
group.

The element of \(E^\times\) used in this formula
has valuations zero above \(\ell\), so every conjugate
factor is a unit. The character and norm factors of such a representative have
\[
 v_\ell(\chi_f(y))=0,\qquad v_\ell(|y|_f)=0,
\]
where \(|y|_f\) denotes the finite norm factor in this normalization.
The same bound thus holds for \(\beta(T_y)\) on all
classes. If a finite order ring class character is added,
its values again have valuation zero. If the separately cleared summands \(a_y\) have nonnegative valuation, the only estimate required for their sum is
\[
 v_\ell\!\left(\sum_y a_y\right)\ge\min_y v_\ell(a_y)\ge0.
\]
Thus the number of classes of the varying order does not change the clearing power.

For the abelian theta module the coefficient at an Artin
label is obtained from the terms \(T_y\) in its fiber,
multiplied by the same nonzero normalization scalar
specified by the toric construction. The fiber map
uses the inverse ideal class. Choose the complex
embedding in the algebraic closure of the row fraction
field as in the toric comparison. Taking twelfth roots gives the individual integral bound
\[
 v_\ell(T_y/\Omega_*)\ge -N\qquad(N\in\mathbb Z,\ N\ge c_E+n_\ell+Wa_0).
\]

One more fixed power clears the product of
\(\Omega_*\) and the normalization scalar in that
algebraic closure. To see this for any algebraic element,
take a monic polynomial over the fraction field and
choose the power large enough that the rescaled monic
polynomial has integral coefficients. Each coefficient
of the theta module is already in the fraction field
of \(O\). Integral closure of the DVR puts its cleared
value in \(O\). For each abelian label with an allowed
\(s\), use one such \(s\) for its fixed coefficient
functions and then take the maximum over the finitely
many labels of this row. This gives the exponent \(N_0\)
in the abelian part of the simultaneous specialization.
\label{cor:toric-twist-1}
\end{proof}

\paragraph{The local jet bound in period coordinates.}
Let \(\omega^{12}\) be the fixed lattice discriminant
and let \(g\) be the defining holomorphic lattice sum
before inserting the powers of the conjugate variables.
The field data here use \(E/\mathbb Q\) abelian,
\(E/F\) unramified at all finite places, and a CM type
containing the chosen image \(i\sqrt d\) of \(\sqrt{-d}\)
at each selected embedding. In particular
\([F:\mathbb Q]\ge2\). If \(F=\mathbb Q\), the
discriminant formula would give \(|D_E|=1\) for an
imaginary quadratic field. The condition on the real
degree applies at weight two as well.

\begin{proof}[Proof of Proposition~\ref{prop:toric-class-denominator}]
\proofstep{Integral idempotents and tangent directions.}
We give the integral comparisons on the lattice
\(E\cap y\widehat{\mathcal O}_{\mathfrak c}\) with
\(y=1\) above \(\ell\). Work over the valuation ring
\(R\) of a finite algebraic coefficient field in
the prescribed embeddings. We enlarge it finitely
to contain the CM embeddings and the good-reduction
model constructed below. The order action extends
to that abelian scheme by the extension property
for homomorphisms. Use \(v_\ell(\ell)=1\) in every
coefficient extension.
Use a generator of the fractional polarization
module locally at \(\ell\). If \(H\) denotes the integral de Rham lattice, perfection of the resulting alternating form means
\[
 H\xrightarrow{\;v\mapsto B(v,-)\;}H^\vee\ \text{is an isomorphism}.
\]

Indeed that local generator gives an isomorphism to
the dual at the coefficient prime, and the integral
de Rham duality comparison gives its perfect form.
The comparison with the ordinary polarization
pairing uses \(2B\), with \(B\) itself the perfect
integral pairing. The generic CM eigenspaces have
rank one, with Hodge subspace given by one
embedding from each conjugate pair. The perfect
pairing is between the opposite eigenspaces.

Let \(D\) be a scalar clearing every primitive CM
idempotent, with
\[
 v_\ell(D)=c_\ell.
\]
For a basis
\((x_j)_j\) of the maximal order and its trace-dual basis
\((x_j^\vee)_j\), the idempotent at an embedding \(\iota_i\)
is \(\sum_j\iota_i(x_j^\vee)x_j\) after coefficient
extension. Multiplication by an element of the different
clears the trace-dual coordinates. Galois uniformity gives
\[
 c_\ell=v_\ell(D_E)/[E:\mathbb Q]
\]
away from the conductor.
At a conductor prime the maximal order is unramified at the
odd prime \(\ell\), so those idempotents are integral there.
Multiplication by \(\ell^{a_0}\) puts their coefficients in
the conductor order. The extra order places are prime to
\(\ell\), so the same local order works in both comparisons.

Let \(e_i,f_i\) be the primitive generators of the two
opposite integral lines in the de Rham lattice and put
\(P_i=B(e_i,f_i)\). They are saturated line lattices.
Perfection of \(B\) supplies an integral \(x\) with
\(B(e_i,x)=1\). Write \(D E_{\bar i}x=s_i f_i\), where
\(E_{\bar i}\) is the idempotent on the opposite line and
\(s_i\in R\). Then
\[
      s_iP_i=B(e_i,D E_{\bar i}x)=D ,
      \qquad 0\le v_\ell(P_i)\le c_\ell .
\]
The equality uses the oppositely paired eigenspaces.
Multiplying the rank-one map \(e_i\mapsto f_i\) by \(D\)
gives an integral operator \(M_i\) killing other eigenlines,
as \(D E_i\) maps into the primitive lattice on the
line of \(e_i\), where \(E_i\) denotes that idempotent.
Each \(M_i\) commutes with the real order and is skew for \(B\).

If \(M=\sum_i t_i M_i\), then \(M^2=0\) and
\[
 B((1+M)u,(1+M)v)
  =B(u,v)+B(Mu,v)+B(u,Mv)+B(Mu,Mv)=B(u,v).
\]
The linear terms cancel by skewness and the last is
\(-B(u,M^2v)=0\). The same expansion applies to the other
polarization forms, since they are obtained by real
multiplication and the \(M_i\) commute with that multiplication.

\proofstep{Comparison with the reference CM periods.}
The primitive integral tangent eigenforms have
complex period scalars \(a_i\). For their comparison
with \(\omega\), let \(A_{\rm ref}\) be the elliptic curve
of the fixed quadratic lattice. Its CM \(j\)-invariant
is integral, so it has good reduction after finite
extension at \(\ell\). Choose an integral generator of
its tangent line. Its analytic coordinate \(t\) in the
given lattice has \(t^{12}\omega^{12}\) equal to the
discriminant of a good Weierstrass model. This is the
weight-twelve change of Weierstrass invariants, using
the full Weierstrass equation at two and three.
Thus \((t\omega)^{-1}\) is algebraic and integral at
the selected place.

Let
\[
 L_y=E\cap y\widehat{\mathcal O}_{\mathfrak c}.
\]

Since the local lattice at each place above \(\ell\)
is the conductor order, \(\ell^{a_0}\mathcal O_E\)
is contained in it locally. At the remaining places
one integer \(m_0\) prime to \(\ell\) clears the
necessary lattice index. Thus
\[
 n_0\mathcal O_E\subset L_y
\]
with
\(n_0=\ell^{a_0}m_0\). Choose \(\theta\in\mathcal O_E\)
separating the complex embeddings and take the
matrix with columns \(n_0(\varphi_i(\theta^j))_i\),
\(0\le j<f_0\).

For \(V=(n_0\varphi_i(\theta^j))_{i,j}\), the embeddings
give distinct values \(\varphi_i(\theta)\), so
\[
 \det V=n_0^{f_0}\prod_{i<i'}
               (\varphi_{i'}(\theta)-\varphi_i(\theta))\ne0 .
\]
The reference lattice consists of simultaneous embedded
quadratic integers. Thus column \(j\) maps such an
integer \(a\) to the lattice element \(n_0a\theta^j\)
of \(n_0\mathcal O_E\subseteq L_y\). The matrix defines an
isogeny from \(A_{\rm ref}^{f_0}\) to the CM variety.
After finite extension its finite generic kernel has
a finite flat schematic closure in the torsion group
of the good source. Quotient by this closure gives
the good target model and an integral isogeny.
The integral tangent vector on the factor in column
\(j=0\) has image with analytic coordinates
\(w_i^0=n_0t\) for all \(i\). Consequently
\[
 v_\ell((w_i^0\omega)^{-1})\ge-a_0.
\]

Evaluating the integral eigenform on this vector
gives \(a_iw_i^0\in R\). This implies the first
of the two bounds below. For the determinant,
the same cleared CM idempotents on the tangent
dual put \(D\) times the full lattice in the span
of the primitive eigenforms. Their wedge is
\(u_0\) times a determinant generator, where
\(0\le v_\ell(u_0)\le f_0c_\ell\) by Smith normal
form. The complex determinant scalar is
\[
 C=(\prod_i a_i)/u_0.
\]


For the lowering coordinates let \(w_i\) be the
complex tangent coordinate of \(f_i\) and choose
the integral unit \(r_i\) with \(P_i=r_i a_i w_i\).
This equality follows by using \(B(e_i,-)\) as
the primitive covector on the tangent quotient.
The Riemann form in the polarized uniformization
has real scale \(h_i\) on this coordinate.
If \(\psi\) is the complex polarization pairing
and \(\xi\) is the imaginary generator of the
polarization module locally at \(\ell\), the
comparison of alternating forms gives
\[
  2i\varphi_i(\xi)\psi(e_i,f_i)=h_i P_i .
\]
The local ideal of polarizations is dual to the
lattice under the trace pairing.

At the chosen
places, \(y=1\) and the extra level is a unit.
The local generator consequently satisfies
\(v_\ell(\varphi_i(\xi))\le0\).
Equivalently one can choose a local anti-invariant
inverse-different generator \(\xi'\) of nonpositive
valuation and use
\(m\xi'=a\xi\) with \(m\) prime to \(\ell\)
and \(a\in\mathcal O_F\) to compare the two generators.
The scale of the lowering direction and its auxiliary
period are
\[
  A_i=\pi h_i\frac{D}{\psi(e_i,f_i)}w_i^2,\qquad
  b_i=\frac{a_i A_i}{D}.
\]
In particular \(b_i/(a_i A_i)=D^{-1}\). Writing
\(\delta=\sqrt{-d}\) with \(\varphi_i(\delta)=i\sqrt d\),
the other period quotient is
\[
   (a_i b_i\sqrt d/(2\pi))^{-1}
       =\frac{r_i^2}{P_i\,\varphi_i(\xi\delta)} .
\]
All of these right-hand sides are algebraic.
Since \(v_\ell(\varphi_i(\delta))=\frac12\mathbf1_{\ell=d}\),
we obtain
\[
 v_\ell(a_i/\omega)\ge-a_0,\qquad
 v_\ell\bigl((a_i b_i\sqrt d/(2\pi))^{-1}\bigr)
     \ge-c_\ell-\tfrac12\mathbf1_{\ell=d}.
\]
The determinant and inverse-idempotent bounds just
proved supply the other two valuations.

\proofstep{The divided-power deformation.}
Use the divided-power thickening of the chosen abelian scheme. In horizontal coordinates the prescribed variation of its Hodge submodule is
\[
 \operatorname{Fil}^1_{\boldsymbol t}=(1+\sum_i t_iM_i)\operatorname{Fil}^1_{0},
\]
with the original filtration extended to the jet ring.
The real action and polarizations lift with the
filtration because the indicated operators commute
with that action and preserve the forms.
Here use the divided-power jet lifting theorem for the integral level.
Its base is a characteristic-zero local domain and its
thickening is the truncated divided-power jet ring,
with split augmentation and nilpotent divided-power
kernel. The input endomorphisms are admissible for the
real action, filtration and alternating forms, as are the
lowering operators above. The theorem gives a polarized
abelian scheme with real multiplication over that jet
ring, together with the pullback to the augmentation.
Under the horizontal de Rham identification, its
Hodge submodule is the span of the prescribed jet lifts
of the original filtration.

If \(s_{\det}\) is the determinant lift and \(s_0\) the given generator at the origin, its normalization in the horizontal identification is
\[
 s_{\det}\bmod(t_1,\ldots,t_d)=s_0,
\]
where \(d\) is the number of deformation parameters.
Complex base change identifies the integral jet and its lifted filtration with the analytic deformation used to calculate the periods.
On the reductions modulo powers of \(\ell\) the
underlying deformation comparison has \(\ell\) nilpotent,
as in Grothendieck--Messing theory.
\cite{messing}\label{cor:toric-twist-2}

Transport the periods horizontally in that family.
In the \(i\)-th real factor take complex combinations
of the parallel cycles with holomorphic coordinate one
and with zero holomorphic coordinate and antiholomorphic
coordinate one, respectively. The latter cycle
represents the Hodge direction at the origin.
After the lowering it lies in the annihilator of
the lifted holomorphic forms.

Let \(P_i^{\rm hol}\) be the nonzero period on the first cycle and \(\varepsilon\) the period pairing sign. Evaluation on the second cycle gives
\[
          P_i^{\rm hol}(\zeta_i-\varepsilon A_i t_i)=0 .
\]
Its constant term is nonzero, so \(P_i^{\rm hol}\) is a unit
in the complex jet ring. Hence \(\zeta_i=\varepsilon A_i t_i\) there.
The determinant scale on the horizontal frame is
\(C\). Indeed the parallel cycles with holomorphic
coordinates the standard basis are on the opposite
CM eigenlines, where the lowering operator is zero.
Evaluating the lifted determinant on their wedge
gives exactly the value at the origin.

\proofstep{Extracting the coefficient of the prescribed jet.}
Put \(g_\kappa=\partial^\kappa g(0)\) and
\[
 J_\kappa=g_\kappa\prod_i((b_i/a_i)^{\kappa_i}/a_i^{2s}).
\]

Let \(m=\ell^{c_E}\) and let \(R\) be the local ring
of algebraic coefficients for the integral CM model,
contained in the chosen valuation ring. It contains
the rational \(\ell\)-integers. The weight-\(2s\) Hilbert
Eisenstein form on each polarized component has
integral nonconstant Fourier coefficients in the
algebraic normalization. Its constant coefficient
is the corresponding partial zeta value divided
by \(2^{f_0}\). The partial zeta rationality theorem
applies since \(F\) is totally real and \(s\ge1\).
The integral model of the Hilbert moduli problem
and the \(q\)-expansion principle on all components
then gives the form multiplied by \(m\) over \(R\).
For weight two the Fourier statement uses the value
at zero of the continued family in real degree at
least two. \cite{deligne-pappas,serre-tate}\label{cor:toric-twist-3}

Use the divided-power algebra in the deformation parameters truncated at total degree
\[
 |\kappa|=\sum_i\kappa_i.
\]

The admitted de Rham operators of the integral CM
model determine its deformation over this algebra.
The lifted determinant generator reads the section
as a scalar in the divided-power basis, with
coefficient \(x\in R\) at the multi-index \(\kappa\).
After base change to \(\mathbb C\), the lowering
direction in the \(i\)-th coordinate has scale
\(\varepsilon A_i\), where \(\varepsilon\in\{1,-1\}\).
The determinant scale is \(C\). Comparing the divided
power coefficient after base change gives
\[
       C^{2s}x=m g_\kappa\prod_i(\varepsilon A_i)^{\kappa_i}.
\]

In particular there is no factorial in this coefficient
identity. The integral model comparison gives
\(v_\ell(C/\prod_i a_i)\ge-f_0c_\ell\) and
\(v_\ell(b_i/(a_i A_i))\ge-c_\ell\). Thus
\[
       J_\kappa=x m^{-1}
         \left(\frac{C}{\prod_i a_i}\right)^{2s}
         \prod_i\left(\frac{b_i}{\varepsilon a_i A_i}\right)^{\kappa_i}
\]
has valuation at least \(-c_E-2sf_0c_\ell-K_\kappa c_\ell\).
At zero depth this is \(-c_E\).
At positive depth the sum of the factorial valuations satisfies
\[
                  \sum_i v_\ell(\kappa_i!)\le K_\kappa/(\ell-1).
\]
Subtract this from the bound for \(v_\ell(J_\kappa)\) to obtain
\[
 v_\ell\!\left(J_\kappa/\prod_i\kappa_i!\right)
    \ge -c_E-2sf_0c_\ell
             -K_\kappa(c_\ell+1+\mathbf1_{\ell=2}) .
\]
\proofstep{Returning to the harmonic lattice sum.}
The normalized coordinate derivative for the harmonic
lattice sum is \((-1)^{K_\kappa}c_{2s}
\prod_i(2s)_{\kappa_i} S_y(0)\), where \(c_{2s}\) is
the defining lattice constant and \((2s)_a\) is the
rising factorial. This identity follows by differentiating
the continued defining sum. In the convergence region each
derivative increases the relevant denominator exponent by
one, supplying its successive factor, and continuity in
the exponent gives the formula at the evaluation point.
The definition of \(\Omega_*\) absorbs the same constant
and rising factorials. Multiplying \(J_\kappa\) by
\[
   \prod_i (a_i b_i\sqrt d/(2\pi))^{-\kappa_i}
                  (a_i/\omega)^{2s+2\kappa_i}
\]
leaves, up to the sign, \(S_y(0)/\Omega_*\).

In the positive-depth case use the divided-power jet
and also multiply by \(\prod_i\kappa_i!\). These
factorials are integral at every prime. The two period
factors cost \(K_\kappa(c_\ell+\tfrac12\mathbf1_{\ell=d})\)
and \(Wa_0\), respectively. At zero depth \(\ell\ne d\)
since the auxiliary field contains the imaginary quadratic
subfield of discriminant \(-d\). The same multiplication
then costs only \(Wa_0\). Adding these contributions
and taking the twelfth power proves the stated local
lattice estimate from the normalized jet inequalities.
The coordinate and jet comparisons use the fixed CM
type and base conductor. The auxiliary order and the
class label may vary with these data and the coefficient
prime fixed. \label{cor:toric-twist-4}
\end{proof}

\paragraph{Distribution at an added split place.}
Write \(q\) for the size of the residue field of a finite place \(v\) of the totally real field
that splits in the CM extension, and \(P,P^c\) for the places above it.
The local completed order at depth one consists of
pairs with equal residues. For a character unramified
at \(v\), let \(u\) be its value on a uniformizer at
\(P\), so its value at \(P^c\) is \(u^{-1}\). In the
absolutely convergent half-plane write \(X=q^{-z}\),
where \(z\) is the complex argument of the Hecke series.

\begin{lemma}[Distribution at a split finite place]\label{lem:toric-split-distribution}
With inverse Artin reciprocity on the ideal classes, removal of an
auxiliary pair multiplies a value unramified at the pair by
\[
       (1-uX)(1-u^{-1}X)+(q-1)X^2,\qquad X=q^{-z},
\]
where \(u\) is the corresponding twisted uniformizer value on
one factor. This relation holds on the remaining group-ring
coordinates with the same period normalization.
\end{lemma}

\begin{proof}
The two unit coordinates with equal residues give one
class modulo the order units. Pairs with both
coordinates divisible by a uniformizer give \(q-1\)
unit cosets for each pair of positive valuations.
There are no pairs with exactly one positive valuation.
The local sum is consequently
\[
     1+(q-1)\frac{X^2}{(1-uX)(1-u^{-1}X)} .
\]
Relative to the original two Euler factors this is the
polynomial
\[
          (1-uX)(1-u^{-1}X)+(q-1)X^2 .
\]
If the tested character is ramified at this place,
only the simultaneous unit class contributes. The other
terms sum to zero because the character on the ratio of
residue units is nontrivial. The usual local factor of
the Hecke series there is one as well. The exponent-one
order is essential to the unit coset count.

\proofstep{The unit group of the order.}
To justify applying this computation to the lattice sum,
we identify the global units of the order precisely.
The base conductor is nonzero and proper and has an odd
prime divisor of residual characteristic \(p\) unramified
in \(E\) over \(\mathbb Q\). Let \(\mathfrak P\) be a prime of
\(E\) above that divisor. If \(u\in\mathcal O_{\mathfrak c}^\times\),
then \(u,\bar u\) are algebraic units and
\[
 u-\bar u\in\mathfrak c\mathcal O_E\subseteq\mathfrak P,
\]
since the image of the order modulo the conductor is real.
All complex conjugates of \(u/\bar u\) have absolute value one.
An algebraic integer with that property is a root of unity,
as the monic polynomials of all its powers have bounded integral
coefficients. At \(\mathfrak P\) the root lies in the principal unit
group. Here the normalized local valuation satisfies \(v(p)=1\).

If \(v(t)\ge1\) and \(t\ne0\), then \(v(pt)=1+v(t)\). The
intermediate terms \(\binom pj t^j\), \(2\le j<p\), have valuation
at least \(1+2v(t)\), and the final term has valuation \(pv(t)>1+v(t)\)
because \(p\) is odd. Thus
\[
 v((1+t)^p-1)=v(pt)=1+v(t).
\] For an exponent \(a\) prime to \(p\) the linear term
gives \(v((1+t)^a-1)=v(t)\). For any positive exponent
\(m=p^j a\) with \((a,p)=1\), iteration therefore gives
\[
                  v((1+t)^m-1)=v(t)+j<\infty .
\]
This is incompatible with \(m\)-torsion.
Thus the principal unit group has no nonidentity torsion, and
\(u=\bar u\). Both \(u\) and its inverse belong to \(\mathcal O_F\),
so
\[
 \mathcal O_{\mathfrak c}^\times=\mathcal O_F^\times.
\]


With this equality,
in the convergence half-plane group the nonzero lattice
vectors by the corresponding integral finite ideles
modulo the completed order units. Dividing each
idele into a class representative and an element of
\(E^\times\) identifies the fiber with the lattice vectors
modulo the global order units. Homogeneity identifies
its summand with the Hecke weight. The local
order conditions and the multiplicativity of the weight
give the Euler product. At a place of the base conductor
the character has the prescribed exact conductor.
In its nonunit layers the sum over the last available
residue unit quotient vanishes by that conductor, leaving
the primitive local factor. At all other unchanged
places the full order gives the usual geometric series.
The equality of the continued sums follows from the
identity theorem, and at the desired value one
substitutes
\[
 X=q^{-s}.
\]


\proofstep{The sign in the inverse-class map.}
Pair removal will compare coefficients at a split prime through their ray classes. The transfer must therefore use the same convention for inversion as the distribution relation.

For removal of one pair of auxiliary places, evaluate
the abelian group-ring data on any character of the
remaining Artin labels. Such a character is unramified
at the removed place. Summing the removed coordinate
amounts to forgetting it in the inverse-class map.
A uniformizer at either place above \(v\) acts on the
remaining labels through its inverse Artin class.
Define \(Y\) to be the inverse of the image of \(\chi(P)\)
in the coefficient field.

With this reciprocity convention the unramified twisting factor at \(P\) is
the value of the remaining character on its Frobenius.
Thus evaluation at a remaining character \(\psi\) gives
\[
 u_\psi=Y^{-1}\psi(\operatorname{Frob}_P),\qquad
 (1-u_\psi X)(1-u_\psi^{-1}X)+(q-1)X^2 .
\]
The second local root is \(u_\psi^{-1}\).
The polynomial gives the same
multiplicative relation at every tested character, so
finite Fourier inversion proves the group-ring
distribution. With all auxiliary places removed, the
sum of the base coefficients is a nonzero scalar times
the untwisted Hecke \(L(s,\chi)\). The character is unitary
and is ramified at the base conductor, so the classical
nonvanishing theorem for unitary Hecke \(L\)-functions
on \(\RePart z\ge1\), away from the pole of a norm
character, applies. Both the coefficient bound and this
distribution refer to the same inverse-class map and
the same period. \label{cor:toric-twist-5}
\end{proof}

\subsubsection*{The theta polynomials from pair removal}
\hypertarget{source-twist-euler}{}

The toric distribution relation describes the effect of adding a split prime. To use it in the simultaneous choice of Section~\ref{subsec:retained-twist}, encode pair removal by specialization in the group ring before evaluating any character.
The first and second halves use the multiplicative group-ring coordinates
\[
 U_i=1+X_i,\qquad V_i=1+Y_i,\qquad U_i^L=V_i^L=1,\qquad L=\ell^M.
\]

The group ring of the coordinates is the quotient by
\(U_i^L-1,V_i^L-1\). In an integral coefficient polynomial one
can use the exponent \(L-1\) for the inverse of a group
coordinate. Equality at all \(L\)-torsion tuples after setting
\(V_i=1\) gives equality modulo \((Y_i,U_i^L-1)_i\) over the
original integers. Indeed division in each variable by the
monic polynomial leaves a remainder of degree less than \(L\)
in each variable. In a characteristic-zero splitting field
successive univariate interpolation kills that remainder.
This proof involves no division by a group order in the
coefficient ring.

\begin{lemma}[Removal of pairs for abelian theta values]\label{lem:abelian-pair-removal}
For an abelian label at an admissible critical integer \(s\ge0\),
use the uniform denominator \(\ell^{N_0}\) and the two halves of the
chosen auxiliary pairs. The cleared full and first-half theta values
have integral polynomial representatives \(f_M,g_M\) such that
\[
 f_M\equiv B_Mg_M\pmod{(Y_i,U_i^L-1)_i},\qquad
 v(g_M(0))=N_0+v(\theta_\lambda(1)).
\]
Each second-half factor of \(B_M\) reduces to a polynomial in
\(X_i\) of order zero or two at zero.
\end{lemma}

\begin{proof}
First consider an abelian label at an admissible critical
integer \(s\ge0\). The input denominator \(\ell^{N_0}\)
makes each coefficient of its ray-class theta element integral,
uniformly in the auxiliary level. Push this element to the
coordinate group ring by the chosen characters and choose a
polynomial lift \(f_M\). More explicitly, lift each ray class
to an Artin element, use its character coordinates as exponents
in \([0,L-1]\), and sum the resulting monomials in \(U_i,V_i\)
with the cleared coefficients. The characters factor through
the ray group. Evaluating gives exactly
\(\ell^{N_0}\theta_\lambda(\psi)\). Use \(g_M\) for the
same construction at the union of the first half of the pairs,
using just the first-half characters.

\proofstep{The second-half factors.}
The pair-removal identity of the theta input at a pair of
norm \(q\) and at a character unramified on that pair gives
the factor
\[
 E_q(A)=(1-Aq^{-s})(1-A^{-1}q^{-s})+(q-1)q^{-2s},
 \qquad A=\lambda(\mathrm{Frob})\psi(\mathrm{Frob}).
\]
All these coefficients are integral at \(\ell\) because
\(q\equiv1\bmod 2L\). Remove the second half of the pairs
with all \(V_i=1\). At pair \(i\) the Frobenius character
value is exactly \(U_i\) on evaluating, by the identity
matrix condition on the first-half characters. Thus \(B_M\)
is the product over the second half of the integral polynomial
representatives of \(E_q(\lambda(\mathrm{Frob})U_i)\).
Successive removal and the interpolation argument above give
\[
 f_M\equiv B_Mg_M\bmod(Y_i,U_i^L-1)_i.
\]


\proofstep{The first-half unit test.}
To compute \(g_M(0)\), remove the first half at the trivial
character. Choose the open normal subgroup inside the kernel
of the residue character \(\bar\lambda\). On the prescribed
tame element \(z\), \(\bar\lambda(z)\ne1\). Indeed at the
distinguished place of norm \(q_0\) the character label there
agrees with \(\delta^{\pm1}\) after taking a \(q_0\)-power
iterate (power \(q_0^n\) for some \(n\)). The finite
quotient values of \(z\) have order prime to \(q_0\ell\).
Writing a bar for reduction at \(\ell\), the choice of \(z\) gives
\[
 (\overline{\delta(z)}-1)(\overline{\delta(z)}+1)=\overline{\delta(z)}^2-1\ne0.
\]
These are the tests on \(z\): take the
prime-to-\(q_0\ell\) part of a generator at the distinguished
order-15 place. Its value retains the prime-to-\(\ell\)
part of a primitive fifteenth root, of order divisible by
three or five and prime to \(q_0\). It follows that
\(\bar\lambda(z)^{q_0^n}\ne1\).

By the first Frobenius test, the factors at the first-half
places reduce at the trivial character to
\((1-\bar\lambda(z))(1-\bar\lambda(z)^{-1})\), a unit.
Consequently
\[
 v(g_M(0))=N_0+v(\theta_\lambda(1)).
\]
Here the value on the right is at empty auxiliary level.
For a second-half factor write \(x=\bar\lambda(\mathrm{Frob})\).
Its reduction is \((1-xU_i)(1-x^{-1}U_i^{L-1})\).
For \(x\ne1\) it has unit constant term. For \(x=1\)
it equals \(X_i^2\sum_{j=0}^{L-2}(1+X_i)^j\), whose last
factor has constant \(L-1\ne0\) in the residue field.
The denominator bound and the pair relation prove the abelian coefficient condition with \(n_*=2\). The distinction between the
two halves of the primes is essential to the constant
coefficient estimate. \label{cor:twist-euler-1}
\end{proof}

\begin{proposition}[Determinant and signed index for the row theta values]\label{prop:row-pair-removal}
Let \(B\) be the prescribed row lattice of rank \(n\), with its
\(n\) independent projected vectors at empty level and the uniform
clearing scalar \(x\). Write \(C_Q\) for their cleared coordinate
matrix at level \(Q\), \(Q_1\) for the first half of the auxiliary
pairs, and \(b=\operatorname{ord}_O(\det C_0)\). Then
\(f_M=\det C_Q\) and \(g_M=\det C_{Q_1}\) satisfy the same
pair-removal congruence as in Lemma~\ref{lem:abelian-pair-removal},
with the second-half determinant factors displayed below, of residual
order at most \(rn\) in each variable. For any lattice containing the specialized columns, let \(\Delta\) denote its excess signed index over the original column index. Then
\[
 \operatorname{ord}_O g_M(0)=b,\qquad f_M(u-1)\ne0\ \Longrightarrow\ \Delta\le v(f_M(u-1))-b/e_O.
\]
\end{proposition}

\begin{proof}
\proofstep{Clearing the columns on both sides.}
For a row, let \(B\) be the prescribed free lattice of rank
\(n\), \(A\) the commuting central operator of the row
calculation, and \(\theta_1,\ldots,\theta_n\) the selected
projected vectors at empty level. They are independent and
span the rational space. Choose a basis of \(B\).
There is a single denominator
\[
 x=\ell^c\ne0
\]
for the
universal lifts of this finite list, valid at every \(M,Q\)
disjoint from their levels. Projection takes the lattice
containing these cleared lifts into \(B\).

Take their basis
coordinates, with the universal characters still present,
and push the group coordinates to the polynomial presentation
above. This gives an integral square matrix \(C_Q\) such
that at every torsion character its columns are the coordinates
of \(x\) times the projected specializations. At empty level
it gives the constant matrix \(C_0\) of the \(x\theta_j\).
In particular \(\det C_0\ne0\). With \(Q_1\) the first-half set, the two determinant polynomials and the base valuation are
\[
 f_M=\det C_Q,\qquad g_M=\det C_{Q_1},\qquad b=\operatorname{ord}_O(\det C_0).
\]

This choice cancels the denominator in the signed index.
Indeed, if \(\det C_Q(u-1)\ne0\), those specialized columns
before clearing span a lattice of normalized relative index
\(v(\det C_Q(u-1))-n v(x)\) in \(B\). Every lattice in
the same rational space that contains them has relative
index at most this, by elementary divisors. Subtract the
corresponding index for the original \(\theta_j\) to get
exactly \(v(f_M(u-1))-b/e_O\). The formula also holds for
fractional lattices extending beyond \(B\).

\proofstep{The primitive local function.}
For a place \(P\) of norm \(q\), with integral Hecke matrices
\(T_{a,P}\) in the basis of \(B\), put
\[
 H_P(Y)=I+\sum_{a=1}^r(-1)^a q^{a(a-1)/2}Y^aT_{a,P}.
\]
Work in the split local model at the chosen member \(P\).
To distinguish the matrix variable from the rank of the row
lattice, denote its number of columns by \(n_{\rm loc}\).
It has \(r\) rows. The split base frames
identify the vacuum with the characteristic function of
integral matrices. Sum over the character coordinate at a
deleted pair. If the reduced rows span an \(r\)-plane, its
complements are the graphs of maps from a fixed complement
to that plane. For a matrix \(z\) the normalized count is
\[
 q^{-r(n_{\rm loc}-r)}
 \#\{\text{complements of dimension } n_{\rm loc}-r\}
 =
 \begin{cases}1,&\operatorname{rk}\bar z=r,\\0,&\operatorname{rk}\bar z<r.\end{cases}
\]
For smaller rank the dimension condition prevents complementarity.
This is the primitive local function on the integral vacuum.

Let \(z\) be integral and \(c=\dim\ker(\bar z^{\mathsf T})\).
The right cosets of the index-\(a\) minuscule Hecke double
coset parameterize the lattices \(D\) between
\(\varpi O_P^r\) and \(O_P^r\) of residual codimension \(a\):
send a representative \(\gamma\) to \(\gamma O_P^r\), and
use elementary divisors for surjectivity. The condition
\(\gamma^{-1}z\) integral means that each column is in
\(D\). Equivalently the annihilator of \(D/\varpi O_P^r\)
is an \(a\)-plane in \(\ker(\bar z^{\mathsf T})\).
Thus there are \(\genfrac{[}{]}{0pt}{}{c}{a}_q\) translates.
The subspace recursion obtained by intersection with a hyperplane
gives the polynomial identity
\[
 \sum_{a=0}^c (-1)^a q^{a(a-1)/2}
    \genfrac{[}{]}{0pt}{}{c}{a}_q Z^a
   =\prod_{j=0}^{c-1}(1-q^jZ).
\]
At \(Z=1\) the product is one for \(c=0\) and zero for \(c>0\).
Hence the alternating Hecke sum is precisely the primitive
indicator. Off the integral vacuum every term vanishes since
\(D\subset O_P^r\).

\proofstep{The normalized Hecke translates.}
In the split Weil action the right translate of the vacuum
by an index-\(a\) representative has the indicator just
counted with modulus and splitting-character multiplier
\[
 q^{a n_{\rm loc}/2}\mu(\varpi)^a.
\]
Here \(n_{\rm loc}\) is even. This is
the change-of-variables factor \(|\det\gamma|^{-n_{\rm loc}/2}\)
on the matrix argument and the unramified splitting
character on its determinant. Dividing the sum of translates
by this same multiplier is the normalized \(T_{a,P}\) on the
row.

The other local components carry the shorter coordinate
datum: the coordinates agree by restriction, and replacing
the deleted test by the constant one on its integral support
restores the vacuum. The identity of local functions passes
to theta series by linearity and right translation. If \(b\) is the remaining Artin label, the determinant of a representative in degree \(a\) changes it by
\[
 b\longmapsto b+a\delta,
\]
where \(\delta\) is the Artin-coordinate vector of the uniformizer at \(P\). This follows
by taking determinants of the representatives, all of which
are \(\varpi^a\) times a local unit and all remaining
characters are unramified at \(P\). With the inverse Artin map of Lemma~\ref{lem:toric-split-distribution} \(\delta\) is minus the vector of
their arithmetic Frobenius values.

For a character descending to the shorter set, write \(Y\) for its Frobenius value. Translating the index back by \(-a\delta\) supplies the factor
\[
 \chi(-a\delta)=\chi(-\delta)^a=Y^a.
\]
Injectivity of
the row realization transfers this equality to the prescribed
Hecke matrices. This gives \(H_P(Y)\) at the signed character
value on the representative selected by the row.

Projection commutes with the
Hecke operators. Each of the \(n\) columns therefore changes
by the very same matrix. Applying the relation successively
and taking determinants gives
\[
 B_M=\prod_{i=1}^{R_*} \det H_{P'_i}(U_i^{c_i}),
       \qquad c_i\in\{1,L-1\},
\]
with \(P'_i\) in the second half. The identity-matrix
Frobenius test gives \(U_i\) on the specified representative,
conjugating the place inverts the value by anticyclotomicity,
and the sign of the row can invert it once more. This proves
the formula at torsion tuples with trivial second half and
hence the required ideal congruence by interpolation.

Over the residue field
\[
 h(Y)=\det H_{P'_i}(Y)
\]
satisfies
\[
 h(0)=1,\qquad\deg h\le rn.
\] The order of
\(h((1+X)^{c_i})\) at \(X=0\) is the multiplicity of the
root \(Y=1\) of \(h\), including multiplicity zero. Indeed
\((1+X)^{c_i}-1\) has nonzero linear coefficient there.
Each factor is therefore \(X_i^{n_i}\) times a polynomial
with unit constant, \(n_i\le rn\), uniformly in \(M\).

\proofstep{Units at the first half of the pairs.}
The factors from the second set of auxiliary primes have already been identified. The determinant coefficient also depends on what survives when the coordinates from the first set are specialized.

It remains to check the first-half determinant units at the
trivial character. For each signature \(s\), write \(H_{P,s}\) for its unramified Hecke polynomial with the prescribed character \(\eta_s\). Evaluation at one gives
\[
 H_{P,s}(1)=\det(1-\eta_s\rho_s(\mathrm{Frob}^{-1})).
\]

At the distinguished place the twist \(\eta_s\rho_s\)
acts by the scalar \(\delta\) or \(\delta^{-1}\) on inertia.
Intersect the kernels of all the reduced twisted matrices and
the conjugate of this intersection. This is an open normal
subgroup because there are finitely many continuous residual
actions. At either place of a first-half pair the inverse residual Frobenius is scalar. Its two possible values are
\[
 \overline{\delta(z)},\qquad \overline{\delta(z)}^{-1},
\]
since conjugation transports inertia and inverts \(\delta\). Thus the displayed determinant is a
unit for every signature.

Suppose \(\bar H_P(1)\) is singular over the algebraically
closed residue field. Its kernel is nonzero and stable under
all commuting Hecke operators. Restrict successively to a
nonzero eigenspace of a nonscalar operator. This gives a common
eigenvector \(w\) in the kernel. Its residual signature \(s\)
therefore satisfies
\[
       0=\bar H_P(1)w=\overline{H_{P,s}(1)}\,w .
\]
The first-half test just proved makes \(H_{P,s}(1)\) a unit,
a contradiction. Thus \(\det H_P(1)\) is a unit on the row.
Deleting the first-half pairs at the trivial character gives
\[
                   g_M(0)=\det C_0\prod_P\det H_P(1).
\]
This proves the coefficient and index bounds for every allowed
lattice. The following comparison identifies the integral lattice,
residual signature and the common geometric denominator.
\label{cor:twist-euler-2}
\end{proof}

\subsubsection*{Uniform integrality on the ray covers}
\label{subsec:theta-cover-integrality}
\hypertarget{source-theta-cover-integrality}{}

The polynomial argument requires an integral bound on the row coefficients as the squarefree ray level varies. We obtain it first on the cover, using its prescribed frames, and then descend to the level on which the determinant is evaluated.

\begin{paramdefinition}[Admitted levels and coefficient place]\label{hyp:ray-cover-integrality}
We prove the coefficient bound on the ray covers used in the
row determinant. Fix the vacuum pair, the base weight and
coefficient map, the pinned orthonormal frame, and the finite
source translate \(s\). The base level and the varying ray level are
\[
 M_0=g_0g_1,\qquad M=M_0d,
\]
where \(d\) is odd,
squarefree, prime to \(M_0\ell\) and to the conductor and
discriminant exclusions of the host. We use the uniformized
cover with its compatible moduli and adelic frames. Its
coefficient field contains the ray field of level \(M\), the
harmonic coefficients under consideration, and the normalized
values of the splitting characters. A cover with these levels and frames is called admissible for the row at the chosen coefficient place.
\end{paramdefinition}

At the fixed place over \(\ell\), the exponent of \(\ell\) in the level is unchanged. The new marking is a basis of prime-to-\(\ell\) torsion. This torsion is finite étale on the good model, and the compatible bases form a finite étale cover there. Pullback to that cover therefore preserves the good-model condition used in the coefficient bound.

\begin{auxlemma}[Integral finite Fourier expansion]\label{lem:cover-vacuum-fourier}
For a level \(M=M_0d\) of Definition~\ref{hyp:ray-cover-integrality},
let \(\Phi\) be an integral level function on the selected half
lattice, periodic modulo \(M\). If \(k\) generators suffice for
the reduction of the tensor lattice, the vacuum \(v\) and its Heisenberg translates \(W(y)v\) satisfy
\[
 M^k f_\Phi\in\sum_y O\,W(y)v,
\]
with \(O\) the coefficient valuation ring.
The clearing cost at \(\ell\) is at most \(k v_\ell(M_0)\).
\end{auxlemma}
\begin{proof}
Let \(\Lambda\) be the adelic tensor lattice of the vacuum
and write \(\Lambda_e\) and \(\Lambda_{\bar e}\) for its
two projected half lattices. Here \(e\) is the sum of the
idempotents at the places of the level selected by the finite
test vector. A level function \(\Phi\) is supported on
\(\Lambda_{\bar e}\) and is periodic modulo
\(M\Lambda_{\bar e}\). Its values belong to the chosen
valuation ring. Fix a list of \(k\) generators for the
reduction of \(\Lambda\) modulo any positive integer, using
generators for the fixed source and target lattices.
Project these to \(\Lambda_e\) and divide them by \(M\), giving
\(y_i\). The character of \(t\) paired with
\[
 y(\alpha)=\sum_i\alpha_i y_i,
\]
for
\(\alpha\in(\mathbb Z/M\mathbb Z)^k\), is a product of
\(M\)-th roots of unity.

By vacuum self-duality the two half lattices pair perfectly.
For \(t\in\Lambda_{\bar e}\), put \(z_i=\psi(\langle t,y_i\rangle)\).
Then \(z_i^M=1\) and
\[
 \sum_{\alpha}\psi(\langle t,y(\alpha)\rangle)
 =\prod_{i=1}^{k}\sum_{a=0}^{M-1}z_i^a
 =M^k\mathbf 1_{t\in M\Lambda_{\bar e}} .
\]
Indeed, a factor vanishes unless \(z_i=1\); all the \(z_i\) equal one
precisely when \(t\) pairs trivially with the generators. This argument
allows repetitions among the generators. Let \(W(y)\) be
Heisenberg translation by \(y\). Fourier inversion
writes the finite theta vector \(f_\Phi\) as
\[
 M^k f_\Phi
   =\sum_{\alpha}
       \left(\sum_{t\in\Lambda_{\bar e}/M\Lambda_{\bar e}}
             \Phi(t)\,\psi(-\langle t,y(\alpha)\rangle)\right)
             W(y(\alpha))f_0 ,
\]
The coefficients are integral, including at a prime dividing
\(M\). Since \(d\) is prime to \(\ell\), the clearing
exponent in this step is \(k v_\ell(M_0)\).
\end{proof}

\paragraph{Moving to the geometric half lattice.}
The moduli level selects a half lattice at each split
place. Switch the halves that differ from this choice by
the central finite Weil operators. A switch exchanges a
place with its conjugate in the idempotent and carries
each translated vacuum to a translated vacuum in that half,
times the splitting-character scalar. At primes of \(d\)
these switches have identity component at \(\ell\). The
scalar there is integral at \(\ell\) because the local norm factors
are powers of rational residue cardinalities prime to \(\ell\),
and the character uses the norm-normalized values in the
cover field. At primes of \(M_0\) there are finitely many
subsets of places to switch. If \(a_J\) denotes the scalar for such a subset \(J\), a common nonnegative clearing exponent satisfies
\[
 c_{\rm flip}\ge\max_J\{-v_\ell(a_J),0\},
\]
and the associated source
translations form a fixed finite set \(Z\).

For each translate \(sz\), \(z\in Z\), use the split
class representatives away from \(2\ell M\). Concretely
the finite adelic point, after the auxiliary switches, is
the product of a rational isometry \(\gamma\), a finite
point \(h\) supported at split primes away from \(2\ell M\),
and a compact element fixing both the standard lattice and
the level translates of the vacuum. Such a factorization
comes from split-place approximation in the open
stabilizer of the lattice and level vectors. The archimedean
factor is the unique correction to the finite adelic
identity. Write the finite factorization as \(sz=\gamma h k\), with \(k\) the compact factor. At \(\ell\) it gives
\[
 \gamma=(sz)k^{-1}h^{-1},\qquad \gamma^{-1}=hk(sz)^{-1},
\]
with \(h,k\) and their inverses integral there. Thus the
entries of \(\gamma\) and its inverse have a bound on the
\(\ell\)-part of their denominators from \(sz\), even though
the rational isometry varies.

When the Weil kernel is
evaluated at zero, rational automorphy removes \(\gamma\)
and the compact element acts trivially. At infinity this is
a linear substitution in each block of the harmonic
polynomial. It preserves the defining harmonic relations
and weights. In each degree \(h\) its coefficients are
sums of products of \(h\) substitution entries. The finitely
many degrees and the finite set \(Z\) therefore give a
common additional clearing exponent \(c_{\rm sub}\).
The rational source space here is the standard one and
the local entries are tested at all embeddings used in
the cover. These conditions are needed in the split
factorization and the bound on both matrices.

\begin{lemma}[The clearing exponent for the short sections]
\label{lem:theta-short-clearing}
Use the cover and fixed data in Definition~\ref{hyp:ray-cover-integrality}
and the split representative away from \(2\ell M\). At a good frame, the short section for an integral
translate mean and harmonic data of total derivative degree \(h\)
is integral after multiplication by \(\ell^{c_{\rm jet}}\), where
\[
 c_{\rm jet}=v_\ell\!\left(
     M_0^{[E:\mathbb Q]nr+n k_W}(2D_E M_0)^h\right).
\]
Here \(n,r\) are the source and target ranks and \(k_W\) is the
\(\mathbb Z\)-rank of the target marking lattice. The section and
its pinned reading use the same harmonic polynomial and
coefficients of the translate mean.
\end{lemma}

\begin{proof}
Let \(\mathcal T\) be the tuple set of torsion translates of the
source letters at level \(M\). A half-level translate has the
same reading as its representative in \(\mathcal T\): their
difference belongs to the integral lattice of that half, and
their Heisenberg pairing vanishes. The reading of a single
translate at \(t_0\) is the normalized mean with coefficients
\[
 u(t)=|\mathcal T|\,\mathbf 1_{t=t_0}.
\]
In particular its value
in the mean does not acquire a factor of \(|\mathcal T|^{-1}\)
at \(\ell\). Lemma~\ref{lem:cover-vacuum-fourier}, after its indicated
clearing, gives integral linear combinations of these coefficient
arrays.

Let \(R\) be the cover valuation ring and \(P\) the polynomial
in the matrix letters of the moduli frame. The finite module over
\(R\) generated by the coefficients of the array \((u(t)P)_t\)
is torsion-free in the coefficient extension. Choose a basis
\((b_j)_j\) and write the array as \(\sum_j b_j P_j\) with
\(P_j\) arrays over \(R\). Write \(D_\lambda\) for any of the linear operators cutting out
the prescribed harmonic weight over the fraction field of \(R\).
Coordinate extraction in the basis \((b_j)\) gives
\[
 D_\lambda\big((u(t)P)_t\big)
        =\sum_j b_jD_\lambda(P_j)=0
        \quad\Longrightarrow\quad D_\lambda(P_j)=0\quad\text{for all \(j\)}.
\]
Each \(P_j\) has the prescribed weight and derivative degree. By linearity of the jet construction \(J\),
\[
 J\big((u(t)P)_t\big)=\sum_j b_j J(P_j).
\]

The short-section construction on the split representative lattice
supplies the algebraic section and its evaluation by that Weil
kernel average, with the fixed Lie, component, and splitting-character
factors. The hypotheses are the full level marking, the standard
source lattice at \(2\ell M\) with self-duality at \(\ell\), the
integral torsion indicators and constant section, and the prescribed
row degrees. Cuspidality uses the dominant weight and strictly
positive last exponent. In the scalar on this class, the determinant
idele has an ideal representative away from \(2\ell M\).
The normalized unramified character times the square root of its
norm is a local unit at the test place. The remaining inverse
norms and lattice indices, including the index from the split
lattice sandwich, are prime to \(\ell\).

Let
\[
 N=Mu_0
\]
clear the lattice displacement, where \((u_0,2\ell)=1\).
The level index divides \(N^{[E:\mathbb Q]nr}\) and the translate
normalization divides \(M^{n k_W}\) by the standard-lattice and
marking conditions. The latter denominator is harmless for the
indicator array computed above, and may still be included in a
common bound. The good-germ calculation of
Lemma~\ref{lem:regular-theta-jets} costs \((2D_E)^h\).
The additional cost of pullback by \([N]\) is \(N^h\) since
\[
 d[N]=N\,\mathrm{id}
\]
and every derivative monomial has total
degree \(h\). Thus the possible denominator divides
\[
      N^{[E:\mathbb Q]nr} M^{n k_W}(2D_E N)^h.
\]
Its \(\ell\)-valuation is \(c_{\rm jet}\), since \(M/M_0\) and
\(N/M\) are prime to \(\ell\). Its prime-to-\(\ell\) part is a
unit at the frame. This proves the bound there. Compatibility
with degenerating integral frames is proved separately below.
\end{proof}

\paragraph{Good-reduction germs and the connection.}

\begin{lemma}[Denominators of derivatives at a good frame]\label{lem:regular-theta-jets}
For an integral regular germ in the specified theta
trivialization at a good frame, let \(h\) be the order of a jet
in integral tangent directions. If \(a_h\) is its value in the modified polarization gauge and \(R\) is the valuation ring of the frame, then
\[
 (2D_E)^h a_h\in R.
\] The bound descends from a finite extension carrying
the required torsion points.
\end{lemma}

\begin{proof}
Work over a complete valuation ring \(R\) for a good
frame. Use the compatible abelian scheme and the rigidified
symmetric theta bundle. A finite extension carries all
\(N^2\)-torsion used in the short sum and the chart with the
pulled-back frame. Multiplication by \(M\) kills the
level sections. Thus if \([N]t\) is a level section and
\(M\mid N\), then \([N^2]t=0\).
The degree-one section is the integral generator specified
by the frame. The cubical isomorphism
\[
 [N]^*\mathcal L\simeq\mathcal L^{N^2}
\]
preserves the
rigidification and extends over \(R\). It is the normalized
pullback in the short-section formula. After this pullback
the two moves are symmetric theta group translates at
\(N^2\)-torsion. The oddness of \(N\) allows the symmetric
lifts to be chosen of finite order. Each torsion section
extends to the abelian scheme by properness.

The extended torsion section lies in the kernel of the polarization of \(\mathcal L^{N^2}\)
over \(R\), because this equality holds on the generic
fiber and the dual abelian scheme is separated. The
theta group over this section is a multiplicative-group
torsor over the local base, so it has an integral lift.
Write
\[
 g_{\rm fin}=c g_R,
\]
where \(g_{\rm fin}\) is the
finite-order generic lift and \(g_R\) is the integral lift.
Take \(m\) divisible by the order of \(g_{\rm fin}\) and the
order of the underlying torsion section. Then \(g_R^m\) is
a scalar in \(R^\times\), and
\[
  1=g_{\rm fin}^m=c^m g_R^m,\qquad m v(c)=-v(g_R^m)=0 .
\]
Thus \(c\in R^\times\) and both translates extend in the
prescribed normalization. This argument permits
torsion of order divisible by the residue characteristic.
Completing in the integral rigidified trivialization
gives power series over \(R\) for these germs.

\proofstep{Changing the biextension trivialization.}
Write \(B(X,Y)\) for the formal biextension in an integral
chart and \(g(X,Y)\) for a logarithmic gauge over the
fraction field. It vanishes on the axes. The identities
for its two addition defects, transpose defect and CM
balance defect are the logarithms of the corresponding
integral units with constant term one. In particular
their constant first derivatives and their mixed
quadratic coefficients are integral. Write \(\beta\)
for the bilinear coefficient matrix of \(g\), and
\(A(a)\) for the action of \(a\in\mathcal O_E\) on the tangent
space. If the tangent rank is \(m\), integrality of the rigidified mixed coefficients gives
\[
 A(a),\quad \beta^t-\beta,\quad A(a)^t\beta-\beta A(\bar a)\ \in M_m(R).
\]

Choose integral pairs \(x_i,y_i\) giving the separability
element times \(D_E\) in \(E\otimes_{\mathbb Q}E\).
One can take a basis and its trace-dual basis with
the latter multiplied by \(D_E\). They satisfy
\[
 \sum_i x_iy_i=D_E,\qquad \sum_i ax_i\otimes y_i=\sum_i x_i\otimes ay_i\quad(a\in E).
\]
Put
\[
  \gamma=\sum_i A(x_i)^t\beta A(\bar y_i),\qquad
  \beta_0=(\gamma+\gamma^t)/(2D_E).
\]
The latter is symmetric and \(A(a)^t\beta_0=\beta_0 A(\bar a)\).
Put \(T=\gamma-D_E\beta\). The balance defect makes \(T\) integral.
Symmetrizing now gives
\[
 2D_E(\beta-\beta_0)=D_E(\beta-\beta^t)-T-T^t .
\]
Each term on the right is integral in the prescribed frame.

Subtract the symmetric balanced bilinear form \(\beta_0\)
in the logarithmic coordinates of the formal group.
The resulting gauge still satisfies the required identities.
In taking the harmonic value this change contributes
contractions of opposite CM types. These are zero
because the polynomial is in the joint harmonic kernel.
The same contraction rule gives the invariance of the
harmonic value under change of chart with the fixed
tangent and theta frames.

\proofstep{Integral derivatives in the modified gauge.}
Let \(D_u\) be the invariant derivation on the formal group
in an integral tangent vector \(u\), and let \(q(X)\)
be the logarithmic correction obtained by evaluating the
gauge at \((X,[w_0]X)\), with \(w_0\) the fixed integral
polarizing element. Let \(d_0=2D_E\). We have
\(q(0)=D_uq(0)=0\) and
\[
 d_0 D_uD_v q\in R[[X]].
\]

For the last statement, differentiate the two addition
identities at the two identity arguments. The
nonconstant terms are derivatives of the logarithms of
the integral units in the addition defects. Such first
derivatives are integral by \(d\log B=B^{-1}dB\).
The remaining terms are the bilinear coefficient
evaluated on \(u,[w_0]v\) and \(v,[w_0]u\), cleared by \(d_0\).
The algebra generated by \(R[[X]]\) and all
\(d_0D_uq\) for the finitely many tangent directions
has integral constant coefficients. It is stable under
each \(d_0D_u\).

For a regular germ \(f\) and integer
bundle power \(j\), applying \(d_0D_u\) to
\(a\exp(jq)\) acts on the coefficient by
\[
  a\longmapsto d_0D_ua+j(d_0D_uq)a .
\]
Starting from \(a=f\) proves integrality of every
degree-\(h\) derivative after multiplication by \(d_0^h\).
The exponential contributes constant term one, so
no bound on its higher coefficients is used. Apply this
letter by letter and use the integral polynomial
coefficients. A finite extension multiplies normalized
discrete valuations by a positive ramification index,
so the resulting inequality descends to the original
good frame.
\end{proof}

\paragraph{Integral tests at degenerating frames.}\label{par:theta-degenerate-arc}
We use the good-at-\(\ell\) level condition of the
admitted cover and the toroidal moduli model over the
selected coefficient place. The degeneration description
uses the split Raynaud extension and polarized period
lattice after finite extension. The relative toroidal
construction over the complete base turns these data
back into a semiabelian family. \cite{faltings-chai}
For a definite signature
all the test frames have good reduction. In the other
signature consider any integral test frame. By the
valuative criterion on the compactification it extends
to a trait.

If the trait lies in the abelian locus, the
preceding estimate applies. Otherwise, after a finite
extension, split the degeneration data and write its
period pairing as
\(\varpi^{Q(\lambda,\mu)}u_{\lambda,\mu}\) with
unit factors and the integral degeneration form \(Q\).
The period pairing on the arc and its specialization are
\[
 X^{Q(\lambda,\mu)}u_{\lambda,\mu}\ \;\longmapsto\;\ \varpi^{Q(\lambda,\mu)}u_{\lambda,\mu}.
\]
The linear
period and polarization data use the same rescaling.
The integral frame lifts on this arc with constant
tangent coefficients and the same degree-one
theta generator. This uses the derivative
identification for the semiabelian uniformization
and the integral generator at the base point.
Inverting \(\varpi X\) puts the family in the
generic abelian locus. Compatibility of algebraic
evaluation with the arc and its specialization
expresses the value of a scaled section as
\(H(X)/(\varpi X)^b\), where \(H\in R[[X]]\)
and \(b\ge0\).

\begin{lemma}[Integrality after specialization of the toroidal arc]\label{lem:theta-arc-integrality}
On the framed toroidal arc of Paragraph~\ref{par:theta-degenerate-arc},
specializing at \(X=\varpi\), let a scaled section be
\(H(X)/(\varpi X)^b\), with \(H\in R[[X]]\) and \(b\ge0\).
Suppose the section is bounded on the normalized complex frame
and integral at the algebraic good-point tests with the same
clearing power. On the arc and the original trait, respectively, one has
\[
 H\in \varpi^b X^b R[[X]],\qquad H(\varpi)/\varpi^{2b}\in R.
\]
\end{lemma}

\begin{proof}
\proofstep{Test the order at \(X=0\) on the complex frame.} For the complex test,
embed the coefficient fraction field abstractly into
\(\mathbb C\) over the specified coefficient embedding.
Use the holomorphic uniformization of the arc with
\(X=\exp(2\pi i z)\) near the puncture. To compare
the two evaluations, take an affine open chart of the
generic moduli scheme containing the Laurent-series point
and use its algebraic frame. The chart-coordinate
comparison for the holomorphic uniformization of the
split degeneration gives convergent Laurent series for
the arc's chart coordinates. The framed comparison uses
the constant tangent periods and the nonzero degree-one
theta generator of this same arc.

The frame comparison gives convergent Laurent series for the entries of the change-of-frame
matrices and for the ratio of theta generators, with
values equal to the corresponding analytic quantities.
The determinants of the frame matrices and the theta
ratio are invertible formally. Each such denominator germ can be written
\[
 d(X)=X^a u(X),\qquad u(0)\ne0,\qquad d(X)^{-1}=X^{-a}u(X)^{-1}.
\]
After shrinking the punctured disk the last germ converges. In the chart, evaluation
of a regular section is evaluation in the chart frame
followed by its algebraic weight factor. The latter uses
just the entries of the frame change, inverse
determinants and inverse theta ratio. Sums and products
of the convergent readings therefore give the analytic
value with Laurent coefficients equal to the transported
formal coefficients.  \label{cor:theta-cover-integrality-1}

To bound its value, normalize the complex
frame by a Borel change to unit Gram frame.
The cusp section is bounded there in its invariant
norm. Along the arc the diagonal tangent ratios
and the generator ratio in the weight character are
bounded as follows. In the symplectic normalization of
the arc the imaginary period matrix is
\[
 A+\operatorname{Im}(z)Q.
\]

Positivity for all sufficiently large \(\operatorname{Im}(z)\)
implies \(Q\ge0\). For a fixed \(z_0\) in that region the
matrix dominates its positive value at \(z_0\). More explicitly, writing
\[
 Y(z)=A+\operatorname{Im}(z)Q,
\]
the order inequalities for \(\operatorname{Im}z\ge\operatorname{Im}z_0\)
are
\[
 Y(z)\succeq Y(z_0)>0,
 \qquad Y(z)^{-1}\preceq Y(z_0)^{-1},
 \qquad \det Y(z)\ge\det Y(z_0).
\]
The inverse bound controls the tangent ratios and the determinant
bound controls the generator norm in the same period frame. The squared Hermitian length of each constant
tangent vector is computed using the inverse imaginary
period matrix. The squared diagonal ratios from
Gram orthogonalization are bounded by these lengths
times the fixed block scalars.

For the theta generator \(e\), the Gaussian integral over a symplectic period parallelepiped is
\[
 \|e\|^2=\frac{|\kappa|^2}{\sqrt{\det(2\operatorname{Im}\tau)}} ,
\]
where \(\kappa\) is the arc's constant generator scalar. The unit frame uses a generator of squared norm
one. Thus the generator ratio is bounded as well.
The dominant weight uses these ratios only with
nonnegative exponents. Thus the complex Laurent
function is bounded as \(\operatorname{Im}z\to\infty\).
Its negative coefficients vanish by Cauchy's formula.
Injectivity of the coefficient embedding gives
\(X^b\mid H\) in \(R[[X]]\).

\proofstep{The nonarchimedean valuation.}
Holomorphic testing on complex frames bounds the order of vanishing. To use this information in division, we still need the valuation of the coefficient in the nonarchimedean frame.

Write
\[
 H=X^b H_1.
\]
For the second test use the
valuation at \((\varpi)\) in \(R[[X]]\), extended to the
completion of its fraction field. The coefficient formula for this valuation is
\[
 v(\varpi)=1,\qquad v(X)=0,\qquad
 v\!\left(\sum_i a_iX^i\right)=\min_i v(a_i)\quad(a_i\in R).
\]
The arc is now in the abelian locus. The frame still
gives an integral point there, since its tangent and
theta generators lift on the arc. The good-point
bound extends to this valuation by the smooth
integral flag charts. Concretely, on an open
abelian chart the section is a rational function
in integral chart coordinates. If it had a pole
at an integral test with arbitrary residue extension,
one can retain its finite pole order on a point with
algebraic coefficients by smooth lifting of the
finite congruences in that chart. The condition to
stay in the open locus is one of these congruences.
Such a pole contradicts the algebraic good-point tests.

Since \(X\) is a unit in this valuation, the valuation of the
section is \(\min_i v([X^i]H_1)-b\). The good-point tests give
\[
 \min_i v([X^i]H_1)\ge b.
\]
Thus \(H=\varpi^bX^b H_2\)
with \(H_2\in R[[X]]\). The value on the original framed trait
is \(H_2(\varpi)\in R\). Both tests used the section after its
fixed multiplication by a power of \(\ell\).
\label{cor:theta-cover-integrality-2}
\end{proof}

Lemma~\ref{lem:theta-arc-integrality} thus extends the good-point
bound to the degenerating frames of the same cover.

\begin{proposition}[One phase and one denominator on the admitted covers]
\label{prop:theta-uniform-cover-clearing}
Fix the data at the start of this subsection and a finitely generated
integral lattice of input polynomials. There is a phase \(\eta\) from the
base frames and a clearing power \(\ell^c\) such that the comparison
with the algebraic reading of each admitted coefficient is integral
at every prescribed place above \(\ell\), after scaling by
\(\eta\ell^c\). One can take
\[
    c=k v_\ell(M_0)+c_{\rm flip}+c_{\rm sub}+c_{\rm jet}+c_P,
\]
with a common maximum over the degrees of the finitely many
polynomial generators when needed. The varying cover, level function,
twist, places and class label do not enter this choice.
\end{proposition}

\begin{proof}
The source frame in the norm comparison changes by a unitary map
preserving the vacuum lattice. On each component compare the
vacuum identities on a nonzero theta value. If \(c_1,c_2\)
are the two component constants, and \(A\) is the product of
that value, the Gaussian and the nonzero character factor, the
comparison of phases is the comparison of \(c_1 A\) and \(c_2 A\).
The common nonzero factor cancels in the phase ratio:
\[
 \frac{c_1 A}{c_2 A}=\frac{c_1}{c_2}.
\]
This reads the base vacuum phase. The dominant pins
on the cover read these same constants after pullback. Multiplying
its inverse by the inverse phase of the fixed lowest weight vector
in the archimedean harmonic map gives \(\eta\).

The matrix coefficients of that harmonic map on a fixed set of
lattice generators form a finite set of algebraic numbers in each
admitted coefficient field. For each coefficient \(a\), take a
positive integer \(d_a\) for which \(d_a a\) is integral. Then
\[
 c_P=\max_a v_\ell(d_a)
\]
works at all places above \(\ell\),
since \(d_a/\ell^{v_\ell(d_a)}\) is a unit there. The short-section
construction over the field gives the rational section. The local
construction of Lemma~\ref{lem:theta-short-clearing}, extended
by Lemma~\ref{lem:theta-arc-integrality}, gives integral comparison
sections. They have the same pinned adelic reading after the indicated
scaling, so injectivity identifies each with the rational section
at that place.

The determinant twist is a finite-order character on the
norm-one determinant quotient, trivial on the rational points
and the cover congruence subgroup. The indicated ray field gives
its values on the components. Each such value is a root of unity,
so multiplication by it preserves the integral section lattice.
It costs no clearing power. Summing the exponents from Fourier
inversion, the half switch, the change of variables, the short
section and the fixed harmonic map gives the displayed value of
\(c\). \label{cor:theta-cover-integrality-3}
\end{proof}

\subsubsection*{Descent to the original row}
\label{subsec:theta-cover-descent}

The preceding coefficient bounds were obtained in frames on the covers. To measure a coordinate in the original cusp lattice we must also account for the different of the coefficient extension in which those frames were evaluated.

\begin{proposition}[Uniform descent from ray covers]\label{prop:ray-cover-descent}
Suppose a section on the prescribed ray cover descends as a
holomorphic form and is integral at all good frames above
\(\mathfrak q\mid\mathfrak p\mid\ell\) of its coefficient
extension \(L/K_0\). Each coordinate \(c\) in an integral cusp basis at \(\mathfrak p\) satisfies
\[
 c\in\mathfrak D_{L_{\mathfrak q}/(K_0)_{\mathfrak p}}^{-1}
\]
under the chosen local embedding.
Unramified descent preserves integrality. For the universal row
covers contained in the fraction field of the fixed coefficient
DVR \(O\), multiplication by \(\ell^{e_O^2}\) suffices at the
chosen place, where \(e_O=v_O(\ell)\).
\end{proposition}

\begin{proof}
\proofstep{Descent of the holomorphic span.}
Let a form on the cover be invariant under the base level group.
The character relating its adelic and holomorphic realizations is invariant
under that group by the level condition on the splitting
character. Hence the weighted holomorphic form factors
through the base quotient. The bounded cuspidal holomorphic
space of the base is spanned by the complex realizations of its algebraic
cuspidal sections in the prescribed dominant weight with
positive last exponent.

In the application its degree is
at least \(8fr+4\), where \(r\) is the target rank and
\(f\) is the base degree parameter.
Pull the sections back to the cover. Linear independence persists under
extension to \(\mathbb C\). Among the cover pullbacks,
choose a basis over the cover field \(L\). If \(s_L\) lay outside the \(L\)-span of independent pullbacks \(F_1,\ldots,F_t\), then
\[
 \dim_L\langle F_1,\ldots,F_t,s_L\rangle=t+1.
\]
Their complexifications would remain independent, contradicting complex descent. Thus the rational
descent expresses the section in the \(L\)-span of the
original base sections, over the base field \(K_0\).

For the integral statement used on the ray coefficients,
let \(\mathfrak q\mid\ell\) be the selected place of \(L\)
and \(\mathfrak p\) its restriction to \(K_0\).
The base cusp sections integral at \(\mathfrak p\) form a
full lattice over \({\cal O}_{K_0,\mathfrak p}\). The
lattice is defined by the valuation bounds at all good
algebraic frames. It is finite over the localized base ring
by the unit ball bound on the finite-dimensional space.
For any section \(v\) in the base-field space, a common scalar clears its finitely many coordinates, so
\[
 \exists\,a\in\mathcal O_{K_0,\mathfrak p}\setminus\{0\}:\quad av\ \text{lies in the integral section lattice}.
\]
Thus one can use a
free lattice basis \(F_1,\ldots,F_t\). Write the descended
cover section as \(\sum_i y_i F_i\) with \(y_i\in L\),
using the pullbacks for this equality.

\proofstep{Testing the dual lattice by local traces.}
Write
\[
 s_L=\sum_i y_iF_i.
\]
For a base frame \(x\),
pullback commutes with evaluating each \(F_i\). At a place
\(\mathfrak q\) where \(s_L\) is integral, this is true
in every local embedding of \(L_\mathfrak q\) into a common
valued field. At the other places over \(\mathfrak p\), choose
\(b\in\mathcal O_L\) with \(v_{\mathfrak q'}(b)\) large
enough to clear the finitely many \(y_i\). Then every local
evaluation of \(b\,s_L\) on the lift of \(x\) is integral.
The trace section is
\[
    \operatorname{Tr}_{L/K_0}(b\,s_L)
       =\sum_i\operatorname{Tr}_{L/K_0}(by_i)F_i .
\]
Its evaluation at \(x\) is the sum over the embeddings
of \(b\,s_L\), so the base lattice test applies to it.
Since the \(F_i\) form an integral basis, each coefficient
of this trace section is integral at \(\mathfrak p\).

By weak approximation take \(b\)'s approaching any fixed
\(b_\mathfrak q\in\mathcal O_{L_\mathfrak q}\) at \(\mathfrak q\)
and approaching zero at the other local factors. Local
trace and the closedness of the integral base ring give
\[
   \operatorname{Tr}_{L_\mathfrak q/K_{0,\mathfrak p}}
             (b_\mathfrak q y_i)\in
                   \mathcal O_{K_{0,\mathfrak p}}^{\wedge}
      \quad(b_\mathfrak q\in\mathcal O_{L_\mathfrak q}).
\]
This is precisely
\[
 y_i\in\mathfrak D_{\mathfrak q/\mathfrak p}^{-1},
\]
with \(\mathfrak D\) the local different. The test used
all frames at the chosen place, as needed for the full trace.
For an unramified extension the different is the unit ideal.

For a finite prescribed set of unramified places, take one such
local decomposition \(s_L=\sum_i y_{v,i}F_{v,i}\) per place \(v\).
Rescale each \(F_{v,i}\) by a base scalar that is a unit
at the place where it was selected and clears its poles at
the other base places, and rescale \(y_{v,i}\) inversely.
The Chinese remainder theorem gives \(\rho_v\in L\) with
\[
 \sum_v\rho_v=1,
\]
\(\rho_v\) a unit at its place and
divisible by sufficient powers at all the others. Indeed one
may choose all but one by their local congruences and set
the last to one minus their sum.

At a prescribed place \(w\), the choice of powers is
\[
 v_w(\rho_v)\ \ge\
      \max_i\max\{0,-v_w(y_{v,i})\}\quad (v\ne w).
\]
The \(F_{v,i}\) are integral there by the initial rescaling.
For \(v=w\), both \(F_{w,i}\) and \(y_{w,i}\) are integral by
local descent and the triviality of the unramified different.
It follows term by term that the section
\[
 \sum_{v,i}\rho_v y_{v,i}F_{v,i}=s_L
\]
is integral at \(w\).

\proofstep{A uniform bound on the different.}
The descent formula expresses the new denominator through a local different. It remains to bound that different uniformly as the allowed ray levels change.

In the application to the universal row coefficients the
cover can have ramification at \(\ell\). The actual cover
field embeds in the fixed fraction field of the row DVR
\(O\). Take \(e_O=v_{\!O}(\ell)\) with \(v_{\!O}(O^\times)=0\)
and a uniformizer of value one. The restriction of \(v_{\!O}\)
to \(L\) is a positive integral multiple of the discrete
valuation at \(\mathfrak q\). Writing \(a\ge1\) for the integer multiplying the valuation at \(\mathfrak q\), evaluation on \(\ell\) gives
\[
 e_O=a\,v_{\mathfrak q}(\ell)\ge v_{\mathfrak q}(\ell).
\]
A bound of
\(e_O^2\) suffices for the relative different exponent.

For instance, over the maximal unramified subextension
of \(L_{\mathfrak q}/\mathbb Q_\ell\), a uniformizer has
an Eisenstein polynomial of degree \(a\), the absolute
ramification index. The different has the valuation of
its derivative. The nonzero terms of the derivative
have distinct valuation residues modulo \(a\), so its
valuation is at most that of its leading term,
\[
 a-1+a v_\ell(a)<a^2,
\]
including \(a=1\).
Transitivity of the different then bounds the relative
exponent by the same number. Multiplying by
 \(\ell^{e_O^2}\) makes the inverse-different coordinates
integral. In fact, writing \(e_{\rm abs}=v_{\mathfrak q}(\ell)\)
and \(d_{\rm rel}=v_{\mathfrak q}(\mathfrak D_{\mathfrak q/\mathfrak p})\),
\[
 v_{\mathfrak q}(\ell^{e_O^2}y_i)
       \ge e_O^2 e_{\rm abs}-d_{\rm rel}\ge0
\]
by \(d_{\rm rel}<e_{\rm abs}^2\) and \(e_{\rm abs}\le e_O\). These two numerical inequalities give the uniform estimate
\[
 d_{\rm rel}<e_{\rm abs}^2\le e_O^2.
\]
Its role is important even
when the additional level is prime to \(\ell\), since
the coefficient field also contains finite-order values.

\proofstep{The fields used before specialization.}
To see which covers this argument serves, let \(Q\) and
\(m_{\rm ray}\) be the auxiliary places and the depth
for the universal coefficient. The construction takes
the split rational primes below the selected conjugate
half of \(Q\), with the designated prime of the row twist.
The avoidance conditions make their product a squarefree
auxiliary level, prime to \(\ell\), the base level and the
host exclusions.

Write \(K_{\rm ray}\) for the ray class field at the joint level, \(A\) for the finite set of character values of the finite type and distinguished twists, and \(K_{\rm row}\) for the row coefficient field under the fixed embedding. The ray-field condition and coefficient choice give
\[
 K_{\rm ray}\subseteq K_{\rm ray}(A)\subseteq K_{\rm row}.
\]
The uniformized, dominantly
pinned cover is formed over the resulting number
field. Thus both the field conditions for integrality
on the cover and the ramification bound for descent
hold for these choices. In particular this construction
uses the cover field before specializing a character of
the variable depth.

The finite universal data on this cover are indicators
times finite-character values. Let \(A_Q\) be the class-label
group, \(b(g)=\operatorname{detchar}_Q(g)\) the determinant
label and \(S_{a'}\) the section for label \(a'\). Then the
assembled section is
\[
              \Xi_a(g)=\sum_{b\in A_Q}
                     \mathbf1_{b(g)=b}\,S_{a-b}(g).
\]
Each determinant fibre is a union of components by the
ray-field comparison. The indicators preserve the section
lattice, and the finite-character values are integral. Thus
the sum uses the same integrality bound without a denominator
\(|A_Q|\). The prescribed finite type and the trivial
upper-left reductions make \(\Xi_a\) invariant under the
original level group, so the trace descent applies.

The resulting base sections are still bounded at their
chosen base place. Under the faithful row realization,
their coordinate vectors lie in the localization of the
supremum lattice at that place. For a coefficient vector \(x\), localization gives \(x=y/s\) with \(y\) a bounded global vector and \(s\) a base scalar of valuation zero. Under the chosen embedding,
\[
 v_\ell(s)=0\quad\Longrightarrow\quad s^{-1}\in O^\times.
\]
Thus passing from \(y\) to \(x\) costs no further denominator. Taking the additional clearing power to be
\(e_O^2\), the descent yields a bounded row for every
coefficient, with the complex phase and the fixed
harmonic comparison scalar supplied by the pinned
coefficient map. The clearance depends on the datum
and on the fixed \(O\), independently of the level,
depth and class label. The passage from this complex
reading to the actual rational row lattice is given next.
\label{cor:theta-cover-integrality-4}
\end{proof}

\subsubsection*{The span of the original theta rows}
\label{subsec:theta-row-rank}

Integrality after descent bounds the coefficients of the theta rows. For their determinant to detect the original row space, these rows must also span its good type at the base tuple with the prescribed unipotent character.

\paragraph{Stable range and the lattice operator.}
Write \(n\) for the source rank, \(r\) for the target
rank and \(D_E=[E:\mathbb Q]r\) for the dimension of a
target lattice over \(\mathbb Q\). The source dimension
condition is \(n\ge 4D_E+4\), with \(n\) even as in the
row construction. In the finite Weil model
for the prescribed level, let \(\Phi(L,A)\) denote the
normalized test vector for a self-dual target lattice
\(L\) and a source lattice class \(A\). The test vectors
are summed over the finite source classes with their
stabilizer weights. Their pairings define the matrix
\[
      C(L,L')=\sum_A w(A)
                \langle\Phi(L,A),\Phi(L',A)\rangle .
\]
With the same weights on both entries, its diagonal
is a positive scalar. If \(m=[L':L\cap L']\), the
ratio of the off-diagonal bound to the diagonal is
\(m^{-n/2}\). We give the common-label calculation
that makes this bound independent of the size of the level.

\begin{lemma}[Invertibility of the overlap matrix in the stable range]\label{lem:theta-overlap-invertible}
At the specified finite level let \(n\ge4D_E+4\) be even.
For distinct lattices let \(a(L,L')\) denote their relative overlap. Then
\[
 |a(L,L')|\le [L':L\cap L']^{-n/2}.
\] The resulting matrix \(C\) has
a bounded inverse on bounded automorphic holomorphic frame functions
of the prescribed weight.
\end{lemma}

\begin{proof}
Let \(\mathfrak m\) be the level ideal, coprime to both
\(2\) and \(\bar{\mathfrak m}\), and write \(L_*\) for the reference
target lattice. A target lattice is commensurable with
\(L_*\) by multiplication in both directions by some
\(s\in\mathcal O_E\) prime to \(\mathfrak m\bar{\mathfrak m}\).
The shifted lattice is
\[
 L^{\sharp}=(\mathfrak m L:\bar{\mathfrak m}).
\]

For a source lattice \(A\), let
\(B_L=A\otimes L^{\sharp}\) be its self-dual tensor lattice
in the Weil space, where the tensor convention uses
conjugate source contractions. The normalized shifted
vacuum \(v_L\) is fixed by the operators
\[
 U(\delta)=\varepsilon(\delta)W(\delta)
\]
for
\(\delta\in B_L\). Here \(W\) is the Heisenberg action,
and \(\varepsilon\) is the common lattice semicharacter.
The level sum for \(L\) uses translates \(W(y)v_L\),
with \(y\) labeled by source rows modulo \(\mathfrak m L_*\)
after localization at \(\mathfrak m\).

The allowed labels
are the level subgroup in the first \(r\) rows and zero
in the remaining rows. The reduced matrices and their
unitary character give the same weights for all \(L\).
The two sums admit common representatives: the intermediate lattices agree at the level because the intervening scalars are prime to \(\mathfrak m\). The Chinese remainder theorem supplies representatives in their intersection. Denominator-lattice shifts fix the translate. Both the lift \(y\) and the difference \(d\) contain the factor \(2\) in the level quotient. Hence the vacuum semicharacter and their integral pairing phase equal one. Consequently
\[
 \varepsilon(d)=1,\qquad W(y+d)v_L=W(y)v_L.
\]


\proofstep{Orthogonality from the Heisenberg characters.}
Suppose \(\langle W(y)v_L,W(y')v_{L'}\rangle\ne0\). For every
\(\delta\in B_L\cap B_{L'}\), the two vectors are eigenvectors
of the same unitary \(U(\delta)\), with eigenvalue ratio
\(\exp(2\pi i\, E_A(\delta,y-y'))\) up to the convention for
the sign of the alternation \(E_A\). Put \(\Delta=y-y'\) and \(B_0=B_L\cap B_{L'}\).
Unitarity gives
\[
 \bigl(1-e^{2\pi i E_A(\delta,\Delta)}\bigr)
         \langle W(y)v_L,W(y')v_{L'}\rangle=0
       \qquad(\delta\in B_0).
\]
For a nonzero overlap this gives
\(\Delta\in B_0^\#\), with \(\#\) denoting the alternating
dual.

Choose \(s,s'\) prime to \(\mathfrak m\bar{\mathfrak m}\)
from the two comparisons with \(L_*\). The inclusion
\(ss' B_L\subseteq B_0\) implies
\[
 B_0^\#\subseteq(ss' B_L)^\#
             =(\overline{ss'})^{-1}B_L ,
\]
using \(B_L^\#=B_L\) and the conjugate adjoint rule. The source lattices in the
level sum contain \(N\mathcal O_E^n\) for an integer \(N\)
prime to a level integer in \(\mathfrak m\).
Contracting with \(Ne_k\) consequently gives
\(N\overline{ss'}(y_k-y'_k)\in L^{\sharp}\) in each row.
Take \(a\in\mathfrak m\), \(b\in\bar{\mathfrak m}\) with
\(a+b=1\). The definition of \(L^\sharp\) gives
\[
  bN\overline{ss'}(y_k-y'_k)\in\mathfrak m L,\qquad
  sbN\overline{ss'}(y_k-y'_k)\in\mathfrak m L_* .
\]
The scalar in the second inclusion is a unit modulo
\(\mathfrak m\), hence \(y_k,y_k'\) have the same level
label. This holds for every row.

For equal labels the same \(W(y)\) occurs in each argument.
Since it is unitary and the level-character weights have
modulus one, the normalized level contribution is
\(\langle v_L,v_{L'}\rangle\) times the common number of
labels. For its absolute value put
\(D'=B_{L'}/(B_L\cap B_{L'})\). The projection on the
\(B_{L'}\)-vacuum line, evaluated on \(v_L\), is
\[
      \Pi_{L'}v_L=\frac1{|D'|}\sum_{\delta\in D'}U(\delta)v_L .
\]
The terms are well defined because the denominator fixes
\(v_L\).

Distinct terms have different \(B_L\)-characters
by self-duality of \(B_L\), hence are orthogonal and each
has norm one. Locally the source lattice is free of rank
\(n\), so after the common shift the relative tensor index
is \(|D'|=[L':L\cap L']^n=m^n\). Therefore
\[
 |\langle v_L,v_{L'}\rangle|^2
   =\|\Pi_{L'}v_L\|^2=|D'|^{-1}=m^{-n}.
\]
The relative bound is the same for every source class. Summation
with the identical nonnegative stabilizer weights gives it for
the matrix \(C\).

\proofstep{Counting neighboring self-dual lattices.}
The number of self-dual \(L'\)
with index \(m\) is at most \(m^{2D_E}\). Set
\(J=L\cap L'\). Self-duality gives equal covolumes
and hence both indices over \(J\) equal to \(m\).
The intersection is an index-\(m\) sublattice of
\(L\). The number of index-\(m\) sublattices of
a free abelian group of rank \(D_E\) is at most
\(m^{D_E}\), by triangular Hermite normal form.
Indeed the recurrence in rank \(d\) is
\[
 a_d(m)=\sum_{k\mid m}k^{d-1}a_{d-1}(m/k)\le\sum_{k\mid m}k^{d-1}(m/k)^{d-1}=m^{d-1}\tau(m).
\]
There are at most \(m\) positive divisors, giving the asserted bound.

Now \(L'\) is an index-\(m\) sublattice of the
dual of \(J\), by the same covolume calculation.
The same count bounds the choices for \(L'\) for
each \(J\). Index one occurs only for \(L'=L\).
The normalized absolute row sum off the diagonal
is therefore at most
\[
       \sum_{m\ge2} m^{2D_E-n/2}
          \le \sum_{m\ge2} m^{-2}<0.65 .
\]
The estimate applies to every target lattice in the
system and every rank meeting the size condition.

\proofstep{The inverse and its multiplier.}
Let \(D\) be the positive diagonal of \(C\).
The unweighted bound on the off-diagonal row sum
makes \(D^{-1}C=1+T\) invertible on the bounded
automorphic functions by the series
\(\sum_{j\ge0}(-T)^j\). Each diagonal entry is
the fixed test-vector norm with the finite class
weights, so its inverse is bounded. The coefficient
matrix is compatible with the unit-modulus
automorphy multiplier on the target lattices.
It acts on bounded holomorphic frame functions
with the prescribed weight. The series converges
uniformly on compact subsets, since the weighted
supremum bound gives an ordinary supremum bound
on compact subsets of the frame domain. Thus its
inverse preserves holomorphicity and the multiplier.
Given the original bounded holomorphic family
\(F\), use \(F'=C^{-1}F\) in the unfolding.
\end{proof}

\paragraph{Reproduction on the thin slab.}

\begin{lemma}[Reproduction of the holomorphic frame functions]\label{lem:slab-reproducing-kernel}
For the Gaussian attached to the given stable-range weight, there is
a sufficiently thin slab and a constant \(\rho>0\) such that its
kernel reproduces each bounded holomorphic frame function of the
specified weight there. At an orthonormal frame \(x_0\) the scalar is
\[
 \rho=\frac{\int_{\rm slab}|K(x_0,y)|^2\,dy}{K(x_0,x_0)} .
\]
\end{lemma}

\begin{proof}
Use the Gaussian kernel with the row highest
polynomials in the prescribed weight. The stable
range contains these polynomials in the harmonic
tensor model. Take a finite family orthonormal
for the Fock pairing on this component. At a real
place use coordinates in which the Hermitian Gram
matrix of an orthonormal frame is
\(J=\operatorname{diag}(I_p,-I_q)\), \(p+q=r\).
The frame domain has invertible frames and
positive first \(p\) columns for this form. A thin
slab consists of frames with Gram matrix in a small
open neighborhood of \(J\), at every real place.
Indefinite Gram--Schmidt writes a frame as \(u b\)
in orthonormal coordinates, with \(u\) unitary for
\(J\) and \(b\) upper triangular.

For the first \(j\) columns of the frame \(y\), let
\(G_j=y_{[j]}^*J y_{[j]}\) and put \(\det G_0=1\). Since
\[
 G_j=b_{[j]}^*J_{[j]}b_{[j]},
\]
where the brackets on the square
matrices denote their leading blocks,
\[
 |b_{jj}|^2=\left|\frac{\det G_j}{\det G_{j-1}}\right|.
\]
On a sufficiently thin slab each \(\det G_j\) remains close to
\(\det J_{[j]}=\pm1\). The diagonal entries and their inverses
are therefore bounded. The Gram--Schmidt recursion then bounds
the remaining entries of \(b\) using these pivots and the
bounded entries of the Gram matrix.
The left unitary action preserves the slab and
Lebesgue measure because each unitary matrix has
determinant of absolute value one.

Let \(K(x,y)\) be the polarized Gaussian after
applying the highest polynomials in both arguments.
It is holomorphic in \(x\), and its right upper
triangular covariance is the weight character
\(\chi_\lambda\). For a unitary change \(u\), the two invariances are
\[
 K(y,x)=\overline{K(x,y)},\qquad K(ux,uy)=K(x,y).
\]
In the Gaussian
these changes preserve the Gram factors. Swapping
the two frames conjugates the polarized factors.
In the first argument the upper triangular change
contributes the determinant power of the Gaussian
and the diagonal powers of the two blocks of the
highest polynomial. Terms from same-argument
quadratic factors are killed by the harmonic
contractions.

At an orthonormal frame \(x_0\) the
diagonal Gaussian, in the harmonic letters \(z_l,w_l\),
is a positive determinant factor times
\(\exp(\sum_l t_l z_l w_l)\), with \(t_l>0\) from the
positive imaginary part of the polarized parameter.
For a nonzero polynomial \(P(z)=\sum_\alpha P_\alpha z^\alpha\)
in the chosen harmonic basis, coefficient extraction gives
\begin{equation}\label{eq:slab-gaussian-diagonal}
 \left.P(\partial_z)\overline P(\partial_w)
       \exp\!\left(\sum_l t_l z_l w_l\right)\right|_{z=w=0}
      =\sum_\alpha |P_\alpha|^2\alpha!\prod_l t_l^{\alpha_l}>0 .
\end{equation}
Only matching multi-indices survive the two differentiations.
Summing \eqref{eq:slab-gaussian-diagonal} over the basis and restoring the positive determinant
factor gives \(K(x_0,x_0)>0\).

\proofstep{The two torus characters.}
Consider the difference \(F\) of two functions with the same
value at \(x_0\), the prescribed right upper triangular weight
and the matching left diagonal unitary character. In coordinates
at \(x_0\), Gaussian elimination reduces a frame in the dense
cell to a lower triangular matrix \(L\). Define
\[
 D(z)=\operatorname{diag}(z,z^2,\ldots,z^r),
 \qquad L(z)=D(z)LD(z)^{-1}.
\]
The \((i,j)\) entry on or below the diagonal is \(z^{i-j}L_{ij}\),
so \(L(z)\) extends holomorphically to an invertible diagonal
matrix at zero.

For \(0<|z|\le1\), the first \(p\) diagonal entries
of \(D(z)\) have squared modulus at least \(|z|^{2p}\);
the others have squared modulus at most \(|z|^{2p}\).
Hence \(D(z)^*JD(z)\succeq |z|^{2p}J\), and the positive
Gram block stays positive after left multiplication. The
right diagonal changes that block by a congruence, so the
arc stays in the domain. For \(|z|=1\) the left and right
characters cancel and
\[
 F(L(z))=F(L).
\]
The maximum modulus
principle applied to \(F(L(z))-F(L)\) gives the equality at
zero. The value there vanishes by the prescribed weight and
\(F(x_0)=0\). The dense cell and continuity give \(F=0\).

For a fixed \(x_0\) this uniqueness expresses \(K(y,x_0)\)
as \(K(x_0,x_0)\) times the product of leading minors
\(\Delta_i\) of the coordinate matrices of \(y\) with
powers \(\lambda_i-\lambda_{i+1}\), with the weight
extended by zero at \(r+1\). This product has precisely
the required two characters. Possible negative powers
occur only at \(p,r\), where the minors do not vanish.
At a place with \(p,q>0\) the weights decrease within
each block and
\[
 \lambda_{p+1}-\lambda_p\ge n.
\]

On an orthonormal coordinate matrix \(U\), all the
leading minors are bounded by \(|\Delta_p(U)|\ge1\)
and \(|\det U|=1\). Write \(P=\mathrm{diag}(I_p,0)\) and \(Q=I-P\).
For \(U^*JU=J\) set \(N=PU+Q\) and \(N'=P+QU\).
Then
\[
 N^*N=U^*PU+Q=U^*(P-Q)U+U^*QU+Q=P+U^*QU=N'^*N'.
\]
The first
\(p\) columns of \(U\), written \(\binom{A}{C}\), satisfy
\(A^*A=I_p+C^*C\), so \(|\det N|=|\det A|\ge1\).
Writing \(W=N'N^{-1}\), we obtain
\[
 W^*W=(N^{-1})^*N'^*N'N^{-1}=(N^{-1})^*N^*NN^{-1}=I .
\] For \(N_i\) retaining the first
\(i\) rows of \(U\) and taking the other rows from the identity,
each row of \(N_iN^{-1}\) is a row of \(I\) or \(N'N^{-1}\).
All have norm one, so Hadamard gives
\[
      |\Delta_i(U)|=|\det N_i|\le|\det N|=|\Delta_p(U)|.
\]

The sum of the minor exponents up to \(r-1\) telescopes
to \(\lambda_1-\lambda_r\le-n\).
It follows that \(|K(y,x_0)|\) is bounded by a
constant times the product of \(|\Delta_p(y)|^{-n}\)
on the slab. The triangular factors there change
this bound only by a bounded factor. At a definite
place the slab is bounded and separated from singular
frames. A left unitary change and right triangular
covariance give the same bounds with any fixed first
frame \(x\).

\proofstep{Integrability in the noncompact directions.}
For integrability at an indefinite place write the
first \(p\) columns of \(y\) as \(\binom{A}{C}\).
The matrix
\[
 G=A^*A-C^*C
\]
stays in a bounded
neighborhood of \(I_p\) in the positive cone.
At fixed first columns use a left unitary isometry
to move them to \(\binom{G^{1/2}}{0}\). On the last
\(q\) columns this change has Jacobian one. The
cross Gram block bounds their top block and the
last Gram block then bounds their lower block.
Their total volume is bounded independently of
the first columns. For a square complex matrix \(A\),
polar coordinates have measure proportional to
\(d(A^*A)\) times compact unitary measure. This follows
from the triangular decomposition with positive
diagonal entries \(r_{ii}\), where both the matrix
and Gram Jacobians have the factor
\(\prod_i r_{ii}^{2(p-i)+1}\).

Since
\[
 A^*A=C^*C+G,
\]
integration in \(A\) amounts
to integration in a bounded set of \(G\) and is bounded
by a constant times \(\det(I_p+C^*C)^{-n/2}\).
To integrate \(C\), remove one row \(v\) and put
\(H=I_p+C_{\rm rem}^*C_{\rm rem}\). The determinant lemma
and the substitution \(v=wH^{1/2}\) give
\[
 \int_{\mathbb C^p}\det(H+v^*v)^{-a}\,dv
    =\det(H)^{1-a}
          \int_{\mathbb C^p}(1+\|w\|^2)^{-a}\,dw .
\]
The last integral is finite for \(a>p\). Starting with
\(a=n/2\), the last of the \(q\) removals has exponent
\(n/2-q+1>p\), since \(n/2\ge p+q\).
This proves absolute integrability of \(K(x,\cdot)\) on the
product slab. The pointwise bounds make the kernel bounded,
and the weighted frame-function norm bounds the ordinary norm
on the slab.


\proofstep{Evaluation by the torus average.}
Average a holomorphic frame function \(F\) against
\(\chi_\lambda^{-1}\) on the left diagonal unitary torus
at \(x_0\). Compactness gives a holomorphic average.
The average retains the right weight and has value
\(F(x_0)\) at \(x_0\). It also leaves the integral against
\(K(x_0,y)\) unchanged, by simultaneous invariance
and covariance in the first frame. Both this average
and \(K(y,x_0)\) have the two characters from the
uniqueness argument. Therefore
\[
 \int_{\rm slab}K(x_0,y)F(y)\,dy
 =\frac{\int_{\rm slab}|K(x_0,y)|^2\,dy}{K(x_0,x_0)}
   F(x_0).
\]
The numerator is finite and positive by boundedness,
integrability and continuity at \(x_0\). Move \(x_0\)
to any orthonormal frame by a left unitary change of
variables, and then to any frame by right weight
covariance. This proves reproduction with the same
positive constant \(\rho\).
\end{proof}

\paragraph{Unfolding the reproducing integral onto the original rows.}

\begin{proposition}[Spanning by the original theta rows]\label{prop:span-original-theta}
For the level and stable range \(n\ge4D_E+4\) of
Lemma~\ref{lem:theta-overlap-invertible}, unfolding
the doubled Weil kernel spans the bounded holomorphic target of the
prescribed weight by the original theta functions. At the base tuple write \(V_{\rm good}\) for the rational good row and \(\Theta_{\rm good}\) for the projected definite-source lift. After the prescribed scalar identifications,
\[
 \operatorname{im}\Theta_{\rm good}=V_{\rm good}.
\]
\end{proposition}

\begin{proof}
In the doubled Weil calculation the sum of the product
kernel over rational changes of target frame is
the finite sum of theta products, with the first
factor read at \((L,x)\) and the second at the other
target lattice and frame. On the unfolded side
the term for \(L'\) is
\[
    C(L,L')\int_{\rm slab} K(x,y)F'(L',y)\,dy .
\]
To obtain the doubled identity, apply the definite
Hermitian Siegel--Weil identity to the admissible
doubled test function at zero in the auxiliary
variables. The difference from the Eisenstein sum
solves the doubled heat equation. The invariant
sum of highest-polynomial differential operators
therefore also kills that difference.

In the
Eisenstein sum a maximal isotropic subspace of
the doubled space is either a unitary graph or
contains a nonzero isotropic vector in one summand.
Writing \(I(U)\) for the term attached to a maximal isotropic space \(U\), the no-coinvariants condition and the support of the finite Weil pairing give, respectively,
\[
 I(U)=0\quad\text{if }U\text{ is not a graph},
 \qquad I(\operatorname{graph}g)=0\quad\text{if }g\text{ is outside the level group}.
\]
On a graph the finite pairing is the overlap of the two test functions. For the
remaining graphs the Weil action on the translated
frame and the conjugate multiplier give the
integrand above. The double sum and
integration may be interchanged.

For fixed \(x\),
the Gaussian at \(x\) is integrable in \(y\) on a
sufficiently thin slab, the functions \(F'\) are
uniformly bounded on that slab, and
\[
 \sum_{L'}|C(L,L')|<\infty
\]
by the overlap
estimate. The group changes preserve the measure
and the slab, and the unit-modulus multiplier
cancels against the conjugate in the theta product.
Integrate over the fundamental domain for these
group changes. The finite sum of theta products
then yields a finite sum of theta functions in
\((L,x)\), whose scalar coefficients are the integrals
of the second factors against \(F'\).
The unfolded expression is
\[
    \rho\sum_{L'} C(L,L')F'(L',x)=\rho F(L,x).
\]
Dividing the coefficients by \(\rho\) gives the
theta spanning assertion. Only the strict bound
on \(T\) is used for an estimate on the rank.
The overlapping lattice system and the doubled
kernel both use the prescribed finite level in this
application. \label{cor:theta-row-rank-1}

\proofstep{Identification of the lowest-vector realization.}
In the adelic realization, evaluate the target
discrete-series map on its nonzero lowest weight
vector. If \(v_0\) is that lowest vector, cyclicity gives the test for an intertwiner \(A\):
\[
 Av_0=0\quad\Longrightarrow\quad A(U(\mathfrak g)v_0)=0\quad\Longrightarrow\quad A=0.
\]
Here \(\mathfrak g\) denotes the complexified Lie algebra on the Harish--Chandra module.
The bottom evaluation has the given compact
torus weight and is killed by the lowering
operators. After extracting the common
determinant character it gives the holomorphic
frame function of the preceding argument. Cusp
rapid decay gives its weighted boundedness. The
nonzero finite Weil vector of the base pair is
a finite translate of the level test vector for
that argument. Conjugating the frames and
undoing this translation preserve the allowed
harmonic source span. The lowest-vector identity spans the unprojected range of \(\Theta\): unfold each term as a definite-source lift summed over source classes. For the good-type projector \(e\),
\[
 \operatorname{im}(e\Theta)=e(\operatorname{im}\Theta).
\]

\proofstep{Projection onto the good row.}
Spanning before taking finite types gives access to all row coordinates. The determinant at the selected place uses the good isotypic component, so the projection has to be applied to the same spanning family.

Now take the base tuple on the actual row. The
trivial auxiliary character and empty auxiliary
level identify its projected lift with the class
sum in the row construction. For its finite unipotent group \(U\) and character \(\chi\),
the projection is
\[
 e_\chi=|U|^{-1}\sum_{u\in U}\chi(u)^{-1}u.
\]

Changing variables \(v=u_0u\) in the sum yields
\[
 u_0e_\chi=|U|^{-1}\sum_{v\in U}\chi(u_0^{-1}v)^{-1}v
            =\chi(u_0)e_\chi .
\]
Conversely \(uf=\chi(u)f\) implies \(e_\chi f=f\). For the equivariant theta map \(T\),
\[
                  T(e_{\chi,\mathrm{src}}v)
                         = e_\chi T(v),
\]
where the source operator uses the dual transformation
rule of the row reading. This imposes that rule on each
preimage of a good row vector. The equivariance of the row reading
makes the average fix the given target vector.
On the source the average imposes the dual
unipotent rule. Thus a good row with this type
is the reading of a class sum of functions on
the source with the base finite transformation
law.

On a finite list of source classes, the
harmonic coefficient lattice spans its allowed
value space over \(\mathbb C\). For the class
stabilizer, the finite trace with its level
multipliers preserves that lattice. For a stabilizer \(H\) on the complex value space, its trace and projector satisfy
\[
 \operatorname{Tr}_H=\sum_{h\in H}h,\qquad |H|^{-1}\operatorname{Tr}_H=\operatorname{pr}_{H\text{-invariants}}.
\]
The operators include the specified level multipliers. Using the unnormalized trace on
integral vectors therefore spans all such class
functions. Their class sums span the given good
row.

A summand at one source class may be moved
to a finite split-supported representative
away from the excluded places. Rational
automorphy and the finite transformation
rule transport its polynomial in the harmonic
lattice. The target character factor is
absorbed by that rule. Each resulting term
is the complex reading of an admitted original
theta datum. This uses the split-support and
level conditions in the definition of the
allowed original data, and the base-series
rational comparison on the same row.

Finally, any vector of the rational good row
whose complex reading is in the span of those
data lies in their rational span. Indeed if \(W\) is their rational span, extension of
scalars over a field is exact, so the kernel of
\[
         V_0\longrightarrow (V_0/W)\otimes_K\mathbb C
\]
is \(W\). The comparison identifies the complex reading
with this scalar extension. Applying the type projection, which
fixes the good row, gives the full span of the
projected original theta elements in \(V_0\).
The finite set of source classes and the
finite harmonic lattices give a finitely generated
row module.

Equivalently, in the rational
comparison the original theta elements lie
in the fixed finite level lattice. Over its DVR,
a generating set for a finite module contains
a subset projecting to a basis modulo the
maximal ideal. Since this module is torsion-free,
the chosen generators form a lattice basis.
This proves the original full-rank comparison
on the fixed tied data and in the stated stable
range. Its use for the division rank later
also requires the identification with the
selected point's period and row dimension.
\label{cor:theta-row-rank-2}
\end{proof}

\paragraph{Containment for the twisted rows.}
Fix the tested character at one admitted auxiliary
level, and work over its coefficient DVR \(O_k\).
The finite type multiplier on the source level
group is a product of its base character, the
splitting character at the two type primes, and
the local components of the fixed and varying
twists. The base character agrees with the
given row character there. All its values are
roots of unity in the coefficient field.
At a place over a type prime \(g_i\) the field
assumption includes the roots of order dividing
\(g_i(g_i-1)\), since the host places there have
residue field \(\mathbb F_{g_i}\).
At the distinguished and auxiliary places use
the prescribed character values in \(O_k\).
If \(\mu\) denotes the product character and \(u\) a level
element, each of its factors has finite order in the admitted
coefficient field. Thus the normalization is
\[
       \mu(u)=\mu_{\rm base}(u)\mu_{\rm spl}(u)
                          \mu_{\rm fixed}(u)\mu_{\rm aux}(u)
                       \in O_k^\times .
\]

For a twisted generator the source translate
differs from its original representative by a
change supported at the chosen split places.
Those places avoid the coefficient prime and
the fixed exclusions. In the harmonic coefficient
lattice the two readings agree over \(O_k\) after
this change. At places over the coefficient prime
the self-dual lattices agree, and all rational
denominators of the other changes are units there.
Write \(h\) for the resulting finite source point
and \(P\) for its integral harmonic polynomial.
The finite trace of \(P\) on the double coset
of \(h\) is the sum of its rational source
translates with the finite level multipliers.

At a given source point the summands are indexed
by a coset of a finite stabilizer, so this is
a finite sum in its harmonic coefficient lattice.
The reindexing uses the bijection of finite indexing sets
\[
 H\longrightarrow h_0H,\qquad h\longmapsto h_0h,
\]
where \(H\) is the stabilizer and \(h_0H\) is the coset indexing a translated sum. Including the level multiplier gives the dual transformation rule. It belongs to the integral module
of class functions for this level.

Unfold its class sum using the finite source
representatives and the same stabilizer weights
as in the original theta formula. On the kernel,
rational left translation fixes the value at
zero and the right level action is the inverse
of the multiplier in the trace. The two factors
cancel.

For one contributing class let \(n_h\) be its stabilizer
order and \(\Theta_h(P)\) the value of the corresponding
flat series term. Each of the \(n_h\) translated terms
has that same value after the two finite multipliers cancel.
The resulting term in the class sum is
\[
       n_h^{-1}\sum_{i=1}^{n_h}\Theta_h(P)=\Theta_h(P).
\]
The reciprocal arises only in the characteristic-zero
class-sum evaluation, after forming the integral trace. The result is the desired single flat
twisted series, multiplied by its norm-one
determinant character on the source point.
That scalar is a root of unity in \(O_k\).
Including it in the integral trace therefore gives
the correct twisted generator. The trace takes
no inverse stabilizer order on the integral
module.

By the comparisons on the fixed row, the class
sum is the reading of an element of the finite
level lattice. The twisted generator has the same
reading. Injectivity over the extension field
identifies these vectors, so the generator belongs
to that lattice before applying the good-type
projection. This proves the twisted containment
used in the signed-index comparison. It applies
to both sides and every admitted character of
the specified depth. \label{cor:theta-row-rank-3}

\subsubsection*{The row lattice measured by the twist}
\label{subsec:twist-row-lattices}

The spanning argument is rational. For a bound on division one must also compare the lattice of bounded forms in that span with the lattice actually used to evaluate the twist.

Fix one row and side at the chosen coefficient place. Let
\(F_0\) be the rational subspace on which the prescribed
unipotent subgroup acts by the prescribed character, and
let \(B_{\rm sup}\) be the integral lattice of bounded
forms. Write \(e\) for the rational central type idempotent,
\(g\) for the order of the finite row group, and put
\[
 M_{\rm all}=B_{\rm sup}\cap F_0,\qquad V_0=eF_0,\qquad
 B=M_{\rm all}\cap V_0,\qquad A=ge .
\]
\begin{lemma}[The original integral type lattice]
\label{lem:projected-row-type}
Over the fraction field \(K\) of the coefficient DVR the intersections satisfy
\[
 M_{\rm all}\otimes K=F_0,\qquad B\otimes K=V_0,\qquad A(M_{\rm all})\subseteq B\quad(A=ge).
\] Any residual
simultaneous Hecke eigencharacter of \(B\) is the reduction of
one of the rational signature characters on the same row
and side.
\end{lemma}

\begin{proof}
\proofstep{Integral intersection and type projection.}
These intersections use the given embeddings in the rational
space. In the row model \(B_{\rm sup}\) is a full finitely
generated lattice stable under the finite group and the
admitted Hecke operators, which commute with the group
action. Clearing denominators in this lattice shows that
\(M_{\rm all}\) spans \(F_0\) and \(B\) spans \(V_0\);
both intersections are finitely generated over the DVR.
The Hecke operators preserve each intersection by commutation.
The integral central representative of \(ge\) preserves \(B_{\rm sup}\) and sends \(F_0\) onto \(V_0\) since \(g\ne0\). Writing \(K\) for the fraction field of the coefficient DVR, we obtain
\[
 ge(M_{\rm all})\subseteq B_{\rm sup}\cap V_0=B,\qquad
 \operatorname{span}_{K}ge(M_{\rm all})=V_0 .
\]
These give the lattice and the operator used in the row
determinant. Dividing by \(g\) on the integral lattice at
this stage could alter a length.

\proofstep{One actual residual signature.}
The residual character test used after pair removal follows
from the rational signature decomposition on this same
\(V_0\). Let \(x\in B\) lift a nonzero residual simultaneous
eigenvector of the Hecke operators. Decompose in the rational simultaneous eigenspaces as
\[
 x=\sum_a y_a,
\]
where the signature \(a\) on this row and side has eigenvalues \(e_a(T)\), integral at the coefficient place. Suppose each one differs from the
residual character on some admitted operator \(T_a\). Then
\[
 \prod_a(T_a-e_a(T_a))
\]
kills \(x\), by commutativity
on the sum of these eigenspaces. Each factor acts on
\(\bar x\) by a nonzero scalar, so its product has nonzero
value on \(\bar x\). This is a contradiction since each
factor preserves \(B\). Thus some signature has every
required residual eigenvalue. A decomposition of the
integral lattice into its characteristic-zero eigenspaces
is unnecessary for this test.

\end{proof}

\begin{lemma}[Clearing the universal row coefficients]
\label{lem:universal-row-denominator}
For the finite original basis there exists \(c\ge0\) such that every individual universal coefficient function \(\Xi\) at every allowed level satisfies
\[
 \ell^c\Xi\in B_{\rm sup}.
\]
This uses the boundedness and common-phase comparisons. The same exponent
clears the sum of coefficient functions whose class labels identify
in the row.
\end{lemma}

\begin{proof}
\proofstep{A denominator for all individual class coefficients.}
For the common denominator, write the universal lift of a
chosen datum at \(Q,M\) as a finite group-ring sum
\(\sum_\delta \Xi_\delta[\delta]\). In the row realization,
the coefficient functions for the individual class labels
lie in the rational row and have the prescribed unipotent
type. If the bounded-coefficient estimate gives
\[
 \ell^{c(a)}\Xi\in B_{\rm sup}
\]
for each such coefficient
of datum \(a\), with \(c(a)\) independent of \(Q,M\),
then the same exponent clears all the \(\Xi_\delta\)
in \(M_{\rm all}\). Indeed the map from class labels to
coordinates can identify labels, and its coefficient at
\(\delta\) is the sum of their coefficients. Thus the exponent used in \(C_Q\) is uniform in the level and can be chosen by
\[
 c=\max_{a\text{ in the original basis}}c(a).
\]
The bound needed here is on the individual
universal coefficients as the modulus varies.

\proofstep{The common phase and its inverse in the coefficient field.}
Propositions~\ref{prop:theta-uniform-cover-clearing}
and~\ref{prop:ray-cover-descent} give one phase and one
clearing power on each generator of the harmonic coefficient
lattice. A finite linear combination of these generators
admits the maximum of their clearing powers, since the phase
is common and the coefficient scalars are integral.
The harmonic polynomial realization is injective in the
range where the sum of the two signature dimensions is at
most the target rank. Consequently the expression for the
coefficient function applies to its given polynomial.

Let
\(\mathrm{rd}\) be the complex realization of the rational
row. With its fixed comparison scalar incorporated, the
integral combination gives a single \(\lambda\in\mathbb C^\times\)
such that, for every coefficient \(x\) under consideration,
there is \(x'\in B_{\rm sup}\) with
\[
 \mathrm{rd}(x')=\lambda\mathrm{rd}(x).
\]
Faithfulness
here includes linear independence over \(\mathbb C\) of
the images of a rational basis. If some \(x\ne0\), comparing
a nonzero coordinate gives \(\lambda\) in the embedded
fraction field. Clear its inverse once by a power of
\(\ell\); injectivity gives that power times every \(x\)
in \(B_{\rm sup}\). If every \(x\) vanishes, exponent zero
suffices. This gives the coefficient estimate in
\(B_{\rm sup}\) used above.

\end{proof}

Let \(L_0=e\Theta_0\) be the projected original theta module
of the specified row indices, and \(L_\psi=e\Theta_\psi\)
the projected integral module generated by their allowed
twisted series on the chosen side, over the DVR \(O_k\)
of the character. Write \(B_k\) for the scalar extension
of \(B\). For full lattices use the signed index
\([U:V]_\ell\), the valuation of the determinant of a basis
of \(V\) in a basis of \(U\), normalized by \(v(\ell)=1\).
Negative values are allowed.

\begin{proposition}[Signed index of the actual twisted row]
\label{prop:twisted-row-index}
With the faithful coefficient comparisons and the selected
character and row defined above, \(L_\psi\) is full in \(V_{0,k}\)
and
\[
          [(L_0)_k:L_\psi]_\ell\le\epsilon_{j,b}.
\]
Its inclusion in the projected level lattice persists under
finite coefficient extension, and the index bound persists after
extension in the normalization \(v(\ell)=1\).
\end{proposition}

\begin{proof}
\proofstep{Specialization of the universal series.}
The chosen basis and full-rank
condition give \(L_0\) full in \(V_0\). The twisted type law puts \(L_\psi\subset V_{0,k}\). It is finitely generated, being a submodule of the projected finitely generated level lattice over a DVR. At a character of order dividing \(\ell^k\)
factor evaluation through its chosen ray coordinates. Avoidance makes every basis datum in \(C_Q\) prime to \(Q\) in level. The specialized row series belongs to the integral span of that datum's twisted series by the following calculation. At a chosen place of residue
field \(k_P\), size \(q\), put \(r_0\le n_0\) for the two
matrix dimensions and let \(y\) have \(r_0\) rows. Sum over
codimension-\(r_0\) planes \(U\subset k_P^{n_0}\).

For each plane the normalized local contribution is
\[
 q^{-r_0(n_0-r_0)}\begin{cases}[d_U(y)],&\operatorname{row}(y)\oplus U=k_P^{n_0},\\ 0,&\text{otherwise}.\end{cases}
\]
Here \([d_U(y)]\) is the group-ring element attached to the quotient determinant. Choose a basis adapted
to each plane. Then
\[
 d_U(y)=c_U\det((yg_U)_{[r_0]})
\]
for a fixed residue unit \(c_U\), an invertible change of
basis \(g_U\), and the first square block of \(yg_U\).
The flat local function at the coordinate character \(\chi\) is
\[
 B\longmapsto\begin{cases}\chi(\det B),& B\text{ invertible},\\
 0,& B\text{ singular}.
 \end{cases}
\]
The preceding formula expresses the evaluation as a sum of its translates. The normalization times the residue character of \(c_U\) is integral since \(\ell\nmid q\). Lift the changes of basis integrally at the chosen places. Multiplication and linearity of theta summation now give an integral combination of the allowed translates. The character of the determinant
agrees with the specialization by Artin reciprocity and the
coordinate restriction law.

In the rational row comparison
both forms have these same complex evaluations; injectivity
after scalar extension gives membership on the prescribed
row. After applying \(A\),
its value belongs to \(g L_\psi\). Hence every full lattice
\(T\) containing \(g L_\psi\) satisfies the same estimate
from the determinant,
\[
           [B_k:T]_\ell\le [B:gL_0]_\ell+\epsilon_{j,b}.
\]
These comparisons concern the whole specialized theta
module; equality between a single universal column and a
chosen basis vector of that module is not needed.

\proofstep{Full rank and scalar cancellation in the index.}
The lattice index only measures a determinant if the specialized rows have full rank. It remains to verify this rank directly before cancelling the scalar common to the two covolumes.

In particular \(L_\psi\) spans \(V_{0,k}\). To see this
without assuming its rank in the index inequality, take
\[
 T_n=L_\psi+\ell^n B_k.
\]
If two consecutive terms of this
descending chain coincide, Nakayama on \(T_n/L_\psi\)
gives \(T_n=L_\psi\), which has full rank. Otherwise every
strict inclusion increases \([B_k:T_n]_\ell\) by at least
\(1/e_{O_k}\); every \(T_n\) contains \(g L_\psi\), so
this contradicts the common bound. We may now take
\(T=gL_\psi\). The valuation of a determinant is unchanged
under normalized finite DVR extension, and indexing two
full lattices after multiplying both by the same nonzero
scalar leaves their relative index unchanged. Additivity
therefore gives
\[
           [(L_0)_k:L_\psi]_\ell\le\epsilon_{j,b}.
\]

For any further finite extension \(O'/O_k\) and full
reference lattice \(M'\) in that space the bound is
\(
 [M':(L_\psi)_{O'}]_\ell
 \le [M':(L_0)_{O'}]_\ell+\epsilon_{j,b}.
\)
The inclusion into the projected level lattice persists under
this flat extension. These are the lattice, inclusion and
index assertions used at every permitted power of the
single selected character. Their geometric inputs are the
bounded-coefficient estimate, type preservation, the faithful row
comparison and integral containment, and the original
full-rank basis on these very levels. \label{cor:twist-row-lattices-1}
\end{proof}


\subsubsection*{Conversion of the split row bound}

The auxiliary character and both row lattices are now fixed. We return to the small split primes and insert their local exponents into the normalized congruence length.

\begin{proposition}[Weighted split loss at the small primes]\label{prop:small-split-loss}
On these split row and column declarations, let
\(H=b+S_c+\eta_c+t+2\), use the full cyclotomic depth \(s\) in Lemma~\ref{lem:small-split-finite-colength}, and
put \(\Lambda=s+\log_\ell(2H)+1\).
Let \(c_{\rm loc}\) denote the normalized local cost after subtracting the scalar and finite Selmer baseline. Then
\[
 c_{\rm loc}\le 6000 f m\,H s+30000 f^2t^2 R\Lambda.
\]
For the ranks and weights on the selected family this contributes
\(O(fmbz^2)\) after summation over \(\mathcal S_{\rm sp}\)
with weight \(\log\ell\). Constants and thresholds are uniform in
the selected point and good eigencharacter.
\end{proposition}

\begin{proof}
\proofstep{Put the lifting exponents in the weight and depth ranges of the selected point.}
Put \(H=b+S_c+\eta_c+t+2\), where \(S_c,\eta_c\) are the
two shift and weight sizes in the selected parameter choice.
Here \(b\le H\le6b\). Write
\[
                    \Lambda=s+\log_\ell(2H)+1 .
\]
The local lifting and invariant estimates before numerical
conversion use
\[
 \lambda_N(h)=s+\lfloor\log_\ell(2h)\rfloor+ \lfloor\log_\ell N\rfloor+1,
\]
with \(h\) the absolute
jump cutoff in each estimate. The selected level satisfies
\(N\le8H\), and both jump cutoffs used here satisfy \(h\le4H\).
Therefore \(\lambda_N(h)\le3\Lambda\), since
\[
 \log_\ell(8H)\le\log_\ell(2H)+2\le\Lambda.
\]

The local lift exponent is \(32(2hs+1)+\lambda_N(h)\).
For the invariant and block estimates, the bounds are
\(20(4nR+1)\lambda_N(h)\) and
\(40(t+1)(4nR+1)\lambda_N(h)\), respectively.
Using \(f=2n\), \(t\ge1\), and
\[
 \Lambda\le4Hs
\]
gives the coefficient \(1000\) below.

\proofstep{Subtract the scalar and rank-complement terms on the same cut.}
The split row inequality bounds the normalized localized
congruence term for a good Hecke character by a sum of a
baseline and the local row and column costs described below.
In the bound used for the numerator, the baseline is
\(s_\pi+V_\ell+v_\ell(h_E)+\ell^{-2}\).
Here \(s_\pi\) is the relaxed global Selmer length,
\(V_\ell\) is the sum of strict localization lengths through
twist \(t-1\), and \(h_E\) is the class number of \(E\).
The scalar and finite Selmer terms have been separated at this point. The remaining tame torsion term is controlled by the local killing condition, which bounds
\[
 t\sum_v v_\ell(\zeta_v-1)
\]
in this formula by \(\ell^{-2}\),
where the \(\zeta_v\ne1\) are the indicated tame test values.

Let \(D_c,D_r\) denote the sums of the local degrees of the
column and row places, and let \(k_c,k_r\) be their numbers.
All these places are over \(\ell\). The degree formula
gives \(k_c\le D_c\le2f\) and \(k_r\le D_r\le2f\).

Write \(e'\) for the lift exponent, \(c_B\) for the block
bound, \(v_\Delta\) for the factorial denominator length
and \(h_0\) for the scalar invariant bound above. The two
column terms involving local weights have bounds
\(r c_0(u)\) and \(h_W(u)\) at a column place \(u\).
The cyclotomic exponent and local invariant estimates used
in the split bound have the following size requirements
with \(c=1000\):
\[
 \begin{gathered}
  e'\le cHs,\qquad c_B\le ctfR\Lambda,\qquad
  c_0(u)\le cfR\Lambda,\qquad h_W(u)\le crfR\Lambda,\qquad
  h_0\le c\Lambda,\\
  v_\Delta\le ts+v_\ell(t!)\le t(s+1)\le t\Lambda .
 \end{gathered}
\]
The inequality for \(h_0\) follows from the scalar
calculation above. The second inequality for \(v_\Delta\)
uses \(v_\ell(t!)=\sum_{a\ge1}\lfloor t/\ell^a\rfloor\le t\).

\proofstep{Sum the column and row terms with their respective local degrees.}
The split row formula then has total cost above its
baseline at most
\begin{align*}
 D_c m(2e'+v_\Delta)
 &+k_c(e'+2v_\Delta+2crfR\Lambda)\
 &+t(D_r+2k_r)c_B+4R(k_c+1)v_\Delta+tk_ch_0 .
\end{align*}
The terms containing \(e'\) contribute at most \(6 c f m\,H s\).
In order, the six remaining terms contribute at most
\(2,4,4,6,12,2\) times \(cf^2t^2R\Lambda\).
For instance \(4R(k_c+1)v_\Delta\le12fRt\Lambda\),
and \(f,t,c\ge1\). Hence the excess over the baseline is
bounded by
\begin{equation}\label{eq:split-row-cost}
             6000 f m\,H s+30000 f^2t^2R\Lambda .
\end{equation}
This calculation uses the cohomological split row
inequality on the same row-and-column declaration.

\proofstep{Take the logarithmically weighted sum only after the local conversion.}
For a uniform size comparison put \(z=x^\epsilon\), as in
Section~\ref{sec:comparison}, and take a fixed \(C_0\ge1\)
with \(f,t\le x^{C_0}\), \(R\le2x^{C_0}\),
and \(Q_S\le n^2\) with \(n\le x^{C_0}\).
For \(\ell\ne q,d\), \(v_\ell(N)=v_\ell(Q_S)\), so
\(s\le3C_0\log x\) for \(\log x\ge1\).
Also \(\log_\ell(2H)\le2z+6\), using \(H\le6b\)
and \(\log\ell>1/2\).
For \(z\ge1\), this gives \(\Lambda\le(3C_0+9)xz\).

For any fixed nonnegative coefficients \(A,B\) we obtain
\[
 A f m\,H s+B f^2t^2R\Lambda
 \le18AC_0 fmb\log x+
             2B(3C_0+9)x^{5C_0+1}z .
\]
Thus \eqref{eq:split-row-cost} is bounded by
\(C fmb\log x+Cx^Cz^2\) with a single \(C\) depending on
\(C_0\). The second term is at most \(C fmb\) past one
threshold since \(b=\exp(x^\epsilon)\). These constants
and the threshold are independent of the point, the
coefficient field and the particular good eigencharacter.
The weights of all the primes just treated satisfy
\[
    \sum_{\ell\in\mathcal S_{\rm sp}}\log\ell
    =\log\!\prod_{\ell\mid 3Q_S}\ell\le\log(3Q_S).
\]
The resulting weighted cost is \(O(fmbz^2)\) beyond the
threshold, using \(\log x\le z\).

 The normalized Galois algebra, tame tests and
global restriction argument apply at the retained points.
The character and finite-class estimates use the same local
lattices as the invariant estimate.
The simultaneous twist selection is given in
Section~\ref{subsec:retained-twist}.
The remaining prime ranges enter the same division numerator.
\label{cor:small-prime-lengths-4}
\end{proof}

